\documentclass[11pt,a4paper]{amsart}

\usepackage[margin=1in]{geometry}
\usepackage[T1]{fontenc}
\usepackage[utf8]{inputenc}
\usepackage{lmodern}
\usepackage{microtype}
\emergencystretch=3em
\usepackage{amsmath,amssymb,amsthm,mathtools}
\usepackage{enumitem}
\usepackage{float}
\usepackage{aliascnt}
\usepackage[colorlinks=true,linkcolor=blue,citecolor=blue,urlcolor=blue]{hyperref}
\usepackage[capitalize,nameinlink]{cleveref}

\numberwithin{equation}{section}

\newtheorem{theorem}{Theorem}[section]
\newaliascnt{proposition}{theorem}
\newtheorem{proposition}[proposition]{Proposition}
\aliascntresetthe{proposition}
\newaliascnt{lemma}{theorem}
\newtheorem{lemma}[lemma]{Lemma}
\aliascntresetthe{lemma}
\newaliascnt{corollary}{theorem}
\newtheorem{corollary}[corollary]{Corollary}
\aliascntresetthe{corollary}
\newaliascnt{definition}{theorem}
\newtheorem{definition}[definition]{Definition}
\aliascntresetthe{definition}
\theoremstyle{remark}
\newaliascnt{remark}{theorem}
\newtheorem{remark}[remark]{Remark}
\aliascntresetthe{remark}

\crefname{theorem}{Theorem}{Theorems}
\Crefname{theorem}{Theorem}{Theorems}
\crefname{proposition}{Proposition}{Propositions}
\Crefname{proposition}{Proposition}{Propositions}
\crefname{lemma}{Lemma}{Lemmas}
\Crefname{lemma}{Lemma}{Lemmas}
\crefname{corollary}{Corollary}{Corollaries}
\Crefname{corollary}{Corollary}{Corollaries}
\crefname{definition}{Definition}{Definitions}
\Crefname{definition}{Definition}{Definitions}
\crefname{remark}{Remark}{Remarks}
\Crefname{remark}{Remark}{Remarks}

\DeclareMathOperator{\RS}{RS}
\DeclareMathOperator{\MCA}{MCA}
\DeclareMathOperator{\CA}{CA}
\DeclareMathOperator{\rank}{rank}
\DeclareMathOperator{\Span}{Span}
\DeclareMathOperator{\Spec}{Spec}
\DeclareMathOperator{\Hom}{Hom}
\DeclareMathOperator{\im}{im}
\DeclareMathOperator{\codim}{codim}
\DeclareMathOperator{\supp}{supp}
\DeclareMathOperator{\Fib}{Fib}
\DeclareMathOperator{\Sym}{Sym}
\DeclareMathOperator{\Gr}{Gr}
\DeclareMathOperator{\Ent}{H}
\DeclareMathOperator{\ord}{ord}
\DeclareMathOperator{\Tr}{Tr}
\DeclareMathOperator{\Gal}{Gal}
\DeclareMathOperator{\Res}{Res}
\DeclareMathOperator{\Aut}{Aut}
\DeclareMathOperator{\diag}{diag}
\DeclareMathOperator{\wt}{wt}

\newcommand{\F}{\mathbb F}
\newcommand{\B}{\mathbb B}
\newcommand{\K}{\mathbb K}
\newcommand{\Z}{\mathbb Z}
\newcommand{\A}{\mathbb A}
\newcommand{\Pj}{\mathbb P}
\newcommand{\C}{\mathbb C}
\newcommand{\Prob}{\mathbb P}
\newcommand{\E}{\mathbb E}
\newcommand{\eps}{\varepsilon}
\newcommand{\Omc}{\Omega^{\circ}}
\newcommand{\Ccal}{\mathcal C}
\newcommand{\Dcal}{\mathcal D}
\newcommand{\Ecal}{\mathcal E}
\newcommand{\Fcal}{\mathcal F}
\newcommand{\Gcal}{\mathcal G}
\newcommand{\Hcal}{\mathcal H}
\newcommand{\Ical}{\mathcal I}
\newcommand{\Jcal}{\mathcal J}
\newcommand{\Lcal}{\mathcal L}
\newcommand{\Mcal}{\mathcal M}
\newcommand{\Ncal}{\mathcal N}
\newcommand{\Pcal}{\mathcal P}
\newcommand{\Qcal}{\mathcal Q}
\newcommand{\Rcal}{\mathcal R}
\newcommand{\Scal}{\mathcal S}
\newcommand{\Tcal}{\mathcal T}
\newcommand{\Ucal}{\mathcal U}
\newcommand{\Vcal}{\mathcal V}
\newcommand{\Wcal}{\mathcal W}
\newcommand{\Xcal}{\mathcal X}
\newcommand{\Ycal}{\mathcal Y}
\newcommand{\Zcal}{\mathcal Z}
\newcommand{\Gord}{\Gamma^{\rm ord}}
\newcommand{\Gsid}{\Gamma^{\rm sid}}
\newcommand{\Gprim}{\Gamma^{\rm prim}}
\newcommand{\barN}{\overline N}
\newcommand{\Renyi}{R\'enyi}
\newcommand{\Hb}{H_2}
\newcommand{\UU}{\mathbb U}
\newcommand{\abs}[1]{\lvert #1\rvert}
\newcommand{\norm}[1]{\lVert #1\rVert}
\newcommand{\floor}[1]{\left\lfloor #1\right\rfloor}
\newcommand{\ceil}[1]{\left\lceil #1\right\rceil}
\newcommand{\defeq}{\coloneqq}
\newcommand{\one}{\mathbf 1}
\newcommand{\Fq}{\mathbb F_q}


\title{Reed--Solomon MCA Thresholds}
\author{Przemek Chojecki}
\hypersetup{pdftitle={Reed--Solomon MCA Thresholds},pdfauthor={Przemek Chojecki}}
\subjclass[2020]{94B35, 11T71, 05B40}
\keywords{Reed--Solomon codes, mutual correlated agreement, Proximity Prize,
syndrome geometry, MDS codes, circle codes, prime density}
\date{July 2026}

\begin{document}

\begin{abstract}
We determine exact finite thresholds in broad regimes and rigorous asymptotic threshold compilers for affine support-wise mutual correlated agreement of Reed--Solomon codes.  If \(C\) has parameters \([n,k]\), \(\Gamma\) is a nonempty challenge set, and the integer \(0\le r\le n-k-1\) satisfies \((n-r)^2\ge n(k+r)\), then \(B_{C,\Gamma}^{\rm MCA}(n-r)=\min\{|\Gamma|,r+1\}\); the same upper bound holds for every linear MDS code.  An exact CA--sparse decomposition, together with the cited lossless unique-decoding theorem, extends the staircase asymptotically through every fixed relative radius below half the minimum distance, while syndrome--secant geometry and locator-prefix pole lines give independent finite upper and lower mechanisms.  We obtain exact target compilers, four certified multiplicative rows at target \(2^{-128}\), a shallow-prefix asymptotic exponent, and the formula \(\delta^*=1-\rho-g+o(1)\) whenever matching certificates meet at agreement \(k+1+gn+o(n)\).  Generalized Reed--Solomon and explicitly diagonalized circle codes inherit the finite conclusions, but no unrestricted near-capacity smooth/circle claim is made.
\end{abstract}

\maketitle

\section{Introduction and main results}\label{sec:introduction}

By a finite Reed--Solomon \emph{row} we mean data
\((\F,D,k)\), where \(\F\) is a finite field,
\(D\subseteq\F\) is an \emph{evaluation domain} of size \(n\), and
\(1\le k<n\).  The associated code
\(C=\RS_\F(D,k)\subseteq\F^D\) consists of the evaluations on \(D\) of
polynomials of degree less than \(k\).  Thus a support is simply a subset
of the coordinate set \(D\).  A word \(u\in\F^D\) is
\emph{explained on a support} \(S\subseteq D\) if its restriction to
\(S\) agrees with one such degree-less-than-\(k\) polynomial.
When the evaluation points lie in a designated coefficient subfield, we
write \(D\subseteq\B\subseteq\F\); \(\B\) controls locator coefficients
and domain structure, while \(\F\) controls codeword coefficients and
slopes.
Throughout,
\(h(x)=-x\log x-(1-x)\log(1-x)\) denotes binary entropy in natural
units, and \(H_2(x)=h(x)/\log 2\) denotes binary entropy in bits, with the
usual convention \(0\log0=0\).

A received pair \(r=(r_0,r_1)\in(\F^D)^2\) determines the
\emph{received line} \(\{r_0+\gamma r_1:\gamma\in\F\}\).  The pair is
simultaneously explained on \(S\) if each of the two received words has an explanation
there, possibly by different polynomials.  A finite slope \(\gamma\) is
\emph{support-wise MCA-bad at agreement \(a\)} if some support
\(S\subseteq D\), \(\abs S\ge a\), explains
\(r_0+\gamma r_1\) but does not simultaneously explain the pair.  The
same support is used in both tests; this is the support-wise convention.
A tuple \((\gamma,S,h)\), where \(\deg h<k\) and
\(h|_S=(r_0+\gamma r_1)|_S\), is a \emph{witness} to that bad slope when
the pair is not simultaneously explained on \(S\), or equivalently is
\emph{support-wise nontrivial} there.  Finally,
\(B_C^{\rm MCA}(a)\) is the maximum, over all received pairs, of the
number of distinct bad slopes.  It counts slopes, not witnesses or
supports.

This is exactly the affine-line definition used by the Proximity Prize
\cite{ProximityPrize} and by Arnon--Boneh--Fenzi
\cite[Definition~4.3]{ABF26}: for a real radius \(\delta\), a witnessing
support has size at least \((1-\delta)n\).  Since support size is integral,
put
\[
 r(\delta)=\floor{\delta n},\qquad a(\delta)=n-r(\delta).
\]
Thus the error is a right-continuous integer staircase.  In particular, a
statement about a ``largest'' real radius needs an endpoint convention;
in the exact rows below the safe set is half open and has a supremum but no
largest element.

Here is the asymptotic formulation used throughout.  If
\(\varnothing\ne\Gamma\subseteq\F\) is the \emph{challenge set}, meaning
the set of slopes allowed in the test, let
\(B_{C,\Gamma}^{\rm MCA}(a)\) denote the same maximum with
\(\gamma\in\Gamma\), and at the radius grid point \(\delta=1-a/n\) put
\[
 e_{\rm MCA,\Gamma}(C,\delta)
   =\frac{B_{C,\Gamma}^{\rm MCA}(a)}{\abs\Gamma}.
 \tag{1.1}\label{eq:intro-normalized-mca}
\]
For the standard full-field challenge we abbreviate
\[
 e_{\rm MCA}(C,1-a/n)=\frac{B_C^{\rm MCA}(a)}{\abs\F},
 \qquad
 \delta_{C,\mathrm{sup}}^*(\eps)=
 \sup\{0\le\delta<1-k/n:e_{\rm MCA}(C,\delta)\le\eps\},
\]
where the error at an arbitrary real radius is evaluated at
\(a(\delta)=n-\floor{\delta n}\).  We set this supremum to zero when the
safe set is empty.  We separately write
\[
 \delta_{C,\mathrm{grid}}^*(\eps)=
 \max\left\{\frac rn:0\le r\le n-k-1,
       e_{\rm MCA}(C,r/n)\le\eps\right\}
\]
when this set is nonempty.  If the first unsafe integer radius is \(r_0\),
then
\(\delta_{C,\mathrm{grid}}^*=(r_0-1)/n\) while
\(\delta_{C,\mathrm{sup}}^*=r_0/n\), and the latter endpoint is unsafe.
Taking
\(\Gamma=\F\) in \eqref{eq:intro-normalized-mca} gives exactly this
customary formulation.
By a \emph{row sequence} we mean a family
\(C_n=\RS_{\F_n}(D_n,k_n)\) indexed so that \(\abs{D_n}=n\); the fields,
domains, and dimensions may all vary with \(n\).  Every claimed upper
bound is uniform over received pairs, whereas a lower construction need
only exhibit the stated received pair.  For such a sequence,
challenge sets
\(\Gamma_n\), and targets \(0\le\eps_n\le1\), define
\[
 \begin{split}
 B_n^*&=\floor{\eps_n\abs{\Gamma_n}},\\
 a_n^*&=\min\{a\in\{k_n+1,\ldots,n\}:
              B_{C_n,\Gamma_n}^{\rm MCA}(a)\le B_n^*\},
 \qquad
 \delta_n^*=1-\frac{a_n^*}{n},
 \end{split}
 \tag{1.2}\label{eq:intro-asymptotic-threshold}
\]
whenever the set defining \(a_n^*\) is nonempty.  Here \(\delta_n^*\)
denotes the largest safe \emph{grid} radius; the corresponding continuous
supremum is one grid step to the right whenever \(a_n^*-1\) is the first
unsafe agreement.  Thus the asymptotic RS--MCA problem is to determine
\(\delta_n^*\), or equivalently the first
\emph{safe agreement}: agreement \(a\) is safe when its bad-slope
numerator is at most \(B_n^*\), and unsafe otherwise.  Since the numerator
is nonincreasing in \(a\), the resulting safe/unsafe boundary is the
\emph{MCA frontier}.  This formulation retains the target and challenge
size; replacing them by a zero-exponent slogan loses real information when
the target is small.

\paragraph{Position relative to current work.}
A proximity-gap theorem says, roughly, that an affine line which is not
entirely close to the code contains only a controlled number of close
points.  For a rate-\(\rho\) Reed--Solomon code, the unconditional
positive theory for arbitrary prescribed domains reaches the Johnson
radius \(1-\sqrt{\rho}\), with recent refinements and MCA conversions
\cite{BCIKS20,BCHKS25,BordageChiesaGuanManzur25,Habock25}.  Capacity-level positive results concern
folded, multiplicity, subspace-design, or random-evaluation models
\cite{GoyalGuruswami25,GoyalGuruswamiSunWootters26}, while fixed smooth
subgroups have explicit near-capacity obstructions
\cite{CS25,KKH26,Kambire26}.  Two priority distinctions are essential.
First, the two-radius theorem \cite[Thm.~1.3]{BCHKS25}, the below-half
equality of the global CA and MCA errors
\cite[Lem.~4.10]{WHIR} (restated as \cite[Lem.~4.6]{ABF26}), and the
matching perturbation already imply the exact \(r+1\) numerator for
constant-rate RS rows through almost half their minimum distance.  The
exact CA/sparse decomposition proved here is a self-contained refinement
that also retains restricted challenge sets.  Thus neither the
constant-rate consequence of the quadratic theorem nor the four large
Proth-row applications below is claimed as a new consequence.  Second, the
syndrome-space viewpoint also occurs independently in \cite{YuanZhu26}.

We are unaware of a prior finite MDS statement of the following form: the
quadratic overlap argument applies without a large-length hypothesis,
gives an MDS upper bound, and covers parameter pockets not certified by
the specialized BCHKS estimate used below.  We make the literature implication integral,
transport it exactly to standard circle codes admitting the displayed
stereographic or diagonal presentation, and record positive-density
families without promoting any of these statements to a near-capacity
theorem.  The unrestricted smooth/circle frontier beyond half distance
still requires new information about lists or received-line incidence.

\begin{theorem}[Self-contained quadratic target theorem]
\label{thm:intro-unconditional-asymptotic-deep}
Let \(C_n=\RS_{\F_n}(D_n,k_n)\) be any Reed--Solomon row sequence over
arbitrary characteristics, let
\(\varnothing\ne\Gamma_n\subseteq\F_n\), and put
\(B_n^*=\floor{\eps_n\abs{\Gamma_n}}\) as in
\eqref{eq:intro-asymptotic-threshold}.
\begin{enumerate}[label=\textup{(\roman*)},leftmargin=*]
\item If \(B_n^*=0\), no agreement is safe and the set defining
      \(a_n^*\) is empty.
\item If \(B_n^*=\abs{\Gamma_n}\), then \(a_n^*=k_n+1\).
\item Suppose eventually
\[
 1\le B_n^*<\abs{\Gamma_n},\qquad
 B_n^*\le n-k_n-1,\qquad
 (B_n^*-1)^2\ge
 n\bigl(3(B_n^*-1)-(n-k_n)\bigr).                  \tag{1.3}
 \label{eq:intro-unconditional-asymptotic-deep}
\]
Then the threshold exists and is exactly
\[
 a_n^*=n-B_n^*+1,\qquad
 \delta_n^*=\frac{B_n^*-1}{n}.                    \tag{1.4}
 \label{eq:intro-unconditional-asymptotic-deep-answer}
\]
In particular, if \(k_n/n\to\rho\),
\[
 \frac{B_n^*}{n}\longrightarrow b<
 \alpha_\rho\defeq\frac{3-\sqrt{5+4\rho}}2,
\]
and
\(1\le B_n^*<\abs{\Gamma_n}\) eventually, then
\(\delta_n^*\to b\).  This conclusion applies, without a smoothness,
odd-characteristic, Fourier, atlas, or ray-compiler hypothesis, to every
multiplicative smooth RS row and to every circle code that has been proved
diagonally agreement-equivalent to an RS presentation.
\end{enumerate}
\end{theorem}

\begin{proof}
If \(B_n^*=0\), the universal tangent floor at agreement \(n\) is one,
so no grid point is safe.  If \(B_n^*=\abs{\Gamma_n}\), the numerator is
always at most the size of the challenge set, so the first grid point
\(k_n+1\) is safe.  In the remaining case put
\(a_0=n-B_n^*+1\).  Then \(n-a_0=B_n^*-1\), and the last inequality in
\eqref{eq:intro-unconditional-asymptotic-deep} places \(a_0\) in the
mean-overlap range.  The exact numerator of
\cref{cor:mean-overlap-exact} is
\(\min\{\abs{\Gamma_n},B_n^*\}=B_n^*\), so \(a_0\) is safe.  The second
inequality in \eqref{eq:intro-unconditional-asymptotic-deep} puts
\(a_0-1\) on the admissible grid, and
\cref{prop:universal-tangent-floor} gives at least
\(\min\{\abs{\Gamma_n},B_n^*+1\}=B_n^*+1\) bad slopes there.  Hence
\(a_0-1\) is unsafe, and monotonicity proves
\eqref{eq:intro-unconditional-asymptotic-deep-answer}.  The asymptotic
specialization follows because, after division by \(n^2\), its strict
limiting margin is \(b^2-3b+1-\rho>0\).  Moreover
\(\alpha_\rho<1-\rho\) for \(0<\rho<1\), since this is equivalent to
\(\sqrt{5+4\rho}>1+2\rho\).  Thus \(B_n^*\le n-k_n-1\) eventually,
and hence
\eqref{eq:intro-unconditional-asymptotic-deep} holds for all sufficiently
large \(n\).
\end{proof}

\begin{theorem}[Positive-density exact MCA staircase]
\label{thm:intro-positive-density}
Let \(C=\RS_\F(D,k)\), put \(R=n-k\), and let
\(\varnothing\ne\Gamma\subseteq\F\).
\begin{enumerate}[label=\textup{(\roman*)},leftmargin=*]
\item If \(r\in\Z\) and
\[
 3r\ge R+1,\qquad 2r<R,\qquad
 n(R-r)<r(r+2)(R-2r),                          \tag{1.5}
 \label{eq:intro-literature-completion}
\]
then
\[
 B_{C,\Gamma}^{\rm MCA}(n-r)=\min\{\abs\Gamma,r+1\}. \tag{1.6}
 \label{eq:intro-literature-exact}
\]
Together with the quadratic theorem, this gives an exact finite
staircase from radius zero through the union of the two displayed ranges.
\item If \(k_n/n\to\rho\in(0,1)\) and the integers
\(r_n\) satisfy \(r_n/n\to\tau\) with
\(0<\tau<(1-\rho)/2\), then
\[
 B_{C_n,\Gamma_n}^{\rm MCA}(n-r_n)
   =\min\{\abs{\Gamma_n},r_n+1\}
\]
for every RS row sequence and every nonempty challenge sequence, once
\(n\) is sufficiently large.  Hence exact radii have normalized lower
density at least \((1-\rho)/2\) among all Hamming radii and relative density one
inside the unique-decoding interval.
\item For every fixed dyadic \(n\), the multiplicative smooth conclusion
holds over a set of primes of relative prime density
\(1/\varphi(n)\).  Standard bivariate circle codes of odd dimension
over \(\F_{p^2}\), indexed by base primes \(p\), occur over a prime set
of density \(1/\varphi(2n)\), and
Chebyshev-smooth circle \(x\)-domain RS rows occur over a prime set of
density \(1/\varphi(4n)\).  These are fixed-length densities, not a
positive-density growing prime-field family.
\item At \(n=32,k=15,r=6\), the quadratic inequality is strict enough to
give \(B^{\rm MCA}_{C,\Gamma}(26)=\min\{\abs\Gamma,7\}\), although
\(3r>R\) and the specialized estimate \eqref{eq:bchks-specialized}
does not certify this value.  In
particular it gives smooth subgroup rows for every
\(p\equiv1\pmod {32}\), a prime set of density \(1/16\), and standard
circle cosets for every \(p\equiv-1\pmod {64}\), a prime set of density
\(1/32\).
\end{enumerate}
\end{theorem}

\begin{proof}
Part~\textup{(i)} is \cref{thm:literature-integral-completion};
part~\textup{(ii)} is \cref{cor:positive-radius-density}; and
parts~\textup{(iii)}--\textup{(iv)} are
\cref{thm:fixed-length-prime-density,cor:circle-32-family}.
\end{proof}

\begin{remark}[Meaning and limitation of density]
The word ``density'' in \cref{thm:intro-positive-density} is used in two
explicit senses: normalized density of radii in part~\textup{(ii)} and
relative density inside the primes in part~\textup{(iii)}.  It does not
assert that a positive proportion of slopes is bad near capacity.
Neither the presently available profile-envelope estimates nor the
literature currently supplies such an all-received-line near-capacity
theorem.  Moreover the fixed-length prime densities tend to zero with
\(n\); \cref{prop:no-growing-prime-density} proves that this loss cannot
be removed for growing dyadic prime-field domains.
\end{remark}

\begin{theorem}[Certified Proximity-Prize rows at all four rates]
\label{thm:intro-prize-rows}
Fix the affine MCA sampler of \cite{ABF26}, in which
\(\gamma\) is uniform in \(\F\), and the target
\(\eps^*=2^{-128}\).  For each
\(\rho\in\{1/2,1/4,1/8,1/16\}\) there is an explicitly certified prime
\(p<2^{256}\), a power of two \(n\), and the multiplicative subgroup
\(D\le\F_p^*\) of order \(n\), such that \(k=\rho n=2^{40}\) and the
complete real-radius safe set of
\(C=\RS_{\F_p}(D,k)\) is
\[
 \{\delta\in[0,1]:e_{\rm MCA}(C,\delta)\le2^{-128}\}
       =\left[0,\frac Bn\right),
 \qquad B=\floor{p\,2^{-128}}.
 \tag{P1}\label{eq:intro-prize-safe-set}
\]
The endpoint \(B/n\) is unsafe.  Equivalently, the largest safe integer
radius is \(B-1\), the first unsafe integer radius is \(B\), and
\[
 B_C^{\rm MCA}(n-B+1)=B,
 \qquad B_C^{\rm MCA}(n-B)\ge B+1.
 \tag{P2}\label{eq:intro-prize-adjacent}
\]
The four transition numerators are
\[
\begin{array}{c|c|c|c}
 \rho&n&B&\delta_{C,\mathrm{sup}}^*(2^{-128})\\ \hline
 1/2&2^{41}&389500552609&389500552609/2^{41}\\
 1/4&2^{42}&1210584858040&1210584858040/2^{42}\\
 1/8&2^{43}&2879806199253&2879806199253/2^{43}\\
 1/16&2^{44}&6233898019554&6233898019554/2^{44}.
\end{array}
\tag{P3}\label{eq:intro-prize-table}
\]
These are smooth-domain rows in the literal sense of the challenge: the
evaluation group has power-of-two order, the dimension is at the stated
upper limit, and the field has fewer than \(2^{256}\) elements.
The numerical consequences are also implied by
\cite[Cor.~1.4]{BCHKS25} together with the below-half decomposition;
their role here is to provide
independent exact certificates and target bookkeeping, not a priority
claim.
\end{theorem}

\begin{proof}
The mean-overlap staircase
\cref{thm:prize-rate-overlap} gives the exact safe numerator
\(B=r_\rho(n)+1\) at radius \(B-1\), while the universal tangent
floor gives at least \(B+1\) slopes at radius \(B\).  The explicit primes,
their Proth certificates, and the exact budget divisions are given and
checked in \cref{tab:prize-proth-rows,prop:proth-row-check}.  Monotonicity
and \(r(\delta)=\floor{\delta n}\) then give
\eqref{eq:intro-prize-safe-set}.  The displayed fractions are the four
values of \(B/n\).
\end{proof}

\begin{corollary}[An extension-field smooth row]
\label{cor:intro-f17-row}
Let \(D\le\F_{17^{32}}^*\) have order \(512\) and let
\(C=\RS_{\F_{17^{32}}}(D,256)\).  At target \(2^{-128}\),
\[
 \{\delta\in[0,1]:e_{\rm MCA}(C,\delta)\le2^{-128}\}
   =[0,3/256),
\]
the largest safe integer radius is \(5\), and radius \(6\) is the first
unsafe one.  In particular the supremal transition is \(3/256\), and it
is not itself safe.
\end{corollary}

\begin{proof}
Here
\(17^{32}=2367911594760467245844106297320951247361\) and
\(\floor{17^{32}2^{-128}}=6\).  Since
\(3(6-1)\le512-256\),
\cref{cor:exact-deep-numerator} gives numerator \(6\) at radius \(5\),
while \cref{prop:universal-tangent-floor} gives at least \(7\) bad slopes
at radius \(6\).  The endpoint statement follows from the closed
integer-ball convention.
\end{proof}

\begin{theorem}[Certificate threshold formula]\label{thm:intro-delta-formula}
Let \(C_n=\RS_{\F_n}(D_n,k_n)\), let
\(\varnothing\ne\Gamma_n\subseteq\F_n\), and suppose \(k_n/n\to\rho\).
Let \(B_n^*=\lfloor\eps_n|\Gamma_n|\rfloor\), and assume the first safe
agreement \(a_n^*\) in \eqref{eq:intro-asymptotic-threshold} exists.
Suppose there is a sequence \(g_n\in[0,1-k_n/n]\) such that, for every
fixed \(\eta>0\), all admissible integer agreements
\[
 a\ge k_n+1+(g_n+\eta)n
\]
are eventually safe and all admissible integer agreements
\[
 a\le k_n+1+(g_n-\eta)n
\]
are eventually unsafe.  Then
\[
 a_n^*=k_n+1+g_n n+o(n),
 \qquad
 \delta_n^*=1-\rho-g_n+o(1).
 \tag{1.7}\label{eq:intro-delta-formula}
\]

Assume in addition that \(D_n\subseteq\B_n\subseteq\F_n\) with
\(\B_n\) a finite subfield and \(\log|\B_n|=o(n)\), and put
\(\beta_n=\log_2|\B_n|\) and
\(\tau_n=n^{-1}\log_2(1+B_n^*)\), and define the target-adjusted
identity crossing
\[
 g_{T,n}=\sup\Bigl(
 \{g\in[0,1-k_n/n]:
 H_2(k_n/n+g)-\beta_ng\ge\tau_n\}\cup\{0\}\Bigr).
 \tag{1.8}\label{eq:intro-target-crossing}
\]
If matching upper and lower certificates satisfy the preceding safe and
unsafe conditions with \(g_n=g_{T,n}+o(1)\), then
\[
 \delta_n^*=1-\rho-g_{T,n}+o(1).
 \tag{1.9}\label{eq:intro-identity-formula}
\]
When \(\tau_n=o(1)\) and \(\rho_n=k_n/n\to\rho\in(0,1)\), put
\[
 g_{0,n}=\sup\{g\in[0,1-\rho_n]:
 H_2(\rho_n+g)\ge\beta_ng\}.
 \tag{1.10}\label{eq:intro-zero-target-crossing}
\]
Then \(g_{T,n}=g_{0,n}+o(1)\), so
\(\delta_n^*=1-\rho-g_{0,n}+o(1)\).
\end{theorem}

\begin{proof}
For every fixed \(\eta>0\), monotonicity places the first safe agreement
between
\[
 k_n+1+(g_n-\eta)n+O(1)
 \quad\text{and}\quad
 k_n+1+(g_n+\eta)n+O(1).
\]
Divide by \(n\), let \(n\to\infty\), and then let \(\eta\downarrow0\).
Since \(\delta_n^*=1-a_n^*/n\), this proves
\eqref{eq:intro-delta-formula}.  The second assertion is the same
argument with the displayed crossing substituted for \(g_n\).

For completeness, the identity-prefix mean at
\(a=k_n+1+gn+O(1)\) is
\[
 \binom na|\B_n|^{-(a-k_n-1)}.
\]
Stirling's formula gives its normalized logarithm as
\[
 \frac1n\log_2\!\left(
 \binom na|\B_n|^{-(a-k_n-1)}\right)
 =H_2(\rho_n+g)-\beta_ng+o(1),
\]
uniformly on compact subintervals of \(0<\rho_n+g<1\).  This explains
\eqref{eq:intro-target-crossing}; it does not supply the assumed upper
certificate.  Finally,
\(f_n(g)=H_2(\rho_n+g)-\beta_ng\) is strictly concave,
\(f_n(0)=H_2(\rho_n)>0\), and \(f_n(1-\rho_n)<0\).  Eventually
\(\tau_n<f_n(0)\).  Hence the \(\tau_n\)-superlevel set is an interval
containing zero, with right endpoint \(g_{T,n}\), whereas the
zero-superlevel set has right endpoint \(g_{0,n}\); in particular,
\(0\le g_{0,n}-g_{T,n}\).  Fix \(\eta>0\).  If \(g_{0,n}<\eta\), this
difference is already less than \(\eta\).  Otherwise, if
\(g_{0,n}-g_{T,n}\ge\eta\), concavity on \([0,g_{0,n}]\) gives
\[
 f_n(g_{0,n}-\eta)
 \ge \frac{\eta}{g_{0,n}}f_n(0)
 \ge \eta H_2(\rho_n).
\]
The left side is at most \(\tau_n\), by the definition of the right
\(\tau_n\)-crossing, while the right side is bounded away from zero.
This contradicts \(\tau_n=o(1)\).  Hence
\(g_{T,n}=g_{0,n}+o(1)\).
\end{proof}

\begin{remark}[Scope of the entropy crossing]\label{rem:intro-delta-scope}
The first part of \cref{thm:intro-delta-formula} is unconditional once
matching certificates have been proved.  The identity crossing in
\eqref{eq:intro-identity-formula} is a specialization, not an
identity-only theorem for every smooth or circle domain.  Quotient,
subfield, planted, and determinant-ray profiles can have a larger
realized-image scale; in such a row \(g_n\) is the crossing of the full
certified profile envelope.  Support or support-pair estimates also require
a separate projection to distinct slopes.  These are precisely the two
points at which an unrestricted near-capacity claim would need additional
input.
\end{remark}



\section{MCA witnesses and exact support reduction}\label{sec:witnesses}

Throughout this section \(\F\) is a finite field, \(D\subseteq\F\) has
size \(n\), and \(C=\RS_\F(D,k)\).  Write
\(\F[X]_{<k}=\{f\in\F[X]:\deg f<k\}\), and let
\(\operatorname{ev}_S:\F[X]_{<k}\to\F^S\) be restriction of the
evaluation map.

\begin{definition}[Explanation and support-wise nontriviality]
\label{def:explanation}
For \(u\in\F^D\) and \(S\subseteq D\), an \emph{explanation of \(u\) on
\(S\)} is a polynomial \(f\in\F[X]_{<k}\) with
\(\operatorname{ev}_S(f)=u|_S\).  We write
\(\operatorname{Expl}_C(u;S)\) when such an \(f\) exists.  A pair
\(r=(r_0,r_1)\) is \emph{simultaneously explained on \(S\)} when
\(\operatorname{Expl}_C(r_0;S)\) and
\(\operatorname{Expl}_C(r_1;S)\) both hold, with no requirement that the
two explaining polynomials coincide.  It is \emph{support-wise
nontrivial on \(S\)} when it is not simultaneously explained there.
We denote this last condition by \(\operatorname{NT}_C(r;S)\).
\end{definition}

Thus support-wise nontriviality is a condition on the same support that
explains the affine combination.  It does not mean that both coordinates
are individually unexplained; failure of either explanation is enough.

\begin{definition}[MCA-bad slope]\label{def:mca-bad}
Let \(r=(r_0,r_1)\in(\F^D)^2\) and let \(a\) be an agreement threshold.
A finite slope \(\gamma\in\F\) is support-wise MCA-bad at agreement \(a\)
if there are \(S\subseteq D\), \(\abs S\ge a\), and
\(h\in\F[X]_{<k}\) such that
\(\operatorname{ev}_S(h)=(r_0+\gamma r_1)|_S\) and
\(\operatorname{NT}_C(r;S)\) holds.  Define
\[
        B_C^{\rm MCA}(a)
        =\max_{r\in(\F^D)^2}\abs{\{\gamma\in\F:
          \gamma\text{ is MCA-bad for }r\text{ at }a\}}.
\]
\end{definition}

The official MCA maximum mixes two logically different sources: received
pairs that are globally far from every codeword pair, and pairs that are
close to one codeword pair but fail on a slope-dependent support.  The
following definitions and theorem separate them exactly.

\begin{definition}[CA numerator and sparse mutual numerator]
\label{def:ca-sparse-numerators}
Assume \(0\le a\le n\) and put \(t=n-a\).  A received pair
\((r_0,r_1)\) is \emph{column-far at
agreement \(a\)} if it is not simultaneously explained on any support of
size at least \(a\).  Let \(B_C^{\rm CA}(a)\) be the maximum, over
column-far pairs, of the number of slopes \(\gamma\) for which
\(r_0+\gamma r_1\) is explained on some support of size at least \(a\).
If there is no column-far pair, this maximum is defined to be zero.
Define the \emph{sparse mutual numerator}
\[
 \sigma_C(a)=\max\left\{
 \#\{\gamma:\gamma\text{ is MCA-bad for }(e_0,e_1)\}:
 \abs{\supp(e_0)\cup\supp(e_1)}\le t\right\}.
 \tag{SP1}\label{eq:sparse-mutual-numerator}
\]
Thus \(B_C^{\rm CA}(a)/\abs\F\) is the ordinary affine-line correlated
agreement error on the integer grid, while \(\sigma_C(a)\) records the
purely mutual obstruction after translation by a nearby codeword pair.
\end{definition}

\begin{theorem}[Exact sparsification of the mutual layer]
\label{thm:exact-sparsification}
For every linear code \(C\subseteq\F^D\) and every agreement threshold
\(0\le a\le n\),
\[
 B_C^{\rm MCA}(a)=
 \max\{B_C^{\rm CA}(a),\sigma_C(a)\}.              \tag{SP2}
\label{eq:exact-sparsification}
\]
Equivalently, after division by \(\abs\F\), the official MCA error is the
maximum of the CA error and the sparse mutual error.
\end{theorem}

\begin{proof}
Both quantities on the right are obtained by maximizing the MCA-bad set
over restricted classes of pairs, so the left side is at least their
maximum.  Conversely fix a received pair.  If it is column-far, every
support that explains a line point necessarily fails to explain the pair;
its MCA-bad set is therefore its CA-bad set and has size at most
\(B_C^{\rm CA}(a)\).

Otherwise choose codewords \(c_0,c_1\in C\) and a common support of size
at least \(a\), and put \(e_i=r_i-c_i\).  The union of the supports of
\(e_0,e_1\) has size at most \(n-a=t\).  At each slope, translation sends
an explaining codeword \(h\) to
\(h-c_0-\gamma c_1\), and is a bijection between witness supports for
\((r_0,r_1)\) and for \((e_0,e_1)\).  It also preserves simultaneous
explanation on that support.  Hence the two pairs have exactly the same
MCA-bad slopes, and their number is at most \(\sigma_C(a)\).
\end{proof}

For a nonempty challenge set \(\Gamma\subseteq\F\), write
\(B_{C,\Gamma}^{\rm CA}(a)\) and \(\sigma_{C,\Gamma}(a)\) for the same
maxima with the bad slopes restricted to \(\Gamma\).  The proof just given
is unchanged after intersecting every bad-slope set with \(\Gamma\), and
therefore gives the exact challenge-restricted identity
\[
 B_{C,\Gamma}^{\rm MCA}(a)
 =\max\{B_{C,\Gamma}^{\rm CA}(a),\sigma_{C,\Gamma}(a)\}.
 \tag{SP3}\label{eq:challenge-sparsification}
\]

\begin{theorem}[Exact mutual layer below half the distance]
\label{thm:exact-half-distance-sparse}
Let \(C\subseteq\F^D\) be a linear code of minimum distance
\(d_{\min}\), let \(\varnothing\ne\Gamma\subseteq\F\), and let
\(0\le r\le n\) satisfy \(2r\le d_{\min}-1\).  Then
\[
 \sigma_{C,\Gamma}(n-r)=\min\{\abs\Gamma,r\},
 \qquad
 B_{C,\Gamma}^{\rm MCA}(n-r)
 =\max\left\{B_{C,\Gamma}^{\rm CA}(n-r),
              \min\{\abs\Gamma,r\}\right\}.       \tag{HD1}
\label{eq:exact-half-distance-sparse}
\]
For a Reed--Solomon row the hypothesis is \(2r\le n-k\).
\end{theorem}

\begin{proof}
Fix a sparse pair \((e_0,e_1)\) whose support union \(E\) has size at
most \(r\), and let \(\gamma\) be bad on a support \(S\) of size at least
\(n-r\), explained there by a codeword \(c\in C\).  On \(S\setminus E\)
the sparse line point is zero, so
\(\supp(c)\subseteq E\cup(D\setminus S)\) has size at most \(2r<d_{\min}\).
Thus \(c=0\).  If \(S\cap E\) contained no nonzero pair coordinate, the
zero codeword pair would simultaneously explain \((e_0,e_1)\) on \(S\),
contrary to badness.  Hence some \(x\in S\cap E\) satisfies
\[
 (e_0(x),e_1(x))\ne(0,0),\qquad e_0(x)+\gamma e_1(x)=0.
\]
Necessarily \(e_1(x)\ne0\), and this coordinate determines the single
slope \(-e_0(x)/e_1(x)\).  Therefore every sparse pair has at most
\(\min\{\abs\Gamma,\abs E\}\le\min\{\abs\Gamma,r\}\) bad slopes.

For the reverse bound put \(t=\min\{\abs\Gamma,r\}\), choose distinct
\(x_1,\ldots,x_t\in D\) and distinct
\(\gamma_1,\ldots,\gamma_t\in\Gamma\), and set
\(e_1(x_i)=1\), \(e_0(x_i)=-\gamma_i\), with both words zero elsewhere.
For \(S_i=D\setminus\{x_j:j\ne i\}\), the line point of slope
\(\gamma_i\) is zero on \(S_i\), and
\(\abs{S_i}=n-t+1\ge n-r\).  The restriction of \(e_1\) to \(S_i\)
cannot be explained by a codeword: such a codeword would be nonzero at
\(x_i\) and zero on \(n-t\) coordinates, hence would have weight at most
\(t<d_{\min}\).  Thus every \(\gamma_i\) is MCA-bad and equality holds.
The second formula is \eqref{eq:challenge-sparsification}; an RS code has
\(d_{\min}=n-k+1\).
\end{proof}

\begin{lemma}[Pair-interleaved Johnson bound]
\label{lem:pair-interleaved-johnson}
Let \(C=\RS_\F(D,k)\), let \(1\le A\le n\), and suppose
\(A^2>(k-1)n\).  For every received pair \(y\in(\F^D)^2\), the number
of codeword pairs \((p_0,p_1)\in C^2\) simultaneously agreeing with
\(y\) on at least \(A\) coordinates is at most
\[
 J_2(A)=\floor{
   \frac{n(A-k+1)}{A^2-(k-1)n}}.                  \tag{HD2}
 \label{eq:pair-johnson}
\]
\end{lemma}

\begin{proof}
Let the list have size \(L\), and for every member retain exactly \(A\)
agreement coordinates.  If \(d_x\) is the number of retained sets
containing \(x\), then \(\sum_xd_x=LA\).  Two distinct codeword pairs
can agree coordinatewise on at most \(k-1\) points, since at least one
nonzero difference polynomial has degree less than \(k\).  Therefore
\[
 \frac{L^2A^2/n-LA}{2}
 \le\sum_{x\in D}\binom{d_x}{2}
 \le\binom L2(k-1),
\]
where the first inequality is Cauchy--Schwarz.  Rearrangement gives
\(L(A^2-(k-1)n)\le n(A-k+1)\), and integrality gives
\eqref{eq:pair-johnson}.
\end{proof}

\begin{theorem}[Self-contained half-Johnson CA and MCA bound]
\label{thm:self-contained-half-johnson}
Let \(C=\RS_\F(D,k)\), let \(\varnothing\ne\Gamma\subseteq\F\), and
let \(r\) be an integer with \(0\le r\le n\).  Put \(A=n-2r\).  If
\(A>0\) and \(A^2>(k-1)n\), then
\[
 B_{C,\Gamma}^{\rm CA}(n-r)
 \le\min\{\abs\Gamma,1+(r+1)J_2(A)\}.             \tag{HD3}
 \label{eq:half-johnson-ca}
\]
If in addition \(2r\le n-k\), the same upper bound holds for
\(B_{C,\Gamma}^{\rm MCA}(n-r)\).
\end{theorem}

\begin{proof}
Fix a column-far received pair and let \(Z\subseteq\Gamma\) be its
CA-bad set.  The assertion is immediate when \(\abs Z\le1\), so fix an
anchor \(\gamma_0\in Z\).  For every
\(\gamma\in Z\setminus\{\gamma_0\}\), choose explaining codewords
\(c_{\gamma_0},c_\gamma\) and agreement sets of size at least \(n-r\),
and define
\[
 P_1=\frac{c_{\gamma_0}-c_\gamma}{\gamma_0-\gamma},
 \qquad P_0=c_{\gamma_0}-\gamma_0P_1.
\]
On the intersection of the two agreement sets, of size at least
\(n-2r=A\), the codeword pair \((P_0,P_1)\) agrees with the received
pair.  Hence all such pairs lie in a list of size at most \(J_2(A)\) by
\cref{lem:pair-interleaved-johnson}.

Fix one pair in that list and let \(E\) be its column-error set, of size
\(e\).  Column-farness gives \(e\ge r+1\).  If a slope \(\gamma\) rides
this pair, meaning its explaining codeword is \(P_0+\gamma P_1\), then
the nonzero affine coordinate errors vanish at at least \(e-r\)
coordinates of \(E\).  Each coordinate vanishes for at most one finite
slope, so double counting gives at most
\(e/(e-r)\le r+1\) riders.  Summing over the pair list and adding the
anchor proves \eqref{eq:half-johnson-ca}.  Under \(2r\le n-k\),
\cref{thm:exact-half-distance-sparse} says that the other branch of MCA
has only \(\min\{\abs\Gamma,r\}\) slopes.  Since
\(n(A-k+1)-(A^2-(k-1)n)=A(n-A)\ge0\), one has \(J_2(A)\ge1\), so this
tangent term is no larger than the displayed CA bound.
\end{proof}

\begin{remark}[The exact remaining prize certificate]
\label{rem:exact-prize-certificate}
At a proposed safe agreement \(a_+\), an unrestricted finite-row proof
must establish both
\[
 B_C^{\rm CA}(a_+)\le B^*,\qquad
 \sigma_C(a_+)\le B^*,\qquad
 B^*=\floor{2^{-128}\abs\F}.
\]
An unsafe construction at \(a_+-1\) does not prove either upper bound.
This identity is the shortest exact description of what remains for the
near-capacity rows not covered by
\cref{cor:half-distance-window-compiler}; the older quadratic compiler
\cref{cor:prize-window-compiler} is a self-contained subrange.
\end{remark}

The finite-slope restriction costs at most the point at infinity, which may
also be treated by exchanging the two coordinates.  The support-wise
condition is essential: explanation and nontriviality are tested on the
same set.

\begin{definition}[Exact-agreement witness incidence]
\label{def:exact-witness-incidence}
Assume \(k+1\le a\le n\).  For a received line \(r=(r_0,r_1)\), its
exact-\(a\) witness incidence is
\[
 \Wcal_a(r)=\left\{(\gamma,S,h):
 \begin{array}{l}
   \gamma\in\F,\quad S\in\binom Da,\quad h\in\F[X]_{<k},\\
   \operatorname{ev}_S(h)=(r_0+\gamma r_1)|_S,\quad
   \operatorname{NT}_C(r;S)
 \end{array}\right\}.
 \tag{2.1}\label{eq:exact-witness-incidence}
\]
Its slope projection is
\(\pi_\gamma(\gamma,S,h)=\gamma\), and its exact bad-slope set is
\(Z_a(r)=\pi_\gamma(\Wcal_a(r))\subseteq\F\).
\end{definition}

\begin{lemma}[Reduction to exact agreement]
\label{lem:exact-agreement-reduction}
For \(a\ge k+1\), the set \(Z_a(r)\) is precisely the set of slopes that
are MCA-bad at threshold \(a\) in \cref{def:mca-bad}.
\end{lemma}

\begin{proof}
An exact-\(a\) witness is a threshold-\(a\) witness.  Conversely, let
\((\gamma,S,h)\) witness badness with \(\abs S\ge a\).  At least one of
\(r_0|_S,r_1|_S\) is outside
\(\operatorname{ev}_S(\F[X]_{<k})\); call it \(u|_S\).  Interpolate \(u\)
on any fixed \(k\)-subset \(K\subseteq S\).  If the resulting polynomial
agreed with \(u\) at every point of \(S\setminus K\), it would explain
\(u\) on \(S\), a contradiction.  Hence some \(x\in S\setminus K\)
makes \(u\) unexplained on \(K\cup\{x\}\).  Extend this set to an
\(a\)-subset \(S'\subseteq S\).  Non-explanation persists on supersets,
while the same \(h\) explains \(r_0+\gamma r_1\) on \(S'\).
\end{proof}

\begin{proposition}[Exact support-atlas upper bound]
\label{prop:exact-support-upper}
For every Reed--Solomon row, every nonempty challenge set
\(\Gamma\subseteq\F\), and every \(k+1\le a\le n\),
\[
 B_{C,\Gamma}^{\rm MCA}(a)
 \le\min\left\{\abs\Gamma,\binom na\right\}.        \tag{2.2}
\]
This is a literal finite-row bound: it has no polynomial or
subexponential loss.
\end{proposition}

\begin{proof}
By \cref{lem:exact-agreement-reduction}, choose one noncommon exact-\(a\)
support for every bad slope.  One noncommon support carries at most one
slope, since two distinct explained slopes would separately explain the
two received words on that support.  The chosen-support map is therefore
injective.  Intersect with the challenge set.
\end{proof}

\section{Exact syndrome geometry and the deep regime}
\label{sec:syndrome-exact}

The \emph{syndrome} of a word is the result of applying a parity-check
matrix; it vanishes exactly on codewords.  Thus syndrome space forgets the
part of a received word already lying in the code and retains only its error
class.  This description separates possible error supports from distinct
line parameters.  Put \(R=n-k\), the redundancy (and syndrome-space
dimension), and choose the generalized Vandermonde parity-check matrix
\(H\) with columns
\[
 h_x=\lambda_x(1,x,\ldots,x^{R-1})^{\mathsf T},\qquad
 \lambda_x=\left(\prod_{y\in D\setminus\{x\}}(x-y)\right)^{-1}.
 \tag{3.1}\label{eq:parity-columns}
\]
Lagrange interpolation gives
\(\sum_{x\in D}\lambda_xx^j=0\) for \(0\le j\le n-2\), so this matrix
annihilates every degree-less-than-\(k\) evaluation.  Every set of at most
\(R\) columns is independent.  For \(E\subseteq D\), write
\(V_E=\Span\{h_x:x\in E\}\); it is exactly the set of syndromes of words
supported on \(E\).  We write \(s(u)=Hu^{\mathsf T}\) for the syndrome of
\(u\).

\begin{proposition}[Syndrome-line normal form]
\label{prop:syndrome-line-normal-form}
Let \(S=D\setminus E\).  A word \(u\) is explained on \(S\) if and only if
\(s(u)\in V_E\).  Consequently \(\gamma\) is MCA-bad on \(S\) for a pair
\((u_0,u_1)\) if and only if
\[
 s(u_0)+\gamma s(u_1)\in V_E,
 \qquad
 \{s(u_0),s(u_1)\}\not\subseteq V_E.
 \tag{3.2}\label{eq:syndrome-line}
\]
For fixed \(E\), there is at most one such finite slope.
\end{proposition}

\begin{proof}
If \(u=c+e\), where \(c\in C\) and \(e\) is supported on \(E\), then
\(s(u)=s(e)\in V_E\).  Conversely, if
\(s(u)=\sum_{x\in E}e_xh_x\), subtract the word supported on \(E\) with
values \(e_x\); the result has zero syndrome and belongs to \(C\).  Apply
this equivalence to the line point and to both coordinates.  Passing
\eqref{eq:syndrome-line} to \(\F^R/V_E\) shows that two distinct solutions
would put both projected syndromes at zero, contradicting badness.
\end{proof}

Since \(\lambda_x\ne0\), one has
\([h_x]=[1:x:\cdots:x^{R-1}]\); these are the points of the rational
normal curve in projective syndrome space, while the weights
\(\lambda_x\) merely choose vector representatives.  A \emph{secant span} is the linear span
\(V_E\) of a selected set of those points.  We call the meeting of a
syndrome line with \(V_E\) \emph{transverse} when the entire line is not
contained in \(V_E\), equivalently when its two generators are not both in
\(V_E\).  The next result is a \emph{compiler} in the sense that it
translates the bad-slope problem, exactly and in both directions, into this
line--secant incidence problem.  A \emph{ray} means a one-dimensional
direction in syndrome or coefficient space; different raw witnesses may
represent the same ray and hence the same slope.

\begin{theorem}[Exact syndrome--secant compiler]
\label{thm:syndrome-secant-exact}
Put \(t=n-a\), and for a syndrome line
\(\ell_{y_0,y_1}=\{y_0+\gamma y_1:\gamma\in\F\}\) define
\[
 \Theta_t(y_0,y_1)=
 \abs{\{\gamma\in\F:\exists E\subseteq D,\ \abs E\le t,
 y_0+\gamma y_1\in V_E,
 \ \{y_0,y_1\}\not\subseteq V_E\}}.
 \tag{3.3}\label{eq:transverse-secant-count}
\]
Then
\[
 B_C^{\rm MCA}(a)=\max_{y_0,y_1\in\F^{n-k}}
                    \Theta_t(y_0,y_1).
 \tag{3.4}\label{eq:exact-secant-numerator}
\]
For each fixed \(E\), the affine line has at most one transverse
intersection parameter with \(V_E\).  For a Reed--Solomon code the
spaces \(V_E\) are the spans of at most \(t\) points of the weighted
rational normal curve in \eqref{eq:parity-columns}.  Thus the
higher-dimensional ray problem is exactly a transverse line--secant
incidence problem, with repeated intersection parameters deduplicated.
\end{theorem}

\begin{proof}
Apply \cref{prop:syndrome-line-normal-form} with \(S=D\setminus E\), and
take the union over \(\abs E\le t\).  The syndrome map is surjective, so
maximizing over received pairs is the same as maximizing over
\((y_0,y_1)\).  In the quotient \(\F^{n-k}/V_E\), a transverse
intersection solves
\(\overline y_0+\gamma\overline y_1=0\) with the two projected vectors
not both zero, and therefore has at most one finite solution.  The last
assertion is the Vandermonde description of the parity columns.
\end{proof}

A \emph{parity-column circuit} is a minimally linearly dependent set of
columns of \(H\).  The next lemma says that many distinct transverse rays
cannot all be supported inside an independent union of columns.

\begin{lemma}[Independent-union ray bound]
\label{lem:independent-union-rays}
Let every set of at most \(R\) parity columns be independent, put
\(t<R\), and fix a syndrome line.  For each slope in a transverse set
\(Z\), choose an error set \(E_\gamma\), \(\abs{E_\gamma}\le t\),
witnessing \eqref{eq:transverse-secant-count}.  If
\(U=\bigcup_{\gamma\in Z}E_\gamma\) has size at most \(R\), then
\(\abs Z\le t+1\).  Consequently, more than \(t+1\) transverse slopes
force the witness union to contain a parity-column circuit.
\end{lemma}

\begin{proof}
If \(\abs Z\le1\), there is nothing to prove; assume that \(Z\) contains
two distinct slopes.  Two distinct line points lie in \(V_U\), hence
\(y_0,y_1\in V_U\).
Column independence gives unique coefficient lifts
\(e_0,e_1\in\F^U\), and every chosen lift is
\(e_\gamma=e_0+\gamma e_1\).  Let \(U'\) be the set of \(s\) coordinates
where the pair \((e_0,e_1)\) is nonzero.  If \(s\le t\) and
\(E_\gamma\) contained all of \(U'\), then
\(y_0,y_1\in V_{U'}\subseteq V_{E_\gamma}\), contradicting the
transversality of this particular witness.  Hence some
\(x\in U'\setminus E_\gamma\) exists, and uniqueness of the lift gives
\(e_0(x)+\gamma e_1(x)=0\).  Each nonzero affine coordinate function has
at most one root, so \(\abs Z\le s\le t\).  If \(s>t\), every chosen lift of
weight at most \(t\) has at least \(s-t\) zero coordinates.  Double
counting the pairs \((\gamma,x)\) with
\(e_0(x)+\gamma e_1(x)=0\) gives
\(\abs Z(s-t)\le s\), whence
\(\abs Z\le\floor{s/(s-t)}\le t+1\).  If \(\abs U\le R\), independence
holds; the contrapositive proves the circuit assertion.
\end{proof}

\begin{theorem}[Bounded residual-kernel ray compiler]
\label{thm:bounded-residual-kernel-ray}
Let \(H\) be a parity-check matrix over \(\F\), let \(U\) be a set of
columns, and put
\[
 \kappa(U)=\dim_\F\ker(H_U:\F^U\longrightarrow\F^{n-k}).
\]
Fix a syndrome line \(y_0+\gamma y_1\), an integer \(t\ge0\), and a set
\(Z\) of finite transverse slopes.  Suppose that for every \(\gamma\in Z\)
there is \(E_\gamma\subseteq U\), \(\abs{E_\gamma}\le t\), such that
\[
 y_0+\gamma y_1\in V_{E_\gamma},\qquad
 \{y_0,y_1\}\not\subseteq V_{E_\gamma}.
\]
Then
\[
                  \abs Z\le(t+1)\abs{\F}^{\kappa(U)}. \tag{RC$_{\rm ker}$}
\]
For a Reed--Solomon parity check of redundancy \(R=n-k\),
\(\kappa(U)=\max\{0,\abs U-R\}\).  Hence a certified residual chart with
\((\abs U-R)_+\log\abs\F=o(n)\) has a direct subexponential ray payment
with the displayed exact finite loss.
\end{theorem}

\begin{proof}
The assertion is immediate when \(\abs Z\le1\).  Otherwise two line
points belong to \(V_U\), so \(y_0,y_1\in V_U\).  Choose lifts
\(b_0,b_1\in\F^U\) with \(H_Ub_i=y_i\).  For every \(\gamma\), choose a
lift \(c_\gamma\) supported in \(E_\gamma\) and put
\(z_\gamma=c_\gamma-b_0-\gamma b_1\in\ker H_U\).

Fix \(z\in\ker H_U\), and consider the slopes for which
\(z_\gamma=z\).  Put
\[
 U_z=\{x\in U:(b_0(x)+z(x),b_1(x))\ne(0,0)\},
 \qquad s=\abs{U_z}.
\]
Every active coordinate is a nonzero affine function of \(\gamma\), and
therefore vanishes for at most one slope.  Hence the total number of active
zero incidences in this kernel class is at most \(s\).  If \(s>t\), the
weight bound on \(c_\gamma\) forces at least \(s-t\) active zeros, so
\[
 \abs{Z_z}(s-t)\le s,\qquad
 \abs{Z_z}\le\floor{s/(s-t)}\le t+1.
\]
If \(s\le t\), transversality forces
at least one active zero: without one,
\(U_z=\supp c_\gamma\subseteq E_\gamma\), while both \(b_0+z\) and
\(b_1\) are supported in \(U_z\), putting \(y_0,y_1\) in
\(V_{E_\gamma}\).  Thus this kernel class contains at most \(s\le t\)
slopes.  Summing over the \(\abs{\F}^{\kappa(U)}\) kernel elements proves
the bound.  The MDS column property gives the Reed--Solomon formula for
\(\kappa(U)\).
\end{proof}

\begin{theorem}[Single MDS-circuit ray compiler]
\label{thm:single-mds-circuit-ray}
Let \(H\) be a Reed--Solomon parity check of redundancy \(R\), let
\(U\subseteq D\) have size \(R+1\), put \(t<R\), and fix a syndrome line
\(y_0+\gamma y_1\).  If every slope in a transverse set \(Z\) has a witnessing error set
\(E_\gamma\subseteq U\) of size at most \(t\), then
\[
                         \abs Z\le\binom{R+1}{2}.   \tag{RC$_{\rm circ}$}
\]
Thus one fixed MDS-circuit profile has a field-independent polynomial ray
payment.
\end{theorem}

\begin{proof}
There is nothing to prove when \(\abs Z\le1\).  Otherwise two distinct
line points lie in \(V_U\); subtracting them and solving for the constant
term gives \(y_0,y_1\in V_U\).  Choose lifts
\(b_0,b_1\in\F^U\).  The circuit kernel is generated by a vector
\(z\) whose coordinates are all nonzero, since a zero coordinate would
give a proper dependent subset of the minimal circuit \(U\).  For each slope choose
\(c_\gamma\in\F\) so that
\(e_\gamma=b_0+\gamma b_1+c_\gamma z\) is supported in
\(E_\gamma\).  For \(x\in U\), set
\[
             f_x(\gamma)=-(b_0(x)+\gamma b_1(x))/z(x).
\]
A zero coordinate is an equality \(c_\gamma=f_x(\gamma)\), and at least
\(R+1-t\ge2\) such equalities hold.

Every transverse slope makes two nonidentical functions \(f_x,f_{x'}\)
equal.  Indeed, otherwise all vanishing coordinates belong to one class
\(C\) of identical affine functions, say
\(f_x=\alpha+\beta\gamma\) on \(C\), and no function outside \(C\) has
that value at \(\gamma\).  Then
\(\supp e_\gamma=U\setminus C\subseteq E_\gamma\), while
\(b_0+\alpha z\) and \(b_1+\beta z\) are also supported in
\(U\setminus C\).  This puts both \(y_0,y_1\) in \(V_{E_\gamma}\), a
contradiction.  Two nonidentical affine functions agree at at most one
finite slope.  Charging each slope to one unordered pair proves the
displayed bound.
\end{proof}



\section{The mean-overlap regime}\label{sec:mean-overlap}

For a fixed received line, the exact bad-slope set is therefore the set of
unique cancelling line parameters contributed by the transverse quotient
equations.  For each \(E\), \(\abs E\le n-a\), project the two
syndromes to \(\F^R/V_E\).  If the second projection is nonzero and the
first is a scalar multiple of it, record the unique cancelling scalar;
otherwise record nothing.  Deduplicating these scalars gives the numerator.
In particular, the number of supports, the number of pairs of supports, and
the number of distinct slopes are three different quantities.

We call \(3(n-a)\le d-1\) the \emph{deep regime}: the allowed error
radius \(n-a\) is so small relative to the minimum codeword weight \(d\)
that the union of three permitted error supports still cannot contain a
nonzero codeword support.

\begin{theorem}[Deep-regime upper bound]\label{thm:deep-regime-upper}
Let \(C\subseteq\F^n\) be linear with minimum nonzero weight \(d\), let
\(0\le a\le n\), put \(r=n-a\), and assume \(3r\le d-1\).  Then
\[
 B_C^{\rm MCA}(a)\le r+1.
 \tag{3.5}\label{eq:deep-upper}
\]
\end{theorem}

\begin{proof}
Fix \((u_0,u_1)\), and let \(Z\) be the slopes for which \(u_0+\gamma u_1\)
is within distance \(r\) of \(C\).  If \(\abs Z\le r+1\), there is nothing
to prove.  Otherwise choose distinct \(\gamma_1,\gamma_2\in Z\), codewords
\(p_1,p_2\), and error sets \(E_1,E_2\), each of size at most \(r\), such
that \(u_0+\gamma_i u_1=p_i\) off \(E_i\).  Solving the two affine
equations gives \(c_0,c_1\in C\) for which the error pair
\(e_i=u_i-c_i\) is supported on \(T=E_1\cup E_2\), \(\abs T\le2r\).

For any \(\gamma\in Z\), let \(p\) be a closest explaining codeword.  The
codeword \(p-(c_0+\gamma c_1)\) is supported on at most \(3r\le d-1\)
positions and is zero.  Thus \(\wt(e_0+\gamma e_1)\le r\).  Let \(T'\) be
the support of the error pair.  If \(\abs{T'}\ge r+1\), every
\(\gamma\in Z\) makes at least \(\abs{T'}-r\) nonzero affine coordinate
functions \(e_{0,x}+\gamma e_{1,x}\) vanish.  Each coordinate has at most
one root, so \(\abs Z(\abs{T'}-r)\le\abs{T'}\) and \(\abs Z\le r+1\).

It remains that \(\abs{T'}\le r\).  If \(\gamma\) is MCA-bad on a set
\(S\), \(\abs S\ge n-r\), its explaining codeword equals
\(c_0+\gamma c_1\), since their difference would otherwise have weight at
most \(2r\le d-1\).  The set \(S\) must meet \(T'\), and at a point of
\(S\cap T'\) the nonzero pair satisfies
\(e_{0,x}+\gamma e_{1,x}=0\).  Such a point determines \(\gamma\), giving
at most \(\abs{T'}\le r\) bad slopes.
\end{proof}

Recall that for a nonempty challenge set \(\Gamma\subseteq\F\),
\(B_{C,\Gamma}^{\rm MCA}(a)\) denotes the same maximum with slopes restricted to
\(\Gamma\).

The following ``tangent floor'' is an elementary family of neighboring
coordinate-span intersections.  The name does not assert a hidden
differentiability or multiplicity condition.

\begin{proposition}[Universal tangent floor]\label{prop:universal-tangent-floor}
For every Reed--Solomon code, every \(a\ge k+1\), and every nonempty
\(\Gamma\subseteq\F\),
\[
 B_{C,\Gamma}^{\rm MCA}(a)
 \ge\min\{\abs\Gamma,n-a+1\}.
 \tag{3.6}\label{eq:tangent-floor}
\]
\end{proposition}

\begin{proof}
Put \(j=n-a\), \(t=\min\{\abs\Gamma,j+1\}\), and choose distinct
\(\gamma_1,\ldots,\gamma_t\in\Gamma\) and independent parity columns
\(h_{x_1},\ldots,h_{x_t}\).  Define
\(s_1=\sum_i h_{x_i}\) and \(s_0=-\sum_i\gamma_i h_{x_i}\), and lift them
to received words.  For \(E_i=\{x_\ell:\ell\ne i\}\),
\(s_0+\gamma_i s_1\in V_{E_i}\), whereas \(s_1\notin V_{E_i}\).
\Cref{prop:syndrome-line-normal-form} makes every \(\gamma_i\) bad on
\(D\setminus E_i\), which has size at least \(a\).
\end{proof}

\begin{corollary}[Exact deep numerator]\label{cor:exact-deep-numerator}
If \(C=\RS_\F(D,k)\), \(1\le k<n\), \(k+1\le a\le n\),
\(\varnothing\ne\Gamma\subseteq\F\), and \(3(n-a)\le n-k\), then
\[
 B_{C,\Gamma}^{\rm MCA}(a)
 =\min\{\abs\Gamma,n-a+1\}.
\]
\end{corollary}

\begin{proof}
The upper bound is \cref{thm:deep-regime-upper}, intersected with
\(\Gamma\).  The hypotheses imply \(a\ge k+1\), so
\cref{prop:universal-tangent-floor} gives equality.
\end{proof}

\begin{remark}[The zero-radius normalization]
Under the formal affine sampler, \cref{cor:exact-deep-numerator} at
\(a=n\) gives \(B_C^{\rm MCA}(n)=1\), and hence
\(e_{\rm MCA}(C,0)=1/\abs\F\) for every proper Reed--Solomon code.  This
also follows directly: in \(\F^D/C\), a noncontained affine line meets
the zero coset in at most one parameter, and the bound is attained by a
line through the zero coset.  A value \(2/\abs\F\) belongs to a different
or loose convention and must not be imported into the affine definition
of \cite{ABF26}.
\end{remark}

The next lemma isolates the deterministic mechanism behind both extensions
of the deep range.

\begin{lemma}[Large-overlap collapse]
\label{lem:large-overlap-collapse}
Fix a received pair and an exact agreement \(a=n-r\ge k+1\).  Choose one
noncommon exact witness support \(S_z\) and explaining polynomial \(p_z\)
for every slope \(z\) in its bad set \(Z\).  If two distinct chosen
supports satisfy
\[
                    \abs{S_z\cap S_w}\ge k+r,       \tag{MO1}
\label{eq:large-overlap-collapse}
\]
then \(\abs Z\le r+1\).
\end{lemma}

\begin{proof}
Put
\[
 c_1=\frac{p_z-p_w}{z-w},\qquad c_0=p_z-zc_1.
\]
On \(S_z\cap S_w\), the two line equations imply that the received pair
equals \((c_0,c_1)\).  Let \(C_0\) be the full common support of this
codeword pair and write \(c=\abs{C_0}\ge k+r\).  For every \(u\in Z\),
\(\abs{S_u\cap C_0}\ge a+c-n\ge k\), so the explaining polynomial
\(p_u\) equals \(c_0+uc_1\).

Subtract this codeword pair from the received pair.  At every coordinate
outside \(C_0\) at least one of the two residual values is nonzero, and a
coordinate can cancel at at most one finite slope.  A noncommon witness
for \(u\) must contain at least \(\max\{1,a-c\}\) such cancellations:
the support-size condition gives \(a-c\), and one is still necessary when
\(a\le c\).  Double counting therefore gives
\[
 \abs Z\le
 \floor{\frac{n-c}{\max\{1,a-c\}}}\le r+1.        \tag{MO2}
\label{eq:large-overlap-count}
\]
Indeed, for \(c\ge a\) the numerator is at most \(r\); for \(c=a-1\)
it is \(r+1\); and, with \(x=a-c\ge2\), the quotient is
\((r+x)/x\le r+1\).
\end{proof}

\begin{theorem}[Quadratic mean-overlap theorem]
\label{thm:quadratic-mean-overlap}
Let \(C=\RS_\F(D,k)\), put \(R=n-k\), let
\(0\le r\le R-1\), and let \(\varnothing\ne\Gamma\subseteq\F\).  If
\[
 (n-r)^2\ge n(k+r),
 \qquad\text{equivalently}\qquad
 r^2\ge n(3r-R),                                  \tag{MO3}
\label{eq:quadratic-overlap-condition}
\]
then
\[
 B_{C,\Gamma}^{\rm MCA}(n-r)=\min\{\abs\Gamma,r+1\}. \tag{MO4}
\label{eq:quadratic-overlap-exact}
\]
The same upper bound holds for every linear \([n,k]\) MDS code.
\end{theorem}

\begin{proof}
The universal tangent construction gives the lower bound.  For the upper
bound, fix a received pair and its challenge-restricted bad set \(Z\).
If \(\abs Z\le r+1\) there is nothing to prove.  Otherwise choose
\(r+2\) slopes and distinct exact witness supports
\(S_1,\ldots,S_{r+2}\), each of size \(a=n-r\).  Suppose every pair had
intersection at most \(k+r-1\), and put
\(d_x=\#\{i:x\in S_i\}\).  Then
\[
 \sum_x\binom{d_x}{2}\le\binom{r+2}{2}(k+r-1),    \tag{MO5}
\]
whereas Cauchy--Schwarz and \(a^2/n\ge k+r\) give
\[
 \sum_x\binom{d_x}{2}
 \ge\frac{(r+2)^2a^2/n-(r+2)a}{2}
 \ge\frac{(r+2)^2(k+r)-(r+2)a}{2}.                \tag{MO6}
\]
The last lower bound exceeds the upper bound in \textup{(MO5)} by
\[
 \frac{r+2}{2}\bigl(k+3r-n+1\bigr).
\]
When \(3r>R\) this is positive.  When \(3r\le R\), the exact deep theorem
already gives the desired upper bound.  Thus in the remaining case some
pair satisfies \eqref{eq:large-overlap-collapse}, and
\cref{lem:large-overlap-collapse} contradicts \(\abs Z\ge r+2\).
The exact-support reduction and large-overlap lemma use only the MDS
information-set property, proving the stated extension.
\end{proof}

\begin{corollary}[Exact quadratic staircase]
\label{cor:mean-overlap-exact}
Put
\[
 r_{\rm quad}(n,k)=
 \floor{\frac{3n-\sqrt{n(5n+4k)}}2}.
 \tag{MO7}\label{eq:quadratic-radius}
\]
For every \(0\le r\le r_{\rm quad}(n,k)\) and every nonempty challenge
set \(\Gamma\),
\[
 B_{C,\Gamma}^{\rm MCA}(n-r)=\min\{\abs\Gamma,r+1\}.
\]
If \(k/n\to\rho\), then
\(r_{\rm quad}(n,k)/n\to
\alpha_\rho=(3-\sqrt{5+4\rho})/2\).
\end{corollary}

\begin{proof}
The smaller root of \(r^2-3nr+n(n-k)=0\) is the real number inside the
floor in \eqref{eq:quadratic-radius}.  Apply
\cref{thm:quadratic-mean-overlap}; the limiting formula follows by
division by \(n\).
\end{proof}

\begin{proposition}[Limit of pairwise-overlap information]
\label{prop:pairwise-overlap-limit}
Let \(k/n\to\rho\in(0,1)\), \(r/n\to\alpha\in(0,1)\), and suppose
\[
 \alpha^2<3\alpha-(1-\rho).
 \tag{MO8}\label{eq:pairwise-overlap-limit}
\]
For all sufficiently large \(n\), there are \(r+2\) distinct
\(r\)-subsets \(T_i\) such that, whenever
\(j=3r-(n-k)\le r\),
\(\abs{T_i\cap T_\ell}\le j-1\) for every \(i\ne\ell\).  If \(j>r\),
that intersection condition is automatic.  Hence pairwise support overlap
alone cannot extend the exact theorem past the quadratic boundary; further
progress must use the syndrome, determinant, or fiber incidence.  This is
a limitation of the method, not a construction of an MCA-bad line.
\end{proposition}

\begin{proof}
When \(j\le r\), choose \(r+2\) independent uniform \(r\)-subsets.
The intersection of a fixed pair is hypergeometric with mean
\(r^2/n=(\alpha^2+o(1))n\), whereas
\(j=(3\alpha-(1-\rho)+o(1))n\).  More explicitly,
\[
 \Pr(\abs{T_1\cap T_2}=s)
 =\frac{\binom rs\binom{n-r}{r-s}}{\binom nr}.
\]
After writing \(s=xn\), Stirling's formula gives the exponent
\[
 \alpha h(x/\alpha)+(1-\alpha)
 h((\alpha-x)/(1-\alpha))-h(\alpha),
\]
which is strictly concave and has its unique maximum zero at
\(x=\alpha^2\).  The strict gap in
\eqref{eq:pairwise-overlap-limit}, followed by summing at most \(n+1\)
terms, therefore gives probability \(e^{-\Omega(n)}\) that one fixed
intersection reaches \(j\).  A union bound over \(O(n^2)\) pairs is
smaller than one.  A
successful family is automatically distinct, because equal members have
intersection \(r\ge j\).  When \(j>r\), choose any \(r+2\) distinct
\(r\)-sets.
\end{proof}

The following complementary theorem can be stronger when the overlap
defect \(j=3r-R\) is small.  Its only new input is a packing inequality;
in particular it uses no smoothness or Fourier estimate.

\begin{theorem}[Exact overlap--packing criterion]
\label{thm:overlap-packing-exact}
Let \(C=\RS_\F(D,k)\), put \(R=n-k\), and fix an integer
\(0\le r\le R-1\).  Set \(a=n-r\) and \(j=3r-R\).  If either
\[
 j\le0,
 \tag{OP0}\label{eq:op-deep}
\]
or
\[
 1\le j\le r,
 \qquad
 \binom nj\le(r+1)\binom rj,
 \tag{OP}\label{eq:op-packing}
\]
then
\[
                    B_C^{\rm MCA}(n-r)=r+1.       \tag{OP$'$}
\label{eq:op-exact}
\]
The same upper bound holds for every linear \([n,k]\) MDS code; the
matching lower bound holds whenever the field contains at least \(r+1\)
distinct slope parameters.
\end{theorem}

\begin{proof}
Condition \eqref{eq:op-deep} is
\cref{cor:exact-deep-numerator}, so assume
\eqref{eq:op-packing}.  Fix an arbitrary received pair \((f,g)\), let
\(Z\) be its bad-slope set, and use
\cref{lem:exact-agreement-reduction} to choose, for every \(z\in Z\), an
exact witness support \(S_z\), \(\abs{S_z}=a\), and an explaining
polynomial \(p_z\in\F[X]_{<k}\).  A single noncommon support cannot carry
two distinct slopes, so the supports chosen for distinct elements of
\(Z\) are distinct.

For the stated MDS extension, the same exact-support reduction is
available: interpolate the unexplained component on any \(k\)-subset of
the original witness support; because every \(k\)-subset is an information
set, nonexplanation is detected at one further coordinate, and that
\((k+1)\)-set can be enlarged to size \(a\).

First suppose that two of them have a large intersection:
\[
 \abs{S_z\cap S_w}\ge n+k-a=k+r
 \qquad(z\ne w).
 \tag{OP1}\label{eq:op-large-overlap}
\]
Then \cref{lem:large-overlap-collapse} gives \(\abs Z\le r+1\).

It remains to consider the complementary case, in which every pair of
chosen supports satisfies
\(\abs{S_z\cap S_w}\le n+k-a-1\).  Put
\(T_z=D\setminus S_z\), so \(\abs{T_z}=r\).  For distinct slopes,
\[
 \abs{T_z\cap T_w}
 =n-2a+\abs{S_z\cap S_w}
 \le3r-R-1=j-1.                                  \tag{OP3}
\label{eq:op-small-overlap}
\]
No \(j\)-subset of \(D\) can therefore lie in two different \(T_z\)'s.
Counting pairs \((z,J)\) with \(J\in\binom{T_z}{j}\) yields
\[
                 \abs Z\binom rj\le\binom nj\le
                 (r+1)\binom rj,
\]
and hence \(\abs Z\le r+1\).  The universal tangent construction gives
the reverse inequality.  The proof used only linearity, the fact that
agreement of two codewords on \(k\) coordinates forces equality, and the
exact-support reduction; these are MDS properties.  The coordinate
tangent construction gives the stated MDS lower bound.
\end{proof}

\begin{proposition}[General overlap--packing upper bound]
\label{prop:general-overlap-packing-upper}
For \(1\le j=3r-(n-k)\le r\), let \(A(n,r,j)\) be the largest size of a
family of \(r\)-subsets of an \(n\)-set whose distinct members intersect
in at most \(j-1\) points.  Then
\[
 B_{C,\Gamma}^{\rm MCA}(n-r)
 \le\min\left\{\abs\Gamma,
     \max\left(r+1,A(n,r,j)\right)\right\}          \tag{OP5}
\label{eq:general-overlap-packing-upper}
\]
and
\[
 A(n,r,j)\le
 \floor{\frac{\binom nj}{\binom rj}}.              \tag{OP6}
\label{eq:packing-number-bound}
\]
\end{proposition}

\begin{proof}
For a fixed received pair, either two chosen exact witness supports have
intersection at least \(k+r\), in which case
\cref{lem:large-overlap-collapse} gives at most \(r+1\) slopes, or every
pair of complementary \(r\)-sets intersects in at most \(j-1\) points.
The latter family has size at most \(A(n,r,j)\).  Finally, no \(j\)-subset
can lie in two members of such a packing, so counting contained
\(j\)-subsets proves \eqref{eq:packing-number-bound}.
\end{proof}

\begin{corollary}[Quadratic staircase at the official rates]
\label{thm:prize-rate-overlap}
Let \(n=2^e\), let
\(\rho=k/n\in\{1/2,1/4,1/8,1/16\}\), and put
\[
 r_\rho(n)=\floor{
 n\frac{3-\sqrt{5+4\rho}}2}.
\]
For every \(n\)-point Reed--Solomon domain over every field,
\[
 B_C^{\rm MCA}(n-r)=r+1
 \qquad(0\le r\le r_\rho(n)).                     \tag{OP4}
\label{eq:prize-rate-staircase}
\]
\end{corollary}

\begin{proof}
This is \cref{cor:mean-overlap-exact}, since the displayed integer is
\(r_{\rm quad}(n,k)\).
\end{proof}

\begin{corollary}[Self-contained quadratic-range target compiler]
\label{cor:prize-window-compiler}
In the setting of \cref{thm:prize-rate-overlap}, put
\(q=\abs\F\) and \(B=\floor{q/2^{128}}\).  If
\[
 1\le B\le M_\rho(n)
 \defeq\min\{r_\rho(n)+1,n-k-1\},
\]
then
\[
 e_{\rm MCA}(C,\delta)\le2^{-128}
 \quad\Longleftrightarrow\quad
 \floor{\delta n}\le B-1.                        \tag{PW}
\label{eq:prize-window-iff}
\]
Thus every such row is solved throughout the field window
\(2^{128}\le q<2^{128}(M_\rho(n)+1)\): its largest safe grid radius is
\((B-1)/n\), while its real transition supremum is \(B/n\) and that
endpoint is unsafe.  If \(q<2^{128}\), even radius zero is unsafe.
\end{corollary}

\begin{proof}
For \(r\le B-1\), \cref{thm:prize-rate-overlap} gives the exact
numerator \(r+1\le B\le q/2^{128}\).  At \(r=B\), the tangent floor
applies because \(B\le n-k-1\), and gives at least
\(B+1>q/2^{128}\) slopes.  Monotonicity proves the
equivalence and the endpoint statements.  If \(B=0\), the tangent floor
at \(r=0\) gives the exact numerator one.
\end{proof}


\section{Literature completion to positive relative radius}
\label{sec:literature-completion}

The quadratic theorem is self-contained and useful at finite parameters,
but it is not the largest known RS range at constant rate.  We now make
the relevant implication of \cite{BCHKS25} explicit at the integer
numerator level.  This section is deliberately separated from
\cref{sec:mean-overlap}: the input below belongs to the literature,
whereas \cref{thm:quadratic-mean-overlap} is the finite MDS argument of
this paper.

Put \(q=\abs\F\).  For a word \(u\) and a received pair \((u_0,u_1)\), write
\(d(u,C)\) and \(d_2((u_0,u_1),C^2)\) for their relative Hamming and
relative column distances from \(C\) and \(C^2\), respectively.  For
\(0\le\alpha<\beta<1\), define the two-radius correlated-agreement error
by
\[
 e_{\rm CA}(C,\alpha,\beta)
 =\frac1q\max_{\substack{u_0,u_1\in\F^D\\
                  d_2((u_0,u_1),C^2)>\beta}}
   \#\{\gamma\in\F:d(u_0+\gamma u_1,C)\le\alpha\}.       \tag{LC0}
 \label{eq:two-radius-ca}
\]
The one-radius notation \(e_{\rm CA}(C,\alpha)\) means that the far-pair
radius is also \(\alpha\).

\begin{theorem}[BCHKS two-radius bound]
\label{thm:bchks-two-radius}
Let \(C=\RS_\F(D,k)\), put \(q=\abs\F\),
\(\rho=k/n\), and \(\Delta=(n-k+1)/n\).  If
\(\Delta/3\le\alpha<(1-\rho)/2\) and
\(\alpha<\beta<1\), then
\[
 e_{\rm CA}(C,\alpha,\beta)
 \le \frac1q\max\left\{
   \frac{1-\rho-\alpha}{\alpha(1-\rho-2\alpha)},
   \frac{\beta}{\beta-\alpha}
                         \right\}.              \tag{BCHKS}
 \label{eq:bchks-two-radius}
\]
\end{theorem}

This is \cite[Thm.~1.3]{BCHKS25}, in the two-radius normalization
restated as \cite[Thm.~4.9\textup{(2)}]{ABF26}.  We have written only
the strict subrange needed below, so all displayed denominators are
positive.

\begin{theorem}[Integral BCHKS completion]
\label{thm:literature-integral-completion}
Let \(C=\RS_\F(D,k)\), let \(q=\abs\F\), put \(R=n-k\), and let
\(r\in\Z\) satisfy
\[
 3r\ge R+1,\qquad 2r<R.                           \tag{LC1}
 \label{eq:literature-range}
\]
Then
\[
 B_C^{\rm MCA}(n-r)
 \le
 \floor{\max\left\{
       \frac{n(R-r)}{r(R-2r)},\ r+1
                  \right\}}.                    \tag{LC2}
 \label{eq:literature-integral-upper}
\]
Consequently, if
\[
             n(R-r)<r(r+2)(R-2r),                \tag{LC3}
 \label{eq:literature-integral-condition}
\]
then, for every nonempty \(\Gamma\subseteq\F\),
\[
 B_{C,\Gamma}^{\rm MCA}(n-r)=\min\{\abs\Gamma,r+1\}.
 \tag{LC4}\label{eq:literature-integral-exact}
\]
\end{theorem}

\begin{proof}
Put \(\delta_{\rm fld}=r/n\) and, for \(0<\theta<1\), put
\(\delta_{\rm int}=(r+\theta)/n\).  Hamming distances belong to
\(\{0,1/n,\ldots,1\}\), so replacing the far-pair radius
\(\delta_{\rm fld}\) by \(\delta_{\rm int}\) does not change the
far-pair condition; this is the grid observation of
\cite[Rem.~4.10]{ABF26}.  The hypotheses in
\eqref{eq:literature-range} give
\(\Delta/3\le\delta_{\rm fld}<(1-k/n)/2\), and
\(\delta_{\rm fld}<\delta_{\rm int}<1\).  Hence
\cref{thm:bchks-two-radius} gives
\[
 q\,e_{\rm CA}(C,r/n)
 \le \max\left\{
  \frac{n(R-r)}{r(R-2r)},\frac{r+\theta}{\theta}
               \right\}.                        \tag{LC5}
 \label{eq:bchks-specialized}
\]
On the integer grid,
\(q e_{\rm CA}(C,r/n)=B_C^{\rm CA}(n-r)\).  The second inequality in
\eqref{eq:literature-range} also places the radius below
\(d_{\min}/(2n)\), because \(d_{\min}=R+1\).  By
\cref{thm:exact-half-distance-sparse},
\[
 B_C^{\rm MCA}(n-r)
   =\max\{B_C^{\rm CA}(n-r),\min\{q,r\}\}.
\]
The right side of \eqref{eq:bchks-specialized} is strictly larger than
\(r\), because \((r+\theta)/\theta>r\).  It therefore bounds both
branches of this maximum and hence bounds the integer
\(B_C^{\rm MCA}(n-r)\).
For the full-field maximum one may alternatively invoke the known
below-half equality \cite[Lem.~4.6]{ABF26}; the sparse decomposition is
kept here because it is exact before maximizing over a restricted
challenge set.

Letting \(\theta\to1^-\) proves that this integer is at most the real
maximum in \eqref{eq:literature-integral-upper}, and taking the floor
proves \eqref{eq:literature-integral-upper}.  Under
\eqref{eq:literature-integral-condition} the first entry of that maximum
is strictly less than \(r+2\).  Since \(r+1\) is integral, the floor in
\eqref{eq:literature-integral-upper} is at most \(r+1\), so the
full-field numerator is at most \(r+1\).
Restriction to \(\Gamma\) gives the upper bound
\(\min\{\abs\Gamma,r+1\}\), and
\cref{prop:universal-tangent-floor} gives the reverse bound.
\end{proof}

\begin{remark}[Priority]
The substance of the upper bound in
\cref{thm:literature-integral-completion} is prior work.  In particular,
\cite[Cor.~1.4]{BCHKS25} states a sharp lossless conclusion under broad
finite hypotheses and \cite[Rem.~2.5]{BCHKS25} discusses the matching
perturbation construction.  Formula \eqref{eq:bchks-specialized} is
included to expose the integer limit and to make the challenge-subset and
target compilers below completely transparent.
\end{remark}

\begin{corollary}[Positive density of exact radii]
\label{cor:positive-radius-density}
Let \(k_n/n\to\rho\in(0,1)\), let \(r_n\in\Z\) satisfy
\(r_n/n\to\tau\), and suppose
\(0<\tau<(1-\rho)/2\).  For every RS row sequence and every sequence of
nonempty challenge sets,
\[
 B_{C_n,\Gamma_n}^{\rm MCA}(n-r_n)
 =\min\{\abs{\Gamma_n},r_n+1\}                   \tag{LC6}
 \label{eq:positive-radius-exact}
\]
for all sufficiently large \(n\).  More uniformly, for every
\(\eta>0\), equality holds throughout
\[
 0\le r\le\left(\frac{1-\rho}{2}-\eta\right)n
\]
for every sufficiently large row whose rate tends to \(\rho\).
Consequently the exact staircase contains an initial interval of
normalized length \((1-\rho)/2-o(1)\), has lower density at least
\((1-\rho)/2\) in the full Hamming-radius interval, and has asymptotic
relative density one in the unique-decoding interval.
\end{corollary}

\begin{proof}
The quadratic theorem covers every sequence whose limiting radius lies
below \(\alpha_\rho=(3-\sqrt{5+4\rho})/2\).  In the remaining compact
subinterval of \((0,(1-\rho)/2)\), the hypotheses
\eqref{eq:literature-range} hold for large \(n\), while
\[
 \frac{n(R-r)}{r(R-2r)}=O_{\rho,\eta}(1)
 \quad\text{and}\quad r+1=\Theta_{\rho,\eta}(n).
\]
Thus \eqref{eq:literature-integral-condition} holds and
\cref{thm:literature-integral-completion} applies.  The uniform statement
follows by the same compactness argument, using the deep theorem near
radius zero.  Counting integer radii gives the density assertions.
\end{proof}

\begin{corollary}[Official-rate staircase]
\label{cor:official-rate-half-staircase}
Let \(n=2^e\ge128\), \(k=\rho n\), and
\(\rho\in\{1/2,1/4,1/8,1/16\}\).  Define
\[
 r_\rho^\sharp=\begin{cases}
 n/4-3,&\rho=1/2,\\
 3n/8-2,&\rho=1/4,\\
 7n/16-2,&\rho=1/8,\\
 15n/32-2,&\rho=1/16.
 \end{cases}                                      \tag{LC7}
 \label{eq:official-half-endpoints}
\]
Then for every RS domain and every nonempty \(\Gamma\subseteq\F\),
\[
 B_{C,\Gamma}^{\rm MCA}(n-r)=\min\{\abs\Gamma,r+1\}
 \qquad(0\le r\le r_\rho^\sharp).              \tag{LC8}
 \label{eq:official-half-staircase}
\]
\end{corollary}

\begin{proof}
Put \(F(r)=r(r+1)(R-2r)-n(R-r)\), where \(R=n-k\), and
\(L=\ceil{(R+1)/3}\).  The deep theorem covers \(r<L\).  On the
positive real axis, \(F'(r)=-6r^2+2(R-2)r+(R+n)\) has exactly one
positive zero, so the minimum of \(F\) on \([L,r_\rho^\sharp]\) occurs
at an endpoint.  Since \(R\ge n/2\ge64\),
\[
 F(L)\ge
 \frac{(R+1)(R+4)(R-8)-9n(2R-1)}{27}>0;
\]
indeed the numerator is at least
\(R\{R(R-4)-18n\}>0\).  Direct substitution at the other endpoint gives,
in the order of the four rates,
\[
 \frac{n^2-84n+288}{8},\quad
 \frac{3n^2-104n+128}{16},\quad
 \frac{21n^2-464n+512}{64},\quad
 \frac{105n^2-1952n+2048}{256},
\]
all positive for \(n\ge128\).  Therefore \(F(r)\ge0\) throughout the
remaining integer interval.  Since \(r>0\) and \(R-2r>0\) there,
\[
 r(r+2)(R-2r)-n(R-r)=F(r)+r(R-2r)>0,
\]
which is \eqref{eq:literature-integral-condition}.  Apply
\cref{thm:literature-integral-completion}.
\end{proof}

\begin{corollary}[Almost-half-distance target compiler]
\label{cor:half-distance-window-compiler}
In the setting of \cref{cor:official-rate-half-staircase}, put
\(B=\floor{\eps\abs\Gamma}\).  If
\[
 1\le B<\abs\Gamma,\qquad B\le r_\rho^\sharp+1,
\]
then the complete safe set is
\[
 \{\delta:e_{\rm MCA,\Gamma}(C,\delta)\le\eps\}
       =[0,B/n),                                   \tag{LC9}
 \label{eq:half-distance-safe-set}
\]
and the endpoint is unsafe.  For the full-field target \(2^{-128}\),
this applies throughout
\(2^{128}\le q<2^{128}(r_\rho^\sharp+2)\).
\end{corollary}

\begin{proof}
At every radius \(r\le B-1\),
\cref{cor:official-rate-half-staircase} gives numerator
\(r+1\le B\).  At radius \(B\), the tangent floor gives at least
\(B+1\) bad slopes.  Monotonicity and
\(r(\delta)=\floor{\delta n}\) prove the assertion.
\end{proof}


\section{Exact thresholds for explicit prize rows}\label{sec:prize-rows}

\begin{table}[H]
\centering
\caption{Explicit maximal-dimension prize-admissible rows.  In each row
\(k=2^{40}\), \(B=\floor{p/2^{128}}=r_\rho(n)+1\), and \((s,a_0)\) is part of
the Proth certificate in \cref{prop:proth-row-check}.}
\label{tab:prize-proth-rows}
\small
\begin{tabular}{@{}ccrrc@{}}
\hline
rate & length & \(B\) & bits of \(p\) & \((s,a_0)\)\\
\hline
\(1/2\)  & \(2^{41}\) & \(389500552609\)  & 167 & \((92,3)\)\\
\(1/4\)  & \(2^{42}\) & \(1210584858040\) & 169 & \((93,13)\)\\
\(1/8\)  & \(2^{43}\) & \(2879806199253\) & 170 & \((95,5)\)\\
\(1/16\) & \(2^{44}\) & \(6233898019554\) & 171 & \((97,5)\)\\
\hline
\end{tabular}
\end{table}

For completeness, the four field orders are
\begin{align*}
p_{1/2}&=132540169958804033333249306710494641010898987122689,\\
p_{1/4}&=411940680852499481698306614369841346700408394874881,\\
p_{1/8}&=979947269755402568812854322316630667196565607677953,\\
p_{1/16}&=2121285573237585848299875619011192262679065433997313.
\end{align*}
Their odd Proth coefficients \(u\), in the same order, are
\begin{align*}
u_{1/2}&=26766274163673319604503,
&u_{1/4}&=41595378994516821279015,\\
u_{1/8}&=24737346889219389259839,
&u_{1/16}&=13387194060291799253121.
\end{align*}
For \(F_{n,k}(r)=r^2-n(3r-(n-k))\), the exact boundary checks are
\[
\begin{array}{c|r|r}
 \rho&F_{n,k}(B-1)&F_{n,k}(B)\\ \hline
 1/2&5154112775168&-663955886271\\
 1/4&7590647904465&-3182321912768\\
 1/8&13908181940112&-6720484728007\\
 1/16&19335616403905&-20973145690236.
\end{array}                                                    \tag{PC0}
\label{eq:quadratic-row-certificates}
\]

\begin{proposition}[Exact arithmetic and primality certificates]
\label{prop:proth-row-check}
For the four rows in \cref{tab:prize-proth-rows}, the displayed integer
\(p\) satisfies
\[
 p=u\,2^s+1,\qquad u\text{ odd},\qquad u<2^s,\qquad
 a_0^{(p-1)/2}\equiv-1\pmod p.                  \tag{PC1}
\label{eq:proth-certificates}
\]
Consequently every \(p\) is prime.  Moreover
\[
 n\mid p-1,\qquad p<2^{256},\qquad
 B\,2^{128}\le p<(B+1)2^{128},                  \tag{PC2}
\label{eq:prize-row-arithmetic}
\]
where \(n\) and \(B\) are the corresponding entries of the table.
They also satisfy \(F_{n,k}(B-1)\ge0>F_{n,k}(B)\).  Since
\(F_{n,k}'(r)=2r-3n<0\) on \([0,n]\), these two signs locate its smaller
root and give \(B-1=r_\rho(n)=r_{\rm quad}(n,k)\).
Thus \(\F_p^*\) contains a subgroup of order \(n\), and each code in
\cref{thm:intro-prize-rows} exists with the asserted slope budget.
\end{proposition}

\begin{proof}
All identities and inequalities in
\eqref{eq:quadratic-row-certificates}--\eqref{eq:prize-row-arithmetic}
are direct
integer calculations; the values of \((s,a_0)\), \(u\), \(p\), \(n\),
and \(B\) have all been printed above so that the calculation is
independently reproducible.  We recall why the certificate proves
primality.  If a prime \(\ell\) divides \(p\), then
\(a_0^{u2^{s-1}}\equiv-1\pmod\ell\), whereas
\(a_0^{u2^s}\equiv1\pmod\ell\).  The multiplicative order of \(a_0\)
modulo \(\ell\) therefore has exact \(2\)-part \(2^s\), so
\(2^s\mid\ell-1\).  Since \(u<2^s\), one has \(2^s>\sqrt p\); hence a
composite \(p\) would have a prime divisor both at most \(\sqrt p\) and
at least \(2^s+1\), a contradiction.  This is Proth's theorem in the
special form used here.  Finally \(s\ge\log_2 n\) gives \(n\mid p-1\),
and cyclicity of \(\F_p^*\) gives the required subgroup.
\end{proof}

\begin{theorem}[Exact first adjacent row over a separating field]
\label{thm:exact-first-adjacent-row}
Let \(C=\RS_\F(D,k)\), put \(R=n-k\) and
\(M=\binom n{R-1}=\binom n{k+1}\), and define
\[
                    Q_{\rm sep}(M)=\max\left\{M,\binom M2\right\}.
\]
If \(\abs\F>Q_{\rm sep}(M)\), then for the full-field challenge
\[
                         B_C^{\rm MCA}(k+1)=M.      \tag{AD1}
\]
If \(R\ge2\) and \(b=\floor{\eps\abs\F}\), the exact target threshold
therefore satisfies
\[
 \begin{array}{ll}
 b\ge M &\Longrightarrow\quad a^*=k+1,\\[2mm]
 \displaystyle\binom n{R-2}\le b<M
   &\Longrightarrow\quad a^*=k+2.
 \end{array}                                       \tag{AD2}
\]
Thus the first finite adjacent comparison has literal binomial constants;
no \(e^{o(n)}\) or unspecified polynomial factor occurs.
\end{theorem}

\begin{proof}
The upper bound in \textup{(AD1)} is
\cref{prop:exact-support-upper}.  For the reverse inequality use the
Reed--Solomon parity check of redundancy \(R\).  For every
\((R-1)\)-set \(E\subseteq D\), its column span \(V_E\) is a hyperplane
in \(\F^R\), and distinct sets give distinct hyperplanes: otherwise the
union of two distinct \((R-1)\)-sets would contain at least \(R\) columns
in one \((R-1)\)-space, contradicting the MDS column property.

Every proper linear hyperplane in \(\F^R\) has
\(\abs{\F}^{R-1}\) points.  Since \(M<\abs\F\), the union of the \(M\)
spaces \(V_E\) has fewer than \(\abs{\F}^{R}\) points; choose \(y_1\)
outside it.  For each \(E\),
choose a nonzero linear functional \(\ell_E\) with kernel \(V_E\).
The line \(y_0+\gamma y_1\) meets \(V_E\) at
\[
                 \gamma_E=-\ell_E(y_0)/\ell_E(y_1).
\]
For \(E\ne E'\), equality \(\gamma_E=\gamma_{E'}\) cuts out a proper
hyperplane in the variable \(y_0\), because the normalized functionals
\(\ell_E/\ell_E(y_1)\) and
\(\ell_{E'}/\ell_{E'}(y_1)\) are distinct.  There are at most
\(\binom M2<\abs\F\) collision hyperplanes; their union also has fewer
than \(\abs{\F}^{R}\) points, so choose \(y_0\) outside it.  Since
\(y_1\notin V_E\), the line is not contained in \(V_E\), and every
intersection is transverse.  The \(M\) parameters \(\gamma_E\) are then
finite and pairwise distinct.  Syndrome surjectivity and
\cref{thm:syndrome-secant-exact} realize them as \(M\) bad slopes at
agreement \(k+1\), proving \textup{(AD1)}.

If \(b\ge M\), the support upper bound makes the first allowed agreement
safe.  If \(\binom n{R-2}\le b<M\), equality \textup{(AD1)} makes
\(k+1\) unsafe, whereas
\cref{prop:exact-support-upper} makes \(k+2\) safe.  Monotonicity proves
\textup{(AD2)}.
\end{proof}

\begin{theorem}[Subfield confinement]\label{thm:subfield-confinement-full}
Let \(\B\subseteq\F\), \(D\subseteq\B\), and let \(u_0,u_1\in\B^D\).  If
\(\gamma\in\F\setminus\B\) and \(u_0+\gamma u_1\) is explained on \(S\),
then the pair is explained on the same \(S\).  Hence every MCA-bad slope of
a \(\B\)-valued received line lies in \(\B\).
\end{theorem}

\begin{proof}
Choose a \(\B\)-basis \(1,\beta_1,\ldots\) of \(\F\).  Decompose the
coefficients of the explaining polynomial and of \(\gamma\) in this basis.
Because \(D\subseteq\B\), comparison of basis components on \(S\) gives a
codeword explaining \(u_1\) from any nonzero nonbase component of
\(\gamma\), and then a codeword explaining \(u_0\) from the base component.
\end{proof}

\section{Exact locator prefixes and collision-aware pole lines}
\label{sec:exact-prefix-pole}

Assume \(D\subseteq\B\subseteq\F\).  For an \(m\)-set \(S\subseteq D\),
write
\[
 Q_S(X)=\prod_{x\in S}(X-x)
       =X^m+c_1(S)X^{m-1}+\cdots+c_m(S),
 \qquad
 \Phi_w(S)=(c_1(S),\ldots,c_w(S)).
\]

The monic polynomial \(Q_S\), whose simple roots are precisely the points
of \(S\), is the \emph{support locator}.  Its \emph{depth-\(w\) locator
prefix} is the vector of the first \(w\) coefficients below the leading
coefficient.  For \(z\in\B^w\), the inverse image
\(\Phi_w^{-1}(z)\) is the \emph{prefix fiber over \(z\)}: the family of
all \(m\)-element supports having that same prefix.  Thus ``fiber'' in this
paper always refers to the inverse image of a displayed map.

\begin{proposition}[Exact prefix-list bijection]\label{prop:exact-prefix-list}
Let \(1\le K\le m\le n\), put \(w=m-K\), and set
\(U_z(X)=X^m+\sum_{i=1}^wz_iX^{m-i}\) for \(z\in\B^w\).  The codewords of
\(\RS_\F(D,K)\) agreeing with \(U_z|_D\) on at least \(m\) points are
exactly
\[
        (U_z-Q_S)|_D,\qquad S\in\Phi_w^{-1}(z).
 \tag{4.1}\label{eq:prefix-bijection}
\]
No such codeword agrees on more than \(m\) points.  In particular, some
received word has list size at least
\(\lceil\binom nm\abs\B^{-w}\rceil\).
\end{proposition}

\begin{proof}
For \(S\in\Phi_w^{-1}(z)\), the leading term and next \(w\) coefficients
cancel, so \(\deg(U_z-Q_S)<m-w=K\).  Conversely, if \(P\), \(\deg P<K\),
agrees with \(U_z\) at \(m\) points, then \(U_z-P\) is monic of degree
\(m\) with those \(m\) roots and equals their locator.  Coefficient
comparison gives the prefix.  A nonzero degree-\(m\) polynomial has no
additional root.  Pigeonhole the \(\binom nm\) supports over
\(\abs\B^w\) prefixes.
\end{proof}

The \emph{simple-pole conversion} chooses an auxiliary point
\(\alpha\notin D\) and divides coordinatewise by \(X-\alpha\).  This
turns one list of nearby polynomials into points on a received affine line
for the dimension-\(k\) code.  A \emph{collision} is an equality
\(P_i(\alpha)=P_j(\alpha)\) for two distinct listed polynomials; the next
formula pays such equalities rather than assuming all slopes are distinct.

\begin{theorem}[Collision-aware simple-pole conversion]
\label{thm:collision-aware-pole}
Let \(C=\RS_\F(D,k)\), \(C^+=\RS_\F(D,k+1)\), \(1\le k<n\), and
\(q=\abs\F>n\).  If a word has \(L\) distinct codewords of \(C^+\)
agreeing on at least \(m\ge k+1\) points, then one received line for \(C\)
has at least
\[
 M(L)=\left\lceil
 \frac{L(q-n)}{q-n+k(L-1)}
 \right\rceil
 \tag{4.2}\label{eq:collision-aware-pole}
\]
distinct MCA-bad slopes at agreement \(m\).
\end{theorem}

\begin{proof}
Let the received word be \(U\in\F^D\), and let the listed polynomials be
\(P_1,\ldots,P_L\), each of degree at most \(k\).  For
\(\alpha\in\F\setminus D\), define the word coordinatewise by
\(f_\alpha(x)=U(x)/(x-\alpha)\) and set
\(g_\alpha(x)=-1/(x-\alpha)\).  On an agreement set for \(P_i\), the line
point of slope \(P_i(\alpha)\) equals
\((P_i(x)-P_i(\alpha))/(x-\alpha)\), a degree-less-than-\(k\) evaluation.
The pair is not simultaneously explained on \(k+1\) points: an explanation
of \(g_\alpha\) by \(G\), \(\deg G<k\), would make
\((X-\alpha)G(X)+1\) vanish at \(k+1\) points although its degree is at
most \(k\) and its value at \(\alpha\) is \(1\).

For \(i\ne j\), the equality \(P_i(\alpha)=P_j(\alpha)\) holds at at most
\(k\) poles.  Some pole has at most \(k\binom L2/(q-n)\) colliding pairs.
If its value multiplicities are \(m_1,\ldots,m_M\), then
\(\sum m_i^2\le L+kL(L-1)/(q-n)\).  Cauchy--Schwarz gives
\(L^2\le M\sum m_i^2\), which yields \eqref{eq:collision-aware-pole}.
\end{proof}

\begin{corollary}[Exact good-pole reservoir]
\label{cor:exact-good-pole-reservoir}
In \cref{thm:collision-aware-pole}, let \(Q=q-n\), and for
\(\alpha\in\F\setminus D\) let \(C_\alpha\) be the number of unordered
pairs \(\{i,j\}\) with \(P_i(\alpha)=P_j(\alpha)\).  For every integer
\(T\ge0\), at least
\[
 \max\left\{0,Q-
   \floor{\frac{k\binom L2}{T+1}}\right\}          \tag{PR1}
\label{eq:good-pole-count}
\]
poles satisfy \(C_\alpha\le T\), and every such pole produces at least
\[
                     \ceil{\frac{L^2}{L+2T}}       \tag{PR2}
\label{eq:good-pole-slopes}
\]
distinct MCA-bad slopes.  In particular at least
\(\max\{0,Q-k\binom L2\}\) poles are fully separating.
\end{corollary}

\begin{proof}
Every nonzero difference \(P_i-P_j\) has at most \(k\) roots, so
\[
                 \sum_{\alpha\in\F\setminus D}C_\alpha
                    \le k\binom L2.
\]
Every pole with \(C_\alpha>T\) contributes at least \(T+1\) to this
sum, which gives \eqref{eq:good-pole-count}.  At a good
pole let \(m_1,\ldots,m_M\) be the multiplicities of the distinct list
values.  Then
\(\sum_im_i=L\) and
\(\sum_im_i^2=L+2C_\alpha\le L+2T\).  Cauchy--Schwarz yields
\(M\ge L^2/(L+2T)\), proving \eqref{eq:good-pole-slopes}; \(T=0\)
gives the last assertion.
\end{proof}

\begin{remark}[One pole, one received line]
The reservoir in \cref{cor:exact-good-pole-reservoir} fixes the same pole
when counting collisions, because different poles define different
received lines.  A cross-pole equation
\(P_i(\alpha)=P_j(\alpha')\) is one equation in two variables and
generically defines a curve; a Bezout degree bound does not turn it into a
finite collision set.  Such cross-line collisions are neither controlled
by the simple-pole argument nor needed by it.
\end{remark}

\begin{corollary}[Identity-prefix MCA floor]\label{cor:identity-prefix-floor}
For \(m\ge k+1\), put \(w=m-k-1\) and
\(L_m=\lceil\binom nm\abs\B^{-w}\rceil\).  If \(q>n\), then
\[
 B_C^{\rm MCA}(m)\ge M(L_m).
 \tag{4.3}\label{eq:identity-prefix-floor}
\]
\end{corollary}

For equal-size sets, the \emph{Johnson distance} is
\(d_J(S,T)=\abs{S\setminus T}=\abs{T\setminus S}\).  It counts the
number of point replacements required to turn one support into the other.

\begin{proposition}[Prefix rigidity and its packing limit]
\label{prop:prefix-rigidity-full}
Two distinct \(m\)-sets in one depth-\(w\) prefix fiber have Johnson
distance at least \(w+1\).  Consequently, with
\(t=\floor{w/2}\), every fiber has size at most
\[
 \frac{\binom nm}{\sum_{i=0}^t\binom mi\binom{n-m}{i}}.
 \tag{4.4}\label{eq:packing-fiber-cap}
\]
At \(w\asymp n/\log\abs\B\), this packing loss is only
\(e^{o(n)}\) and does not supply the random factor \(\abs\B^{-w}\).
\end{proposition}

Prefix rigidity says that sharing many leading locator coefficients forces
distinct supports to be far apart in this metric.  Here the
\emph{polynomial-field regime} means \(\abs\B=n^{O(1)}\), equivalently
\(\log\abs\B=O(\log n)\).

\begin{proof}
If \(S,T\) have Johnson distance \(j\le w\), factor the locator of their
intersection, of degree \(m-j\), from \(Q_S-Q_T\).  The residual factor is
nonzero, so the difference has degree at least \(m-j\ge m-w\), whereas the
common prefix makes its degree at most \(m-w-1\), a contradiction.  Radius
\(t\) Johnson balls about the fiber members are disjoint and each has the
denominator in \eqref{eq:packing-fiber-cap}.  For
\(w=n/\log\abs\B\), the logarithm of that ball volume is
\(O(w\log(n/w))=o(n)\) in the polynomial-field regime.
\end{proof}

\begin{theorem}[Prefix fibers embed in received lines]
\label{thm:prefix-to-line-hardness}
Let \(\mathcal G\subseteq\Phi_w^{-1}(z)\) be a fiber of \(N\) many
\(m\)-sets, with \(0\le w\le m-2\), and put \(k=m-w-1\).  Every finite
extension \(\F/\B\) satisfying
\[
        \abs\F-n>k\binom N2
        \tag{4.5}\label{eq:extension-separation}
\]
admits one received line for \(\RS_\F(D,k)\) with at least \(N\) distinct
MCA-bad slopes furnished by \(\mathcal G\).
\end{theorem}

\begin{proof}
For distinct \(S,T\in\mathcal G\), the polynomial \(Q_S-Q_T\) is nonzero
of degree at most \(k\).  Condition \eqref{eq:extension-separation} leaves a
pole \(\alpha\in\F\setminus D\) avoiding the roots of all these
differences.  Thus the values \(Q_S(\alpha)\) are pairwise distinct.  By
\cref{prop:exact-prefix-list}, the polynomials \(U_z-Q_S\) form one list for
the dimension-\(k+1\) code; the pole line makes their values, and hence the
resulting MCA slopes, pairwise distinct.
\end{proof}

\begin{theorem}[Exact list--line bijection at a separating pole]
\label{thm:exact-list-line-bijection}
Let \(D\subseteq\B\subseteq\F\), let \(1\le k<n\), let \(m\ge k+1\),
and let \(U\in\B^D\).  Write
\[
 \Lcal_m(U)=\{P\in\B[X]:\deg P\le k,
       \ \abs{\{x\in D:U(x)=P(x)\}}\ge m\}.
\]
Suppose \(\alpha\in\F\setminus D\) and the evaluation map
\(P\mapsto P(\alpha)\) is injective on \(\Lcal_m(U)\).  Define
\[
 f_\alpha(x)=\frac{U(x)}{x-\alpha},\qquad
 g_\alpha(x)=-\frac1{x-\alpha}\qquad(x\in D).
\]
Then \(P\mapsto P(\alpha)\) is a bijection from \(\Lcal_m(U)\) onto
the set of finite MCA-bad slopes of the received line
\((f_\alpha,g_\alpha)\) for \(\RS_\F(D,k)\) at agreement \(m\).  The
corresponding explaining polynomial is
\[
       h_P(X)=\frac{P(X)-P(\alpha)}{X-\alpha},
\]
and its agreement set with
\(f_\alpha+P(\alpha)g_\alpha\) is exactly the agreement set of \(P\)
with \(U\).

If \(L=\abs{\Lcal_m(U)}\), such a pole exists in every finite extension
satisfying
\[
       \abs\F>n+k\binom L2.                         \tag{4.6}
\]
\end{theorem}

\begin{proof}
For \(P\in\Lcal_m(U)\), the polynomial \(h_P\) has degree less than
\(k\), and on every point where \(U=P\) one has
\[
 f_\alpha+P(\alpha)g_\alpha
   =\frac{P-P(\alpha)}{X-\alpha}=h_P.
\]
The word \(g_\alpha\) cannot be explained by a polynomial \(G\) of degree
less than \(k\) on \(k+1\) points, since then
\((X-\alpha)G+1\) would have degree at most \(k\), at least \(k+1\)
roots, and value one at \(\alpha\).  Thus every displayed slope is
MCA-bad.

Conversely, suppose \(f_\alpha+\gamma g_\alpha\) is explained by
\(h\in\F[X]_{<k}\) on at least \(m\) points.  Then
\(P(X)=(X-\alpha)h(X)+\gamma\) has degree at most \(k\) and agrees with
\(U\) there.  Interpolation on any \(k+1\) of those points uses a
Vandermonde system over \(\B\) with right-hand side in \(\B\), so
\(P\in\B[X]\).  Hence \(P\in\Lcal_m(U)\) and
\(\gamma=P(\alpha)\).  Injectivity proves the bijection, and the same
calculation proves equality of the agreement sets.

For distinct \(P,Q\in\Lcal_m(U)\), the nonzero polynomial \(P-Q\) has
degree at most \(k\), and therefore at most \(k\) roots.  The union of
\(D\) and the root sets of all \(\binom L2\) differences has size at most
\(n+k\binom L2\); a larger extension contains an admissible pole.
\end{proof}

\begin{corollary}[Exact prefix-fiber ray realization]
\label{cor:exact-prefix-ray-realization}
Let \(0\le w\le m-2\), put \(k=m-w-1\), and let
\(\Fcal_z=\Phi_w^{-1}(z)\) be a complete depth-\(w\) fiber of
\(m\)-subsets of \(D\).  If \(\alpha\) separates the values
\(Q_S(\alpha)\), \(S\in\Fcal_z\), then the pole line associated with
\(U_z=X^m+\sum_{i=1}^wz_iX^{m-i}\) has exactly
\(\abs{\Fcal_z}\) MCA-bad slopes at agreement \(m\).  The maps
\[
 S\longmapsto U_z(\alpha)-Q_S(\alpha)
 \longmapsto
 \left(U_z(\alpha)-Q_S(\alpha),S,
 \frac{U_z-Q_S-(U_z(\alpha)-Q_S(\alpha))}{X-\alpha}\right)
\]
are bijections between the prefix fiber, the bad-slope set, and the
exact-\(m\) witness incidence.  In particular, the complete prefix fiber,
its bad-slope set, and its exact-agreement witness incidence have the same
cardinality.
\end{corollary}

\begin{proof}
The exact prefix-list bijection identifies \(\Fcal_z\) with the complete
dimension-\((k+1)\) list for \(U_z\).  Apply
\cref{thm:exact-list-line-bijection}.  Since
\(U_z-(U_z-Q_S)=Q_S\) has precisely the roots in \(S\), every agreement
is exactly \(m\) and supplies exactly that support.
\end{proof}



\section{Further exact constructions}\label{sec:further-exact}

The main threshold theorem needs only the syndrome correspondence and the
mean-overlap argument.  This section records complementary exact
constructions that clarify the support-to-slope projection and the
structured domains used in the examples.  Every statement is a finite
identity, inequality, or construction; no equidistribution or inverse
theorem is assumed.

\subsection{Determinantal ray normalization}\label{drc:sec-determinantal-ray}

The final projection from supports to slopes can be made explicit without
postulating a support-pair multiplicity.  The construction below applies to
every Reed--Solomon row, in every characteristic.  It also identifies the
exact additional boundary maps that would govern extensions beyond the
proved overlap range.  No estimate for those maps is used in the main
threshold theorem.

The folding atlas is stated for agreement supports
\(S\in\binom Da\), whereas this subsection uses their complementary error
supports \(E=D\setminus S\in\binom D{n-a}\).  Complementation is a
bijection.  On a degree-\(c\) complete fold it sends a fiber occupied in
\(j\) agreement points to one occupied in \(c-j\) error points, so
\(\nu^E_{\pi,j}=\nu^S_{\pi,c-j}\); it also complements the recorded subset
of the bounded exceptional set.  Hence the joint-signature census and its
exact support add-back transfer to the error-support slices without loss.

Let \(C=\RS_\F(D,k)\), put \(R=n-k\), and use the weighted Vandermonde
parity-check columns
\[
 h_x=\lambda_x(1,x,\ldots,x^{R-1})^{\mathsf T},
 \qquad \lambda_x\ne0,
\]
as in the syndrome--secant normal form.  For \(E\subseteq D\), write
\(V_E=\Span\{h_x:x\in E\}\).  Fix an exact agreement \(a\ge k+1\), and set
\[
             t=n-a,\qquad d=R-t=a-k.
             \tag{DR0}\label{drc:eq-depth}
\]
In particular, \(d-1=a-k-1\) is precisely the identity locator-prefix
depth.  For \(E\in\binom Dt\), write its error locator in ascending order as
\[
 Q_E(X)=\prod_{x\in E}(X-x)=\sum_{\ell=0}^{t}q_\ell(E)X^\ell,
 \qquad q_t(E)=1.
\]
For \(y=(\eta_0,\ldots,\eta_{R-1})\in\F^R\), define the
\emph{locator-recurrence boundary}
\[
 \mathcal A_E(y)=
 \left(\sum_{\ell=0}^{t}q_\ell(E)\eta_{j+\ell}\right)_{0\le j<d}
 \in\F^d.
 \tag{DR1}\label{drc:eq-recurrence-boundary}
\]
If \(f(X)=\sum_{i=0}^{R-1}f_iX^i\), we use the coefficient pairing
\(
 \langle y,f\rangle=\sum_{i=0}^{R-1}y_if_i
\).

\begin{lemma}[Locator recurrence and a secant span]
\label{drc:lem-locator-recurrence}
For every \(E\in\binom Dt\) and \(y\in\F^R\),
\[
                         y\in V_E
             \quad\Longleftrightarrow\quad
                         \mathcal A_E(y)=0.
             \tag{DR2}\label{drc:eq-recurrence-membership}
\]
\end{lemma}

\begin{proof}
Identify a polynomial of degree less than \(R\) with its coefficient
vector in \(\F^R\).  Since
\(
 \langle h_x,f\rangle=\lambda_xf(x)
\), every coefficient vector of a polynomial in
\(Q_E\F[X]_{<d}\) annihilates \(V_E\).  The MDS column property gives
\(\dim V_E=t\), while multiplication by the nonzero monic polynomial
\(Q_E\) gives
\(\dim Q_E\F[X]_{<d}=d=R-t\).  Consequently
\[
 V_E^\perp=
 \{\operatorname{coeff}(Q_EP):P\in\F[X]_{<d}\}.
\]
Testing against the basis \(1,X,\ldots,X^{d-1}\) of
\(\F[X]_{<d}\) gives exactly the \(d\) coordinates in
\eqref{drc:eq-recurrence-boundary}.
\end{proof}

For a syndrome pair \(y_0,y_1\in\F^R\), and a support family
\(\Omega\subseteq\binom Dt\), define its support-restricted finite bad-slope
set by
\[
 Z_\Omega(y_0,y_1)=
 \left\{\gamma\in\F:\begin{array}{l}
   \text{there is \(E\in\Omega\) with }
       y_0+\gamma y_1\in V_E,\\[-1mm]
   \{y_0,y_1\}\not\subseteq V_E
 \end{array}\right\}.
 \tag{DR3}\label{drc:eq-restricted-rays}
\]
For \(E\in\Omega\), abbreviate
\[
              u_E=\mathcal A_E(y_0),\qquad
              v_E=\mathcal A_E(y_1).
\]
Let \(\mathcal P_j(y_1)\) be the universal pivot locus on which \(j\) is
the first nonzero coordinate of \(v_E\), and put
\(\Omega_j=\Omega\cap\mathcal P_j(y_1)\):
\[
 \begin{aligned}
 \mathcal P_j(y_1)
   &=\{E\in\tbinom Dt:
          v_{E,i}=0\ \text{for }0\le i<j,\quad v_{E,j}\ne0\},\\
 \Omega_j&=\Omega\cap\mathcal P_j(y_1),
       &&0\le j<d.
 \end{aligned}
 \tag{DR4}\label{drc:eq-pivot-chart}
\]
On this chart define the \emph{determinantal ray boundary}
\[
 \Delta_j(E)=
 \bigl(v_{E,j}u_{E,\ell}-u_{E,j}v_{E,\ell}\bigr)_
       {\substack{0\le\ell<d\\ \ell\ne j}}
 \in\F^{d-1}.
 \tag{DR5}\label{drc:eq-pivot-boundary}
\]
When \(d=1\), the target in \eqref{drc:eq-pivot-boundary} is the trivial
group and the displayed tuple is empty.

\begin{theorem}[Exact determinantal ray compiler]
\label{drc:thm-exact-compiler}
With the notation above, an error support \(E\in\Omega\) supplies a finite
transverse slope if and only if
\[
 v_E\ne0,
 \qquad
 \rank\begin{pmatrix}u_E\\v_E\end{pmatrix}\le1.
 \tag{DR6}\label{drc:eq-rank-one-condition}
\]
On the pivot chart \(\Omega_j\), this is equivalent to
\(\Delta_j(E)=0\), and its unique slope is
\(
 \gamma_E=-u_{E,j}/v_{E,j}
\).  Consequently
\[
 Z_\Omega(y_0,y_1)=
 \bigcup_{j=0}^{d-1}
 \left\{-\frac{u_{E,j}}{v_{E,j}}:
       E\in\Omega_j,\ \Delta_j(E)=0\right\},
 \tag{DR7}\label{drc:eq-exact-ray-union}
\]
and
\[
 \abs{Z_\Omega(y_0,y_1)}
 \le
 \sum_{j=0}^{d-1}
 \abs{\Omega_j\cap\Delta_j^{-1}(0)}.
 \tag{DRC}\label{drc:eq-direct-compiler}
\]
The construction uses \(d\le n\) pivot charts and contains no unnamed
finite or asymptotic multiplicity.
\end{theorem}

\begin{proof}
By \cref{drc:lem-locator-recurrence}, \(E\) supports the line parameter
\(\gamma\) precisely when
\[
                         u_E+\gamma v_E=0.
                         \tag{DR8}\label{drc:eq-vector-slope}
\]
If \(v_E=0\), a solution of \eqref{drc:eq-vector-slope} forces \(u_E=0\),
so \(y_0,y_1\in V_E\) and the intersection is common rather than
transverse.  If \(v_E\ne0\), equation
\eqref{drc:eq-vector-slope} has a solution if and only if \(u_E\) is a
scalar multiple of \(v_E\), which is
\eqref{drc:eq-rank-one-condition}; the scalar, and therefore the finite
slope, is unique.  The charts \(\Omega_j\) partition the locus
\(v_E\ne0\).  On \(\Omega_j\), vanishing of all \(2\)-by-\(2\) minors
using the pivot column \(j\) is equivalent to proportionality, and the
proportionality factor is \(-u_{E,j}/v_{E,j}\).  This proves
\eqref{drc:eq-exact-ray-union}.  Taking the cardinality of a union and then
forgetting the slope projection gives \eqref{drc:eq-direct-compiler}.
\end{proof}

\begin{proposition}[Exact slope-incidence linearization]
\label{drc:prop-exact-slope-linearization}
With the notation of \cref{drc:thm-exact-compiler}, put
\[
 \Ical_\Omega=
 \{(\gamma,E)\in\F\times\Omega:
                 u_E+\gamma v_E=0,\ v_E\ne0\}.
 \tag{DRC1}\label{drc:eq-slope-incidence}
\]
Projection \(\Ical_\Omega\to\Omega\) is injective, its slope projection
is \(Z_\Omega(y_0,y_1)\), and
\[
 \abs{Z_\Omega(y_0,y_1)}
 \le\abs{\Ical_\Omega}
 =\sum_{\gamma\in\F}
   \abs{\{E\in\Omega:u_E+\gamma v_E=0,\ v_E\ne0\}}. \tag{DRC2}
\label{drc:eq-slope-sum}
\]
Thus a fixed-slope zero-fiber estimate must be aggregated over all slopes.
If every fixed-slope zero fiber were bounded by \(\abs\Omega/q^d\),
summing over the \(q\) slopes would give \(\abs\Omega/q^{d-1}\), not
\(\abs\Omega/q^d\).
\end{proposition}

\begin{proof}
For \(v_E\ne0\), the vector equation determines at most one
\(\gamma\), proving injectivity.  The exact determinantal compiler says
that its slope image is precisely the transverse bad-slope set.  Counting
the incidence first by \(E\) and then by \(\gamma\) gives
\eqref{drc:eq-slope-sum}; the final scale is the corresponding uniform
fiber calculation.
\end{proof}

\begin{proposition}[The recurrence boundary is not pointwise additive]
\label{drc:prop-recurrence-nonadditive}
The exact proportionality condition above does not, by itself, turn the
locator-recurrence boundary into a subset-sum map.  Already for
\(E=\{x,y\}\), one coordinate of \(\mathcal A_E(\eta)\) is
\[
 F(x,y)=\eta_jxy-\eta_{j+1}(x+y)+\eta_{j+2}.       \tag{DRC3}
\label{drc:eq-mixed-boundary}
\]
If \(\eta_j\ne0\), this function cannot be written as
\(c+g(x)+g(y)\) whenever the four admissible two-sets
\(\{x,y\},\{x',y\},\{x,y'\},\{x',y'\}\) lie in the chart and
\((x-x')(y-y')\ne0\).
Consequently, a proposed Fourier reduction to additive maps
\(E\mapsto\sum_{t\in E}g_\gamma(t)\) requires a separate coordinate
theorem; it is not a consequence of slope linearization.
\end{proposition}

\begin{proof}
For a two-set,
\(Q_E(X)=X^2-(x+y)X+xy\), which gives
\eqref{drc:eq-mixed-boundary}.  Its mixed difference is
\[
 F(x,y)-F(x',y)-F(x,y')+F(x',y')
   =\eta_j(x-x')(y-y').
\]
Every function \(c+g(x)+g(y)\) has zero mixed difference, proving the
claim whenever the displayed product is nonzero.
\end{proof}

The maps in \eqref{drc:eq-pivot-boundary} are explicit quadratic maps in
the locator coefficients.  They depend on the received syndrome pair, as
they must in a worst-line theorem, but their number and degree are literal.
The next corollary shows how they replace the abstract ray interface.


\begin{theorem}[Fixed-core determinacy ray compiler]
\label{thm:fixed-core-determinacy-ray}
Let \(C_0\subseteq D\) have size \(c\le k\), let \(a\ge k+1\), and
fix a received line \(r=(r_0,r_1)\).  Suppose that for every
\(\gamma\) in a set \(Z\) of distinct finite MCA-bad slopes we have
chosen an exact-\(a\) noncommon witness
\((\gamma,S_\gamma,h_\gamma)\) with \(C_0\subseteq S_\gamma\).  Then
\[
             \abs Z\le \binom{n-c}{k-c+1}.          \tag{RC$_{\rm core}$}
\]
In particular, a fixed-core cell with \(k-c=o(n)\) has a direct
subexponential ray payment.  This statement is uniform in the received
line and valid in every characteristic.
\end{theorem}

\begin{proof}
Put \(d=k-c\), let
\(Q_{C_0}(X)=\prod_{x\in C_0}(X-x)\), and choose
\(g_0,g_1\in\F[X]_{<k}\) interpolating \(r_0,r_1\), respectively, on
\(C_0\).  Write \(e_i=r_i-\operatorname{ev}_D(g_i)\).  Since
\(h_\gamma-g_0-\gamma g_1\) vanishes on \(C_0\), there is a unique
\(P_\gamma\in\F[X]_{<d}\) such that
\[
 h_\gamma=g_0+\gamma g_1+Q_{C_0}P_\gamma .
\]
For \(x\in D\setminus C_0\), the witness equation is
\[
 e_0(x)+\gamma e_1(x)-Q_{C_0}(x)P_\gamma(x)=0.     \tag{*}
\]
Homogenize the parameters as
\([z:\gamma:p_0:\cdots:p_{d-1}]\in\Pj^{d+1}\), and let the row
indexed by \(x\in D\setminus C_0\) be the linear functional
\[
 z e_0(x)+\gamma e_1(x)
   -Q_{C_0}(x)\sum_{j=0}^{d-1}p_jx^j.
\]
The rows indexed by \(S_\gamma\setminus C_0\) have rank exactly
\(d+1\).  Indeed, the displayed witness parameter is a nonzero kernel
vector, so the rank is at most \(d+1\).  If the rank were at most \(d\),
the kernel would contain an independent vector.  Subtracting a multiple
of the witness vector gives a nonzero kernel vector with \(z=0\).  Its
\(\gamma\)-coordinate cannot vanish: otherwise a nonzero polynomial of
degree less than \(d\) would vanish at all
\(a-c\ge d+1\) points of \(S_\gamma\setminus C_0\).  Adding a nonzero
multiple of this kernel vector to the witness vector therefore produces a
second finite slope explained on the same support.  By fixed-support slope
uniqueness that support would be common, contrary to the witness
hypothesis.

Choose, in a fixed order, the first \((d+1)\)-subset
\(T_\gamma\subseteq S_\gamma\setminus C_0\) on which the rows have rank
\(d+1\).  Their common projective kernel is the single witness parameter,
so \(T_\gamma\) determines \(\gamma\).  Hence
\(\gamma\mapsto T_\gamma\) is injective into
\(\binom{D\setminus C_0}{d+1}\), proving the bound.  If
\(d=o(n)\), then
\(\log\binom{n-c}{d+1}=o(n)\).
\end{proof}

\begin{remark}[Cores are internal unless their census is small]
The theorem pays one fixed core.  It does not permit an atlas to use every
positive-density subset \(C_0\) as a separate profile: that census can be
exponential.  A complete atlas either routes only an algebraically
specified subexponential family of cores or keeps the actual core internal
to a coarse intersection-size profile.
\end{remark}

\subsection{Sharpness of support-only ray bounds}
\begin{theorem}[Aperiodic ray obstruction below the profile frontier]
\label{thm:aperiodic-ray-obstruction}
There is a sequence of smooth multiplicative rows
\[
 D_n=\B_n^\times\subseteq\F_n^\times,
 \qquad \log\abs{\B_n}=O(\log n)=o(n),
\]
in characteristic two, with rates tending to \(1/2\), and agreements
\(m_n=k_n+O(n/\log n)\), for which one received line has
\(2^{\Omega(n)}\) distinct finite MCA-bad slopes.  Every threshold-
\(m_n\) witness support of that line has trivial multiplicative
stabilizer, and after enlarging \(\F_n\) with only
\(\log\abs{\F_n}=O(n)\), the bad slopes and exact-agreement supports are
in bijection.  Thus smoothness, quotient-aperiodicity, Q/SP support data,
and exact agreement do not force retained-support occupancy greater than
one.  This is a sharpness result for occupancy-based compilers, not a
counterexample to a direct distinct-slope theorem at the correctly
normalized profile scale.
\end{theorem}

\begin{proof}
Let \(q_r=2^r\), put \(\B_r=\F_{q_r}\),
\(D_r=\B_r^\times\), and \(n_r=q_r-1\).  Set
\[
 m_r=2^{r-1}-1=\frac{n_r-1}{2}.
\]
Fix \(0<\eta<\log2\), take
\(w_r=\floor{\eta n_r/\log q_r}\), and put
\(k_r=m_r-w_r-1\).  Then \(k_r/n_r\to1/2\) and
\(m_r-k_r=w_r+1=O(n_r/\log n_r)\).

Pigeonholing the \(m_r\)-set locators over their \(q_r^{w_r}\) prefixes
gives a fiber \(\Fcal_r\) satisfying
\[
 \log\abs{\Fcal_r}
 \ge \log\binom{n_r}{m_r}-w_r\log q_r
 \ge(\log2-\eta)n_r-o(n_r).
\]
Moreover
\[
 \gcd(m_r,n_r)=
 \gcd(2^{r-1}-1,2^r-1)=2^{\gcd(r-1,r)}-1=1.
\]
If an \(m_r\)-support were invariant under a nontrivial subgroup of
\(D_r\) of order \(c>1\), it would be a union of \(c\)-element orbits,
forcing \(c\mid m_r\) and \(c\mid n_r\), a contradiction.  Thus every
member of \(\Fcal_r\) is aperiodic.

Choose a finite extension \(\F_r/\B_r\) with
\[
 \abs{\F_r}>n_r+k_r\binom{\abs{\Fcal_r}}2.
\]
Since \(\abs{\Fcal_r}\le2^{n_r}\), the least extension degree meeting
this inequality has \(\log\abs{\F_r}=O(n_r)\).  By
\cref{cor:exact-prefix-ray-realization}, a separating pole gives exactly
\(\abs{\Fcal_r}=\exp(\Omega(n_r))\) distinct bad slopes, one for each
aperiodic support, and every occupancy equals one.
\end{proof}

\begin{remark}[Scope of the aperiodic obstruction]
The construction lies below the entropy--profile crossing: its natural
prefix scale is itself exponential.  It therefore does not contradict a
profile-envelope upper bound proved at the correct scale.  It does rule
out a field-independent polynomial ray bound, division of SP by the full
support slice without a retained multiplicity theorem, and the assertion
that quotient-aperiodicity alone forces ray collisions.
\end{remark}

The opposite projection extreme occurs over the coefficient field itself.
It shows that a large aperiodic prefix fiber need not supply a lower
reserve after first-match saturation.

\begin{theorem}[Aperiodic one-ray saturation]
\label{thm:aperiodic-one-ray-saturation}
For an integer \(\nu\ge5\), put
\[
 q=2^\nu,\quad D=\F_q^\times,\quad n=q-1,\quad
 m=2^{\nu-1}-1,\quad w=\floor{n/\nu^2},\quad k=m-w-1.
\]
Refine the depth-\(w\) prefix by the locator constant coefficient:
\[
 \Psi_\nu(S)=\bigl(c_1(S),\ldots,c_w(S),c_m(S)\bigr)
 \qquad\left(S\in\binom Dm\right).
\]
There are \(z\in\F_q^w\) and \(c\in\F_q^\times\) for which
\[
 \Gcal_{z,c}=\{S:\Psi_\nu(S)=(z,c)\},
 \qquad
 \abs{\Gcal_{z,c}}
 \ge\ceil{\binom nm q^{-(w+1)}}
 =\exp\bigl((\log2-o(1))n\bigr).                  \tag{SAT1}
\label{eq:one-ray-large-fiber}
\]
Every support in \(\Gcal_{z,c}\) has trivial multiplicative stabilizer,
but all witnesses arising from the exact prefix list project at the pole
\(\alpha=0\) in the coefficient field \(\F_q\) to the single MCA slope \(-c\) for
\(\RS_{\F_q}(D,k)\).

More precisely, if \(\Fcal_z=\Phi_w^{-1}(z)\) and
\(\Ccal_d=\{S\in\Fcal_z:c_m(S)=d\}\), then the nonempty
\(\Ccal_d\), \(d\in\F_q^\times\), partition the supports of all
exact-\(m\) witnesses of this pole line, and the witnesses supported on
\(\Ccal_d\) have slope image exactly \(\{-d\}\).  Thus at most
\(q-1=\exp(o(n))\) explicit saturation cells
pay the complete line even though one cell contains exponentially many
aperiodic supports.
\end{theorem}

\begin{proof}
The map \(\Psi_\nu\) has at most \(q^{w+1}\) values, so pigeonholing
proves the first inequality in \eqref{eq:one-ray-large-fiber}.  Moreover
\[
 \log\binom nm=n\log2-O(\log n),\qquad
 (w+1)\log q\le(n/\nu+\nu)\log2=o(n),
\]
which proves its asymptotic form.  Also
\[
 \gcd(m,n)=\gcd(2^{\nu-1}-1,2^\nu-1)
           =2^{\gcd(\nu-1,\nu)}-1=1.
\]
An invariant support would be a union of orbits of a nontrivial subgroup
whose order divides both \(m\) and \(n\), so every \(m\)-support is
multiplicatively aperiodic.

Let \(U_z=X^m+\sum_{i=1}^wz_iX^{m-i}\) and
\(P_S=U_z-Q_S\).  Prefix cancellation gives \(\deg P_S\le k\), and
\(U_z-P_S=Q_S\) shows that \(P_S\) agrees with \(U_z\) exactly on \(S\).
At the only pole of the coefficient field outside \(D\), namely zero,
\[
                     P_S(0)=U_z(0)-Q_S(0)=-c_m(S).
\]
The pole-line calculation therefore maps every \(\Ccal_d\) to \(-d\).
Conversely, an exact-\(m\) explanation on this line gives a polynomial
\(P=Xh+\gamma\) of degree at most \(k\) agreeing with \(U_z\) on
\(m\) points.  The exact prefix-list bijection makes it the unique
\(P_S=U_z-Q_S\) for some \(S\in\Fcal_z\), and then
\(\gamma=P(0)=-c_m(S)\).  This proves coverage and the asserted
one-slope images.
\end{proof}

\begin{remark}[Two incompatible aperiodic inferences]
\label{rem:two-aperiodic-ray-extremes}
\Cref{thm:aperiodic-ray-obstruction} makes a separating extension pole
inject an exponentially large aperiodic fiber into distinct slopes, while
\cref{thm:aperiodic-one-ray-saturation} makes a coefficient-field pole collapse
such a refined fiber to one slope.  Aperiodicity alone therefore implies
neither ray expansion nor ray collision.  A lower compiler needs a
separating-pole estimate or a quantitative exclusion of earlier saturation
cells; an upper compiler needs an actual determinant or ray-image bound.
\end{remark}

\subsection{Smooth and circle evaluation domains}\label{sec:smooth-circle-domains}

This subsection records the domain geometry behind the exact quotient
constructions.  The multiplicative and circle cases are treated simultaneously through
uniform-fiber folding maps.  The point is not merely that such maps exist,
but that their fibers are complete inside the evaluation domain and their
locators have a \emph{lacunary} top part, meaning that prescribed gaps occur
among the highest-degree coefficients.  These two facts make quotient
families visible in the witness moduli and keep their locator prefixes
stable under planted cores.

\begin{definition}[Complete-fiber folding map and smooth multiplicative coset]
\label{def:structured-folding}
Let \(D\subseteq\B\) be finite.  A polynomial map
\(\varphi:D\to Q:=\varphi(D)\) is a \emph{\(c\)-fold
complete-fiber folding map} if every \(y\in Q\) has exactly \(c\)
preimages in \(D\).  A support is \emph{complete at this folding scale}
if it is a union of whole fibers, equivalently if it is
\(\varphi^{-1}(E)\) for some \(E\subseteq Q\).  Passing from \(D\) to
\(Q\), and from a complete support to \(E\), is the corresponding
\emph{quotient} or \emph{folding}.

A \emph{multiplicative coset} is a set \(D=\theta H\subseteq\B^\times\),
where \(H\le\B^\times\) is a subgroup and \(\theta\in\B^\times\).
For every divisor \(c\mid\abs H\), the restriction
\[
             \pi_c:D\longrightarrow D^{(c)},\qquad
             \pi_c(x)=x^c,\qquad
             D^{(c)}=\{x^c:x\in D\},
\]
is a \(c\)-fold complete-fiber folding map.  We call \(D\)
\emph{smooth at scale \(c\)} when this power folding is among the
quotients retained by the row, and \emph{smooth} for a family of scales
when it is smooth at every scale in that family.  Thus the adjective
\emph{smooth} here is an explicit uniform-fiber property, not a claim
about nonsingularity of an algebraic variety.  The locator of a complete
support is \(V_E(X^c)\); its nonleading terms occur only in degree gaps
divisible by \(c\).  This controlled pattern of missing top coefficients
is what \emph{lacunary} means in this section.
\end{definition}

\begin{definition}[Circle domain and circle code]
\label{def:circle-domain-code}
Let \(\B_0\) have odd characteristic and let
\(\mathcal C/\B_0\) be the affine conic
\(\mathcal C:x^2+y^2=1\).  A \emph{standard circle domain} is a finite
set \(E\subseteq\mathcal C(\B_0)\).  If \(\F\supseteq\B_0\), the
degree-\(w\) \emph{circle code} on \(E\) is
\[
 \mathcal C_w(\F,E)
 =\{(f(P))_{P\in E}:
      f\in\operatorname{im}(\F[x,y]_{\le w}
          \to\F[x,y]/(x^2+y^2-1))\}.
\]
The degree is the total-degree filtration before passage to the quotient
ring.  A \emph{circle \(x\)-domain} is instead a set of first
coordinates obtained from the torus construction below; an RS code on
that \(x\)-domain and a bivariate circle code are different presentations
until a diagonal equivalence such as
\cref{lem:circle-uniformization} has been exhibited.
\end{definition}

\begin{definition}[Norm-one torus, twin coset, and Chebyshev domain]
\label{def:circle-twin-domain}
Following the circle-domain geometry of \cite{HLP24}, let
\(p\equiv3\pmod4\), choose \(i\in\F_{p^2}\) with \(i^2=-1\), and
write
\[
 \UU=\ker\!\left(N_{\F_{p^2}/\F_p}:\F_{p^2}^{\times}
                      \to\F_p^\times\right)
     =\{u\in\F_{p^2}^{\times}:u^{p+1}=1\}.
\]
The algebraic kernel of the norm morphism is the
\emph{norm-one torus}; \(\UU\) is its cyclic group of
\(\F_p\)-rational points.  Put
\(\chi(u)=(u+u^{-1})/2\in\F_p\).  If \(H\le\UU\) and
\(g\in\UU\) satisfies \(g^2\notin H\), then
\[
 \Dcal(g,H)=gH\sqcup g^{-1}H
\]
is a disjoint union exchanged by \(u\mapsto u^{-1}\); it is called a
\emph{twin coset}.  Its projected set
\[
             D(g,H)=\chi(\Dcal(g,H))\subseteq\F_p
\]
has cardinality \(\abs H\) and is the \emph{Chebyshev twin-coset
\(x\)-domain}.

For \(c\ge1\), let \(T_c\) be the Chebyshev polynomial normalized by
\[
       T_c\!\left(\frac{u+u^{-1}}2\right)
       =\frac{u^c+u^{-c}}2 .
\]
Thus \(T_c\circ\chi=\chi\circ(u\mapsto u^c)\).  Whenever
\(c\mid\abs H\) and the two image cosets remain disjoint, the restriction
of \(T_c\) to \(D(g,H)\) is the \emph{Chebyshev folding map}; a
Chebyshev twin-coset domain equipped with these complete-fiber maps is
the circle analogue of a smooth multiplicative coset.
\end{definition}

\begin{definition}[Map-smooth domain]\label{def:map-smooth}
Let \(\B\subseteq\F\) be finite fields, let \(\varphi\in\B[X]\) have
degree \(c\ge1\), and let \(D\subseteq\B\) have size \(n\).  The set
\(D\) is \((\varphi,c)\)-smooth if every nonempty fiber
\(\varphi^{-1}(y)\cap D\), \(y\in\varphi(D)\), has exactly \(c\)
elements.  We write \(Q=\varphi(D)\) and \(N=n/c\).
\end{definition}

A multiplicative coset \(D=\theta H\subseteq\B^\times\) is \((X^a,a)\)-smooth for every divisor \(a\mid\abs H\).  The circle case uses Chebyshev maps rather than power maps; see \cref{lem:cheb-smooth}.  The following lemma is the common quotient-support construction.  It is included because it produces explicit large lists and because the same lacunarity controls their locator prefixes.

\begin{lemma}[Map-smooth fiber construction]\label{lem:map-smooth-fiber}
Let \(D\) be \((\varphi,a)\)-smooth, \(\abs D=n\), and \(C^+=\RS_\F(D,k+1)\).  Put \(\ell=\lfloor k/a\rfloor+2\) and \(A=a\ell\).  If \(\ell\le N-1\), then there is \(z\in\B\) and a word \(u_z\in\B^D\) such that \(u_z\) has at least \(\binom N\ell/\abs\B\) distinct codewords of \(C^+\) agreeing with it on at least \(A\) positions.  Moreover \(k+a+1\le A\le k+2a\), with equality \(A=k+2a\) if \(a\mid k\).
\end{lemma}

\begin{proof}
For each \(\ell\)-subset \(E\subseteq Q\), let \(P_E(Y)=\prod_{b\in E}(Y-b)=Y^\ell-e_1(E)Y^{\ell-1}+R_E(Y)\), with \(\deg R_E\le\ell-2\).  The composed locator \(P_E(\varphi(X))\) vanishes on the union of the \(a\)-point fibers above \(E\), hence on \(A=a\ell\) points of \(D\).  Set \(z_E=-e_1(E)\) and \(c_E=-R_E(\varphi)\).  Since \(\deg c_E\le a(\ell-2)\le k\), the word \(c_E\) lies in \(C^+\), and it agrees on that fiber union with \(u_{z_E}(x)=\varphi(x)^\ell+z_E\varphi(x)^{\ell-1}\).  Among the \(\binom N\ell\) subsets \(E\), some value of \(z_E\in\B\) occurs at least \(\binom N\ell/\abs\B\) times.  If two such subsets with the same \(z\) gave the same codeword, then \(R_E(\varphi)=R_{E'}(\varphi)\), hence \(R_E=R_{E'}\) because composition with a nonconstant polynomial is injective, and then \(P_E=P_{E'}\).  Thus the codewords are distinct.  The bounds on \(A\) follow from \(\ell=\lfloor k/a\rfloor+2\).
\end{proof}

\begin{proposition}[Collision-aware lower cap for map-smooth domains]\label{prop:map-smooth-cap}
Let \(D\) be \((\varphi,a)\)-smooth, \(n<\abs\F=q\), and let
\(C=\RS_\F(D,k)\).  With \(\ell,A\) as in
\cref{lem:map-smooth-fiber}, assume \(\ell\le N-1\), and put
\[
 L=\left\lceil\binom N\ell\abs\B^{-1}\right\rceil.
\]
Then for every integer threshold \(k+1\le m\le A\),
\[
 B_C^{\rm MCA}(m)\ge
 \left\lceil\frac{L(q-n)}{q-n+k(L-1)}\right\rceil .
\]
\end{proposition}

\begin{proof}
The list in \cref{lem:map-smooth-fiber} has size at least \(L\) at
agreement \(A\).  Apply the collision-aware pole conversion
\cref{thm:collision-aware-pole}.  Its line pair is not simultaneously
explained on any \(k+1\) points, so an agreement set of size \(A\) also
witnesses badness for every threshold \(m\) in the displayed range.
\end{proof}

We now verify the complete-fiber property for the construction in
\cref{def:circle-twin-domain}.  Thus
\(\Dcal=gH\sqcup g^{-1}H\subseteq\UU\),
\(D=\chi(\Dcal)\), and \(T_a(\chi(u))=\chi(u^a)\).

\begin{lemma}[Chebyshev smoothness]\label{lem:cheb-smooth}
Let \(\Dcal=gH\cup g^{-1}H\subseteq\UU\) be a twin coset, let \(D=\chi(\Dcal)\), and let \(a\mid\abs H\).  If \(g^{2a}\notin H^{(a)}=\{h^a:h\in H\}\), then \(D\) is \((T_a,a)\)-smooth.  If \(\ord(g)\nmid2\abs H\), this holds for all \(a\mid\abs H\).
\end{lemma}

\begin{proof}
The map \(u\mapsto u^a\) has kernel of size \(a\) on \(\UU\), and that kernel lies in \(H\).  Hence it maps \(gH\) and \(g^{-1}H\) onto \(g^aH^{(a)}\) and \(g^{-a}H^{(a)}\), respectively, with fibers of size \(a\).  The hypothesis says these image cosets are disjoint, so the image is again a twin coset.  Since \(\chi\) is exactly two-to-one on both twin cosets and identifies only inverse points, the equation \(T_a(x)=y\) on \(D\) has exactly \(a\) solutions for each \(y\in T_a(D)\).  Taking \(a=\abs H\) shows that the all-scale condition is equivalent to \(g^{2\abs H}\ne1\), and the implication to every smaller divisor follows by raising to \(\abs H/a\).
\end{proof}

\begin{lemma}[Diagonal uniformization of circle codes]\label{lem:circle-uniformization}
Assume \(\F\supseteq\F_{p^2}\).  Let
\(E\subseteq\{(x,y)\in\F_p^2:x^2+y^2=1\}\), define
\(\iota(x,y)=x+iy\), put \(E'=\iota(E)\subseteq\UU\), and let
\(\Ccal_w(\F,E)\) be the circle code of
\cref{def:circle-domain-code}.  Under the coordinate bijection \(\iota\)
and the diagonal multiplier \(t_u=u^{-w}\), the code
\(\Ccal_w(\F,E)\) is \(t\circ\RS_\F(E',2w+1)\).
Consequently list sizes, ordinary correlated-agreement (CA) numerators,
and mutual-correlated-agreement (MCA) numerators are preserved.
\end{lemma}

\begin{proof}
The map \(u=x+iy\) identifies \(\F[x,y]/(x^2+y^2-1)\) with \(\F[u,u^{-1}]\).  The subspace of total degree at most \(w\) maps to \(u^{-w}\F[u]_{\le2w}\), and both spaces have dimension \(2w+1\).  Multiplication by the nonzero diagonal vector \((u^{-w})_{u\in E'}\) preserves coordinatewise equality and hence every agreement notion used here.  Therefore it preserves list, CA, and MCA counts exactly.
\end{proof}

The preceding uniformization is convenient when \(i\) belongs to the
challenge field.  The following form works over every odd-characteristic
field and will also be used for scalar extensions of odd degree.

\begin{lemma}[Invariance under monomial equivalence]
\label{lem:monomial-mca-invariance}
Let \(C\subseteq\F^D\), let \(\pi:D'\to D\) be a bijection, and let
\((d_x)_{x\in D'}\in(\F^\times)^{D'}\).  Define
\[
 C'=\{(d_xc_{\pi(x)})_{x\in D'}:c\in C\}.
\]
The same coordinate permutation and diagonal multiplication give a
bijection between received pairs for \(C\) and \(C'\), preserving the
bad-slope set on every support and for every challenge set.  In
particular \(B_{C,\Gamma}^{\rm MCA}(a)=B_{C',\Gamma}^{\rm MCA}(a)\).
\end{lemma}

\begin{proof}
For every word \(u\), codeword \(c\), and support \(S\), equality
\(u|_S=c|_S\) is equivalent after the displayed coordinate change to
equality on \(\pi^{-1}(S)\).  The transformation is \(\F\)-linear, so it
commutes with \(r_0+\gamma r_1\).  It therefore preserves both the
single-line-point explanation and simultaneous explanation of the pair
on the corresponding support, slope by slope.
\end{proof}

\begin{corollary}[Generalized Reed--Solomon rows]
\label{cor:grs-transfer}
Every generalized Reed--Solomon code obtained from \(\RS_\F(D,k)\) by
a coordinate permutation and nonzero column multipliers has exactly the
same CA and MCA numerators as that RS code.  In particular all exact
staircases and target compilers in this paper transfer verbatim.
\end{corollary}

\begin{proof}
This is \cref{lem:monomial-mca-invariance} applied to the defining
monomial equivalence.
\end{proof}

\begin{proposition}[Stereographic circle equivalence]
\label{prop:stereographic-circle-equivalence}
Let \(\F\) have odd characteristic, let
\(E\subseteq\{(x,y)\in\F^2:x^2+y^2=1\}\) avoid \((-1,0)\), and put
\[
 t=\frac{y}{1+x},\qquad T=t(E).
\]
If \(2w+1\le\abs E\), then
\[
 \Ccal_w(\F,E)
  =\diag\bigl((1+t^2)^{-w}\bigr)_{t\in T}
       \RS_\F(T,2w+1)                             \tag{SC1}
 \label{eq:stereographic-code-equivalence}
\]
after the coordinate bijection \(E\to T\).  Hence every list, CA, and
support-wise MCA bad-slope set is preserved exactly.
\end{proposition}

\begin{proof}
The inverse parametrization is
\[
 x=\frac{1-t^2}{1+t^2},\qquad
 y=\frac{2t}{1+t^2},
\]
and \(1+t^2=2/(1+x)\ne0\) on \(E\).  Modulo the circle equation, the
degree-at-most-\(w\) filtered space is
\[
 \{f_0(x)+yf_1(x):\deg f_0\le w,\ \deg f_1\le w-1\},
\]
with dimension \(2w+1\).  After substitution and multiplication by
\((1+t^2)^w\), every such function becomes a polynomial of degree at
most \(2w\).  The substitution is injective in the localized coordinate
ring and the dimensions agree, so all of \(\F[t]_{\le2w}\) occurs.
Equation \eqref{eq:stereographic-code-equivalence} follows, and
\cref{lem:monomial-mca-invariance} proves the last assertion.
\end{proof}

\begin{corollary}[Exact circle staircase]
\label{cor:exact-circle-staircase}
Let \(E,w\) be as in \cref{prop:stereographic-circle-equivalence}, put
\(n=\abs E\) and \(k=2w+1<n\).  For every nonempty
\(\Gamma\subseteq\F\) and every \(0\le r\le n-k-1\),
\[
 (n-r)^2\ge n(k+r)
 \quad\Longrightarrow\quad
 B_{\Ccal_w(\F,E),\Gamma}^{\rm MCA}(n-r)
      =\min\{\abs\Gamma,r+1\}.                  \tag{SC2}
\]
The same equality holds under the hypotheses of
\cref{thm:literature-integral-completion}.  Consequently standard circle
codes admitting such a projection chart have the positive-relative-radius staircase of
\cref{cor:positive-radius-density}.
\end{corollary}

\begin{proof}
Apply \cref{lem:monomial-mca-invariance} and then
\cref{thm:quadratic-mean-overlap,thm:literature-integral-completion} to
the RS presentation in
\eqref{eq:stereographic-code-equivalence}.
\end{proof}

For a set \(\Pcal\) of primes, its \emph{relative prime density} is
\[
\operatorname{dens}_{\rm pr}(\Pcal)=
 \lim_{X\to\infty}\frac{\#\{p\le X:p\in\Pcal\}}{\pi(X)},
\]
when the limit exists.  The prime number theorem in arithmetic
progressions gives density \(1/\varphi(M)\) to every reduced residue
class modulo \(M\) \cite{Davenport}.
We write \(\overline{\operatorname{dens}}_{\rm pr}\) for the same
expression with \(\limsup\) in place of the limit.

\begin{theorem}[Fixed-length positive-density prime families]
\label{thm:fixed-length-prime-density}
Fix \(n=2^m\), \(m\ge2\).
\begin{enumerate}[label=\textup{(\alph*)},leftmargin=*]
\item The primes \(p\equiv1\pmod n\), of relative prime density
\(1/\varphi(n)=2/n\), support the unique subgroup
\(H\le\F_p^\times\) of order \(n\).  Every exact RS conclusion of
\cref{thm:quadratic-mean-overlap,thm:literature-integral-completion}
holds on \(H\) and on each of its multiplicative cosets.
\item The primes \(p\equiv-1\pmod {2n}\), of density
\(1/\varphi(2n)=1/n\), support a standard circle coset \(E\) of size
\(n\).  For every odd \(k=2w+1<n\), the circle code
\(\Ccal_w(\F_{p^2},E)\) has both exact staircases.
\item The primes \(p\equiv-1\pmod {4n}\), of density
\(1/\varphi(4n)=1/(2n)\), support a Chebyshev-smooth circle
\(x\)-domain \(D\subseteq\F_p\) of size \(n\), and
\(\RS_{\F_p}(D,k)\) has both exact staircases for every \(k<n\).
\end{enumerate}
\end{theorem}

\begin{proof}
Part~\textup{(a)} follows from cyclicity of \(\F_p^\times\).  For
part~\textup{(b)}, let \(\UU\le\F_{p^2}^\times\) be the norm-one group,
of order \(p+1\), choose \(g\in\UU\) of order \(2n\), and put
\(H=\langle g^2\rangle\).  The coset \(E'=gH\) has size \(n\), is
stable under inversion, and avoids \(-1\).  Its inverse image under
\(u=x+iy\) consists of \(n\) points of the circle over \(\F_p\).
\Cref{lem:circle-uniformization} identifies its circle code with a
diagonal translate of \(\RS_{\F_{p^2}}(E',2w+1)\).

For part~\textup{(c)}, choose \(g\in\UU\) of order \(4n\), put
\(H=\langle g^4\rangle\), and use the twin coset
\(gH\sqcup g^{-1}H\).  Its \(\chi\)-image has size \(n\), and
\(\ord(g)=4n\nmid2n\), so \cref{lem:cheb-smooth} proves smoothness at
every dyadic scale.  The density assertions are the prime number theorem
in the three displayed reduced residue classes.
\end{proof}

\begin{corollary}[Finite smooth and circle families beyond the prior window]
\label{cor:circle-32-family}
Let \(n=32,k=15,r=6\).  For every prime
\(p\equiv1\pmod {32}\), let \(H\le\F_p^\times\) have order \(32\)
and put \(C_p^+=\RS_{\F_p}(H,15)\).  For every prime
\(p\equiv-1\pmod {64}\), construct the circle coset \(E\) of
\cref{thm:fixed-length-prime-density}\textup{(b)} and put
\(C_p^-=\Ccal_7(\F_{p^2},E)\).  Then
\[
 B_{C_p^+}^{\rm MCA}(26)=B_{C_p^-}^{\rm MCA}(26)=7. \tag{SC3}
\]
At targets \(6/p\) and \(6/p^2\), respectively, both complete safe
sets are \([0,3/16)\), with unsafe endpoint.  The two supporting prime
sets have relative densities \(1/16\) and \(1/32\).
\end{corollary}

\begin{proof}
Here \(R=17\) and
\[
 (n-r)^2=26^2=676>672=32(15+6).
\]
Thus \cref{thm:quadratic-mean-overlap,cor:exact-circle-staircase} give
the numerator seven for the smooth and circle rows.  Radius five is in
the deep range and has numerator six; the tangent floor at radius six and
monotonicity give the stated safe sets.  This example lies
beyond the elementary one-third range because \(3r=18>R\).  It is not
certified by the specialized BCHKS estimate
\eqref{eq:bchks-specialized}, whose first term is \(176/15>7\).  The respective densities
are \(1/\varphi(32)=1/16\) and \(1/\varphi(64)=1/32\).
\end{proof}

\begin{corollary}[Concrete Mersenne-31 circle rows]
\label{cor:mersenne-circle-rows}
Let \(p=2^{31}-1\), the Mersenne prime used in \cite{HLP24}, and let
\(\UU\) be its norm-one circle group.
Choose a generator \(g\), put \(H_0=\langle g^2\rangle\), and let
\(E\subseteq\{(x,y)\in\F_p^2:x^2+y^2=1\}\) be the circle domain
corresponding to \(gH_0\), so \(\abs E=2^{30}\).  This domain avoids
\((-1,0)\).  Over \(\F_{p^5}\), the circle code with
\(w=2^{28}\) and \(k=2^{29}+1\) has, at target \(2^{-128}\),
the exact safe set
\[
\left[0,\frac18-2^{-30}\right).                 \tag{SC4}
\label{eq:mersenne-circle-safe}
\]
Moreover, if \(H=\langle g^4\rangle\), the Chebyshev twin-coset
\(gH\sqcup g^{-1}H=gH_0\) has \(x\)-domain
\(D=\chi(gH_0)\) of size \(2^{29}\), and
\(\RS_{\F_{p^5}}(D,2^{27})\) has exact safe set
\[
\left[0,\frac14-2^{-29}\right).                 \tag{SC5}
\label{eq:mersenne-x-safe}
\]
The first code has rate slightly above \(1/2\); the second has exact
rate \(1/4\) but is a circle-smooth \(x\)-domain rather than a literal
multiplicative-domain Prize row.
\end{corollary}

\begin{proof}
The group \(\UU\) has order \(p+1=2^{31}\).  Since
\(-1=g^{2^{30}}\in H_0\), neither \(1\) nor \(-1\) belongs to the
disjoint coset \(gH_0\); in particular \(E\) avoids the pole
\((-1,0)\) of the stereographic chart.  Also
\(gH_0=gH\sqcup g^{-1}H\).  Inversion acts freely on this coset and
\(\chi(u)=(u+u^{-1})/2\) identifies precisely the pairs
\(\{u,u^{-1}\}\), proving \(\abs D=2^{29}\).  The relation
\(T_a\circ\chi=\chi\circ(u\mapsto u^a)\) and
\cref{lem:cheb-smooth} give the asserted dyadic smoothness.

Binomial expansion gives
\(B=\floor{p^5/2^{128}}=2^{27}-1\).  At radius \(B-1\), the quadratic
margin \((n-r)^2-n(k+r)\) for the bivariate row is
\(162129591417176068>0\), so
\cref{cor:exact-circle-staircase} and the tangent floor give
\eqref{eq:mersenne-circle-safe}.  Here \(p^5\equiv3\pmod4\), so
\(i\notin\F_{p^5}\); the applicable equivalence is the stereographic
one in \cref{prop:stereographic-circle-equivalence}, not torus
uniformization.  For the \(x\)-domain row the same budget lies inside
the quadratic staircase at rate \(1/4\), since its margin at radius
\(B-1\) is \(18014401193836548>0\).  The identical target argument gives
\eqref{eq:mersenne-x-safe}.
\end{proof}

\begin{corollary}[Every odd-characteristic extension-field tower]
\label{cor:density-one-extension-towers}
Fix an odd prime \(p\), put \(n_e=2^e\) and \(q_e=p^{n_e}\), and assume
\(e\ge7\).  Then \(4n_e\mid q_e-1\).  Over \(\F_{q_e}\) there exist
both a multiplicative subgroup RS row of length \(n_e\) and a
Chebyshev-smooth circle \(x\)-domain row of that length.  At every
official rate and every \(0\le r\le r_\rho^\sharp\), both satisfy
\[
 e_{\rm MCA}(C_e,r/n_e)=\frac{r+1}{q_e}.           \tag{SC6}
\]
Thus the construction exists in every odd characteristic and beats every
fixed error target throughout every compact subinterval below half the
minimum distance for large \(e\).
\end{corollary}

\begin{proof}
The lifting-the-exponent identity gives
\[
 v_2(p^{2^e}-1)=v_2(p-1)+v_2(p+1)+e-1\ge e+2.
\]
Cyclicity produces the multiplicative subgroup and an element of order
\(4n_e\).  Since \(n_e\) is even, \(q_e\equiv1\pmod4\); choose
\(i\in\F_{q_e}\) with \(i^2=-1\).  The map
\[
 u\longmapsto
 \left(\frac{u+u^{-1}}2,\frac{u-u^{-1}}{2i}\right)
\]
identifies \(\F_{q_e}^\times\) with the split circle over
\(\F_{q_e}\).  If \(g\) has order \(4n_e\) and
\(H=\langle g^4\rangle\), the twin coset
\(gH\sqcup g^{-1}H\) projects to an \(n_e\)-point \(x\)-domain.
The proof of \cref{lem:cheb-smooth}, which uses only cyclicity and
\(T_a\circ\chi=\chi\circ(u\mapsto u^a)\), applies verbatim and proves
Chebyshev smoothness at every dyadic scale.
Apply \cref{cor:official-rate-half-staircase}.  Since
\(q_e=p^{n_e}\), the normalized error tends to zero exponentially.
\end{proof}

\begin{proposition}[No growing positive-density prime-field family]
\label{prop:no-growing-prime-density}
Let \(\Pcal\) be a set of primes and assign a dyadic integer
\(n_p\to\infty\) along \(\Pcal\).  If \(n_p\mid p-1\) for every
\(p\in\Pcal\), or \(n_p\mid p+1\) for every \(p\in\Pcal\), then
\(\Pcal\) has upper relative prime density zero.  If in addition
\(p/n_p\) is bounded, there are only \(O(\log X)\) possible primes of
the indicated form below \(X\).
\end{proposition}

\begin{proof}
For every fixed \(E\), all sufficiently large assigned lengths are
divisible by \(2^E\).  Hence the primes eventually lie in the single
class \(1\pmod {2^E}\), or the single class \(-1\pmod {2^E}\), whose
relative prime density is \(1/\varphi(2^E)=2^{1-E}\).  Letting
\(E\to\infty\) proves the first assertion.  Under the bounded-ratio
hypothesis, \(p\mp1=a2^e\) with \(a\) in a fixed finite set; there are
only \(O(\log X)\) choices of \((a,e)\) below \(X\).
\end{proof}

\begin{remark}[Prize scope]
The prime-density statements fix \(n\); their densities tend to zero as
\(n\) grows.  The extension towers have exponential field size and soon
leave the \(q<2^{256}\) Prize box.  Finally the target \(6/p^2\) in
\cref{cor:circle-32-family} varies with \(p\).  None of these facts may be
rephrased as a positive-density growing prime-field solution at the fixed
target \(2^{-128}\).
\end{remark}

\subsection{Quotient--remainder profiles and partial occupancy}
\label{sec:quotient-obstruction}

The quotient and remainder variables can be separated exactly at the
locator-prefix level.  This removes a spurious factor coming from summing
over marked remainder sets, but it also exposes a genuine scaled quotient-
coefficient-field term which prevents an unrestricted comparison with the
identity profile.

\begin{definition}[Quotient/remainder profile and field drop]
\label{def:quotient-remainder-profile}
Let \(\phi:D\to Q\) be a \(c\)-fold complete-fiber folding map in the
sense of \cref{def:structured-folding}.  A
\emph{complete-fiber-plus-remainder support} is
\[
                  S=\phi^{-1}(E)\sqcup R,
\]
where \(E\subseteq Q\) and
\(R\subseteq D\setminus\phi^{-1}(E)\).  The set \(R\) is the
\emph{support remainder}: it records the selected points in fibers that
are not declared complete.  A \emph{quotient/remainder profile} fixes
\(\phi\), its degree \(c\), the size \(r=\abs R\), the relevant
coefficient field, and any planted remainder data, while \(E\) and the
unfixed part of \(R\) range subject to the displayed disjointness.
A fixed-\(R\) profile is the subprofile in which the remainder support
itself is part of the label.  The adjective \emph{Euclidean} used below
refers to separating the locator factor of degree \(r<c\) from the
composed quotient locator; it does not mean that an arbitrary support
remainder has size less than \(c\).

We call \(\B_\phi\subseteq\B\) a \emph{scaled quotient coefficient
field} when, for some \(\eta\in\B^\times\),
\[
                         Q=\phi(D)\subseteq\eta\B_\phi .
\]
When this field is recorded in a profile, it is taken minimal among the
subfields with this property.
Then the \(j\)-th elementary quotient coefficient lies in
\(\eta^j\B_\phi\).  A \emph{field drop} is the use of a proper such
subfield.  It is a \emph{positive-rate field drop} along a sequence when
\[
       1-\frac{\log\abs{\B_\phi}}{\log\abs\B}
\]
stays bounded away from zero and a positive-density number on the
entropy scale of quotient-prefix slots is live.  A field drop changes the number of available coefficient values and is therefore recorded explicitly.
\end{definition}

\begin{definition}[Balanced quotient core]
\label{def:balanced-quotient-core}
A \emph{split locator} is a monic squarefree polynomial all of whose roots
lie in the evaluation domain \(D\).
After common locator factors, fixed quotient data, and planted remainder
points have been removed, let \(A\) and \(B\) be two monic split
locators of the same degree \(e\).  They form a
\emph{depth-\(w\) balanced core} if they have the same first \(w\)
nonleading coefficients, equivalently
\[
                         \deg(A-B)\le e-w-1.
\]
It is a \emph{balanced quotient core} when the remaining locator factors
are pullbacks through the same folding map \(\phi\).  The word
\emph{balanced} records equality of degree and leading normalization;
the \emph{core} consists of the roots still allowed to move after common
and planted factors are removed.  If the resulting projective locator
family is one-dimensional it is a split pencil
\(Q_0+\lambda Q_1\); in higher dimension this definition by itself
supplies no distinct-ray bound.  The exact determinant-pivot compiler
still applies, but its boundary maps require the analytic payment stated
in \textup{(A6)}.
\end{definition}

For a monic polynomial
\[
 A(X)=X^d+a_1X^{d-1}+\cdots+a_d
\]
and an integer $0\le u\le d$, write
\(\operatorname{pref}_u(A)=(a_1,\ldots,a_u)\).  Equivalently, if
\(A^\vee(Z):=Z^dA(Z^{-1})\), then
\(\operatorname{pref}_u(A)\) is the coefficient vector of
\(A^\vee(Z)\bmod Z^{u+1}\), with its constant coefficient omitted.

\begin{theorem}[Exact quotient--remainder prefix normal form]
\label{thm:exact-quotient-remainder-normal-form}
Let $\B\subseteq\F$ be finite fields, let $D\subseteq\B$ have $n$
distinct elements, and let $\phi\in\B[X]$ be monic of degree $c\ge2$.
Assume that the restriction
\[
             \phi:D\longrightarrow Q:=\phi(D)
\]
has complete fibers of size $c$; in particular, $N:=|Q|=n/c$.
Fix integers $m,r$ with $0\le r<c$ and $cm+r\le n$.  For
\[
 E\in\binom Qm,
 \qquad
 R\in\binom{D\setminus\phi^{-1}(E)}r,
\]
put
\[
 \begin{split}
 V_E(Y)&:=\prod_{y\in E}(Y-y)
       =Y^m+\sum_{j=1}^m v_j(E)Y^{m-j},\\
 P_R(X)&:=\prod_{x\in R}(X-x)
       =X^r+\sum_{i=1}^r p_i(R)X^{r-i},\\
 L_{E,R}(X)&:=P_R(X)V_E(\phi(X)).
 \end{split}                                      \tag{QR1}\label{eq:qr-locators}
\]
Then $L_{E,R}$ is the monic locator of the support
\(\phi^{-1}(E)\sqcup R\), of degree $A=cm+r$.

Let $0\le w\le A$ and put $d=\lfloor w/c\rfloor$.

\begin{enumerate}[label=\textup{(\roman*)},leftmargin=*]
\item If $w\ge r$, every nonempty fiber of
      \((E,R)\mapsto\operatorname{pref}_w(L_{E,R})\) determines $R$
      uniquely and, after $R$ is recovered, is exactly one quotient-prefix
      fiber.  More precisely, for every prefix value $z$ there are unique
      data $R_z$ and $\eta_z\in\B^d$, whenever the fiber is nonempty, such
      that
      \[
      \begin{split}
       &\{(E,R):\operatorname{pref}_w(L_{E,R})=z\}\\
       &\quad=
       \{(E,R_z): E\in\tbinom{Q\setminus\phi(R_z)}m,
           \ \operatorname{pref}_d(V_E)=\eta_z\}.
      \end{split}                                  \tag{QR2}\label{eq:qr-fiber-exact}
      \]
      Thus no factor counting possible remainder sets occurs in a fixed
      global prefix fiber.

\item If $w<r$, then necessarily $w<c$, and the quotient set $E$ is
      invisible to the first $w$ coefficients.  There is a fixed
      unitriangular bijection $T_{\phi,m,w}:\B^w\to\B^w$ such that
      \[
          \operatorname{pref}_w(L_{E,R})
          =T_{\phi,m,w}(\operatorname{pref}_w(P_R))                 \tag{QR3}
      \]
      for every admissible $(E,R)$.  Consequently, with the convention
      \(\binom uv=0\) for $v>u$,
      \[
      \begin{split}
       &\#\{(E,R):\operatorname{pref}_w(L_{E,R})=z\}\\
       &\quad=
       \sum_{\substack{R\in\binom Dr\\
          \operatorname{pref}_w(P_R)=T_{\phi,m,w}^{-1}(z)}}
          \binom{N-|\phi(R)|}{m}.
      \end{split}                                  \tag{QR4}\label{eq:qr-shallow-remainder}
      \]
      Thus this case is a remainder-prefix problem, multiplied by the
      number of quotient sets still available after $R$ is fixed.
\end{enumerate}
\end{theorem}

\begin{proof}
For a monic degree-$e$ polynomial $A$, set
\(A^\vee(Z)=Z^eA(Z^{-1})\).  Also set
\[
       \Phi(Z):=Z^c\phi(Z^{-1})
       =1+\phi_1Z+\cdots+\phi_cZ^c.
\]
Taking reciprocal polynomials in \eqref{eq:qr-locators} gives the exact
formal-power-series identity
\[
 L_{E,R}^\vee(Z)=P_R^\vee(Z)
 \left(
   \Phi(Z)^m+
   \sum_{j=1}^m v_j(E)Z^{cj}\Phi(Z)^{m-j}
 \right).                                          \tag{QR5}\label{eq:qr-infinity}
\]
All factors on the right have constant coefficient one except for the
displayed scalar coefficients $v_j(E)$.

Suppose first that $w\ge r$.  Since $r<c$, no term containing a $v_j$
occurs below degree $c$.  For $1\le i\le r$, the coefficient of $Z^i$
on the right of \eqref{eq:qr-infinity} is
\(p_i(R)\) plus an expression depending only on
\(p_1(R),\ldots,p_{i-1}(R)\), on $m$, and on $\Phi$.  The first $r$
coefficients therefore recover $p_1(R),\ldots,p_r(R)$ successively.
They recover the monic polynomial $P_R$, hence its root set $R$, uniquely.

We may now divide \(L_{E,R}^\vee\) by the unit
\(P_R^\vee\) modulo $Z^{w+1}$.  In the remaining factor of
\eqref{eq:qr-infinity}, the coefficient of $Z^{cj}$ is
\(v_j(E)\) plus an expression depending only on
\(v_1(E),\ldots,v_{j-1}(E)\), on $m$, and on $\Phi$.  Hence the
coefficients at degrees $c,2c,\ldots,dc$ recover
\(v_1(E),\ldots,v_d(E)\) successively.  Conversely, $P_R$ and these $d$
coefficients determine the right side of \eqref{eq:qr-infinity} modulo
$Z^{w+1}$, because every term containing $v_j$ with $j>d$ starts in degree
$cj>w$.  This proves \eqref{eq:qr-fiber-exact}; admissibility of $E$ is
exactly the condition $E\cap\phi(R)=\varnothing$.

If $w<r<c$, every term involving a $v_j$ starts in degree at least
$c>w$, so modulo $Z^{w+1}$ the bracket in \eqref{eq:qr-infinity} is the
fixed unit $\Phi(Z)^m$.  Multiplication by this unit sends the first $w$
coefficients of $P_R^\vee$ to the first $w$ coefficients of
$L_{E,R}^\vee$ by a unitriangular transformation.  It is independent of
$E$ and invertible.  For a fixed $R$, the admissible $E$ are precisely the
$m$-subsets of $Q\setminus\phi(R)$, of which there are
\(\binom{N-|\phi(R)|}{m}\).  Summing over the indicated remainder-prefix
fiber proves \eqref{eq:qr-shallow-remainder}.
\end{proof}

\begin{proposition}[Complete support factorization and its limitation]
\label{prop:complete-support-factorization}
Under the complete-fiber hypotheses of
\cref{thm:exact-quotient-remainder-normal-form}, every support
\(S\subseteq D\) has the unique decomposition
\[
 E(S)=\{y\in Q:\phi^{-1}(y)\subseteq S\},\qquad
 R(S)=S\setminus\phi^{-1}(E(S)),
 \tag{QR5a}\label{eq:complete-support-factorization}
\]
and
\(Q_S(X)=P_{R(S)}(X)V_{E(S)}(\phi(X))\).  The remainder meets every
nonfull \(\phi\)-fiber in at most \(c-1\) points, but its total size may
be \(\Theta(n)\).  The reciprocal identity \eqref{eq:qr-infinity} remains
valid for this arbitrary remainder.  When \(\abs{R(S)}\ge c\), however,
the quotient coefficients and the remainder coefficients begin to
interlace at degree \(c\), so the triangular recovery theorem above does
not give a comparison without an additional partial-occupancy atlas.
\end{proposition}

\begin{proof}
A fiber is put into \(E(S)\) exactly when all of its points lie in \(S\);
all other selected points form \(R(S)\).  This proves existence,
disjointness, the occupancy bound, and uniqueness.  Multiplying the
corresponding linear factors gives the locator identity.  The reciprocal
calculation used for \eqref{eq:qr-infinity} is purely multiplicative and
does not require \(\deg P_R<c\).  If \(\deg P_R\ge c\), both
\(p_c(R)\) and \(v_1(E)\) occur in degree \(c\), which proves the stated
limitation of the triangular separation.
\end{proof}

\begin{theorem}[Exact partial-occupancy disintegration]
\label{thm:exact-partial-occupancy}
Let \(D=D_0\sqcup X\), where \(\abs X=b\), and let
\(\phi:D_0\twoheadrightarrow Q\) have exactly \(N=\abs Q\) fibers, each
of cardinality \(c\).  For \(S\subseteq D\), put
\[
 X_S=S\cap X,\qquad
 E(S)=\{y\in Q:\phi^{-1}(y)\subseteq S\},\qquad
 R(S)=(S\cap D_0)\setminus\phi^{-1}(E(S)).
\]
Write
\[
 t=\abs{X_S},\qquad m=\abs{E(S)},\qquad
 p=\abs{\phi(R(S))},\qquad r=\abs{R(S)}.
\]
Then this datum is unique, every fiber meeting \(R(S)\) contains between
one and \(c-1\) points of \(R(S)\), and \(\abs S=t+cm+r\).  If
\(\Omega_{t,m,p,r}\) denotes the family of supports with these four
parameters, then
\[
 \abs{\Omega_{t,m,p,r}}
 =\binom bt\binom Np\binom{N-p}m
   [x^r]\bigl((1+x)^c-1-x^c\bigr)^p.                 \tag{PO1}
\]
Consequently the nonempty families \(\Omega_{t,m,p,r}\) partition
\(\binom Da\), and their exact add-back is
\[
 \sum_{\substack{t,m,p,r\\t+cm+r=a}}\abs{\Omega_{t,m,p,r}}
 =\binom{cN+b}{a}.                                   \tag{PO2}
\]
If \(\phi\) is the restriction of a monic degree-\(c\) polynomial and
every fiber \(\phi^{-1}(y)\) is the complete simple root set in \(D_0\)
of \(\phi(X)-y\), then their locators satisfy
\[
             Q_S=Q_{X_S}P_{R(S)}V_{E(S)}(\phi).
\]
\end{theorem}

\begin{proof}
Full fibers, partial fibers, and empty fibers are mutually exclusive, so
the displayed decomposition is canonical.  Choose the \(t\) exceptional
points, the \(p\) partial fibers, and the \(m\) full fibers in
\(\binom bt\binom Np\binom{N-p}m\) ways.  On one partial fiber the
weight enumerator of the allowed nonempty proper subsets is
\((1+x)^c-1-x^c\); hence the coefficient in \textup{(PO1)} counts the
choices having total partial weight \(r\).  Summing \textup{(PO1)} and
extracting the coefficient of \(x^a\) gives
\[
 [x^a](1+x)^b
 \left(1+x^c+\bigl((1+x)^c-1-x^c\bigr)\right)^N
 =[x^a](1+x)^{b+cN}=\binom{b+cN}{a}.
\]
Multiplying the linear factors belonging to the three disjoint parts of
the support proves the locator identity.
\end{proof}

\begin{proposition}[Partial-occupancy Fourier factorization]
\label{prop:partial-occupancy-fourier}
In the setting of \cref{thm:exact-partial-occupancy}, let \(G\) be a
finite abelian group, let \(g:D\to G\), and put
\(\Phi(S)=\sum_{x\in S}g(x)\).  For a character
\(\chi\in\widehat G\) and a regular fiber \(F_y=\phi^{-1}(y)\), define
\[
 A_y(\chi)=\prod_{x\in F_y}\chi(g(x)),\qquad
 P_y(\chi;z)=
 \prod_{x\in F_y}\bigl(1+z\chi(g(x))\bigr)-1-z^cA_y(\chi).
\]
Then
\[
 \sum_{S\in\Omega_{t,m,p,r}}\chi(\Phi(S))
 =[u^tv^m\zeta^pz^r]
 \prod_{x\in X}\bigl(1+u\chi(g(x))\bigr)
 \prod_{y\in Q}
 \bigl(1+vA_y(\chi)+\zeta P_y(\chi;z)\bigr).          \tag{PO3}
\]
Consequently, for every \(s\in G\),
\[
 \abs{\{S\in\Omega_{t,m,p,r}:\Phi(S)=s\}}
 =\frac1{\abs G}\sum_{\chi\in\widehat G}
   \overline{\chi(s)}
   \sum_{S\in\Omega_{t,m,p,r}}\chi(\Phi(S)).         \tag{PO4}
\]
At the trivial character, \textup{(PO3)} is exactly the support count
\textup{(PO1)}.  Thus support add-back is unconditional, while a finite
flatness theorem is exactly a bound for the nontrivial-character sum in
\textup{(PO4)}.
\end{proposition}

\begin{proof}
For a fixed regular fiber the three terms in the second product encode,
respectively, an empty fiber, a full fiber, and a nonempty proper subset.
The variables \(v,\zeta,z\) record the numbers \(m,p,r\), while \(u\)
records the exceptional weight \(t\).  Multiplicativity of \(\chi\) turns
the character of the support sum into the product of its selected-point
characters, proving \textup{(PO3)} by coefficient extraction.
Formula \textup{(PO4)} is character orthogonality on \(G\).
\end{proof}

\begin{remark}[Effective normalization of a partial-occupancy slice]
Fix a nonempty slice \(\Omega_\lambda=\Omega_{t,m,p,r}\), choose
\(S_0\in\Omega_\lambda\), and let
\[
 G_\lambda=
 \left\langle\Phi(S)-\Phi(S_0):S\in\Omega_\lambda\right\rangle
 \le G.
\]
The attained boundary values lie in the affine coset
\(\Phi(S_0)+G_\lambda\), and the effective inversion is
\[
 \abs{\{S\in\Omega_\lambda:\Phi(S)=s\}}
 =\frac1{\abs{G_\lambda}}
   \sum_{\chi\in\widehat{G_\lambda}}
   \overline{\chi(s-\Phi(S_0))}
   \sum_{S\in\Omega_\lambda}
      \chi(\Phi(S)-\Phi(S_0)).                       \tag{PO5}
\]
Every character of \(G_\lambda\) extends to the ambient finite abelian
group, so the inner sum may be evaluated from \textup{(PO3)} after
multiplication by the fixed factor
\(\overline{\chi(\Phi(S_0))}\).  Different extensions agree on support
differences.  Equivalently, ambient characters group by their restriction
to \(G_\lambda\), and the annihilator multiplicity cancels exactly against
the ambient denominator.  This identifies the correct Fourier denominator;
it does not assert that the realized image fills the affine group.
\end{remark}

\begin{remark}[What partial-occupancy add-back does and does not prove]
The exponentially many local proper-subset choices are already contained
inside the coarse slices \(\Omega_{t,m,p,r}\); they incur no separate
profile-count loss in \textup{(PO2)}.  The theorem does not bound the
nontrivial Fourier terms in \textup{(PO4)}, the size of a realized
boundary image, or the projection of a witness cell to distinct slopes.
Those are separate analytic and ray payments.
\end{remark}


\subsection{Asymptotic and finite regimes}\label{sec:complementary}


\begin{theorem}[Unconditional shallow-prefix RS--MCA exponent]
\label{thm:unconditional-asymptotic-rs-mca}
Let \(C_n=\RS_{\F_n}(D_n,k_n)\), where
\(D_n\subseteq\B_n\subseteq\F_n\), \(\abs{D_n}=n\), and
\(k_n/n\to\rho\in(0,1)\).  Let \(a_n\ge k_n+1\), put
\(w_n=a_n-k_n-1\), and assume
\[
 \frac{a_n}{n}\to\alpha\in(0,1),\qquad
 w_n\log\abs{\B_n}=o(n),                           \tag{AS1}
\]
\(\abs{\F_n}>n\), and, for some finite \(\phi\ge0\),
\[
 \frac1n\log\abs{\F_n}\to\phi,
 \qquad
 \frac1n\log(\abs{\F_n}-n)\to\phi.               \tag{AS2}
\]
For the full-field challenge,
\[
 \frac1n\log B_{C_n}^{\rm MCA}(a_n)
   =\min\{h(\alpha),\phi\}+o(1),                   \tag{AS3}
\]
and therefore
\[
 \frac1n\log e_{\rm MCA}
   \left(C_n,1-\frac{a_n}{n}\right)
   =-\bigl(\phi-h(\alpha)\bigr)_++o(1).            \tag{AS4}
\]
Because \(\abs{\B_n}\ge n\), condition \textup{(AS1)} forces
\(\alpha=\rho\).  No smoothness, Fourier estimate, or auxiliary ray
hypothesis is assumed.
\end{theorem}

\begin{proof}
Put \(M_n=\binom n{a_n}\) and
\(L_n=\ceil{M_n\abs{\B_n}^{-w_n}}\).  Exact-agreement reduction and
fixed-support slope uniqueness give the unconditional upper bound
\[
 B_{C_n}^{\rm MCA}(a_n)\le\min\{\abs{\F_n},M_n\}.  \tag{AS5}
\]
The exact locator-prefix list and collision-aware pole compiler give
\[
 B_{C_n}^{\rm MCA}(a_n)\ge
 \left\lceil
 \frac{L_n(\abs{\F_n}-n)}
 {\abs{\F_n}-n+k_n(L_n-1)}
 \right\rceil.                                     \tag{AS6}
\]
By Stirling's formula and \textup{(AS1)},
\(\log L_n=h(\alpha)n+o(n)\).  For positive \(x,y\),
\(xy/(x+k_ny)\ge\frac12\min\{y,x/k_n\}\); hence
\textup{(AS2)}, \(\log k_n=o(n)\), and the ceiling show that the
right side of \textup{(AS6)} has logarithm
\(n\min\{h(\alpha),\phi\}+o(n)\).  The upper bound
\textup{(AS5)} has the same exponent, proving \textup{(AS3)}; subtracting
\(\log\abs{\F_n}\) proves \textup{(AS4)}.  Finally
\(w_n\log n=o(n)\), so \((a_n-k_n)/n\to0\) and
\(\alpha=\rho\).
\end{proof}

\begin{corollary}[Unconditional large-challenge RS--MCA threshold]
\label{cor:unconditional-rs-mca-capacity}
Under the row hypotheses of
\cref{thm:unconditional-asymptotic-rs-mca}, let the full-field target
satisfy \(\eps_n=\exp(-\lambda n+o(n))\) for a finite
\(\lambda\ge0\).  If
\[
                         \lambda<\phi-h(\rho),      \tag{AS7}
\]
then, for every sufficiently large \(n\),
\[
 a_n^*=k_n+1,\qquad
 \delta_n^*=1-\frac{k_n+1}{n}=1-\rho+o(1).         \tag{AS8}
\]
Conversely, if \(\lambda>(\phi-h(\rho))_+\), every agreement sequence
satisfying \((a_n-k_n-1)\log\abs{\B_n}=o(n)\) is unsafe.  The critical
boundary \(\lambda=(\phi-h(\rho))_+\) is governed by the
exact finite constants and is not decided by exponent asymptotics.  This
is a large scalar-extension regime.  If \(\abs{\F_n}\) is polynomial in
\(n\), then \(\phi=0\) and \textup{(AS7)} cannot hold; the corollary does
not settle that Proximity-Prize parameter regime.
\end{corollary}

\begin{proof}
At \(a=k_n+1\), \textup{(AS5)} is
\(B_{C_n}^{\rm MCA}(a)\le\binom n{k_n+1}
=\exp(h(\rho)n+o(n))\).  Under \textup{(AS7)},
\(\eps_n\abs{\F_n}=\exp((\phi-\lambda)n+o(n))\) has strictly larger
exponent, and taking the integer floor does not change the comparison.
Thus the first permitted agreement is safe
and \textup{(AS8)} follows.  The converse follows from \textup{(AS4)}
and the strict reverse exponential inequality.
\end{proof}


\begin{corollary}[Exact binary additive-domain examples]
\label{cor:binary-additive-examples}
Let \(k=2^{40}\), let \(D\) be any \(\F_2\)-linear subspace of the
indicated field having the indicated size, and take the target
\(2^{-128}\).  The four rows
\[
\begin{array}{c|c|c|c}
 k/n&\F&n&\delta_{C,\mathrm{sup}}^*\\ \hline
 1/2&\F_{2^{166}}&2^{41}&1/8\\
 1/4&\F_{2^{168}}&2^{42}&1/4\\
 1/8&\F_{2^{169}}&2^{43}&1/4\\
 1/16&\F_{2^{170}}&2^{44}&1/4
\end{array}
\]
have exactly the displayed supremal MCA transitions, and every endpoint
is unsafe.
\end{corollary}

\begin{proof}
For \(q=2^m\), the slope budget is
\(B=\floor{q2^{-128}}=2^{m-128}\).  In the four rows this is respectively
\(2^{38},2^{40},2^{41},2^{42}\), and
\(3(B-1)\le n-k\).  Apply
\cref{thm:intro-unconditional-asymptotic-deep} and translate the largest
safe grid radius \(B-1\) to the real closed-ball supremum \(B/n\).
\end{proof}

\begin{proposition}[Exact unsafe test]\label{prop:simple-pole-lower}
Let
\[
 L_n=\left\lceil
 \binom n{a_n}\abs{\B_n}^{-(a_n-k_n-1)}
 \right\rceil,
 \qquad q_n=\abs{\F_n}>n.
\]
If
\[
 \left\lceil \frac{\abs{\Gamma_n}}{q_n}
 \left\lceil
 \frac{L_n(q_n-n)}{q_n-n+k_n(L_n-1)}
 \right\rceil\right\rceil>B_n^*,
 \tag{13.3}\label{eq:exact-unsafe-budget}
\]
then agreement \(a_n\) is unsafe at the target.  The same statement holds
with \(L_n\) replaced by any larger identity, quotient, Chebyshev, or
remainder-profile list proved for the dimension-\(k_n+1\) code.
\end{proposition}

\begin{proof}
The identity list is \cref{prop:exact-prefix-list}, and
\cref{thm:collision-aware-pole} produces a bad-slope set \(Z\subseteq\F_n\)
of the inner-ceiling size.  Replacing a line \((r_0,r_1)\) by
\((r_0+\delta r_1,r_1)\) translates \(Z\) by \(-\delta\).  Averaging over
\(\delta\in\F_n\) gives \(\abs Z\abs{\Gamma_n}/q_n\) challenge slopes, so
some translate has at least the outer-ceiling number.  The profile variants
use the same conversion after their list sizes have been established.
\end{proof}

\begin{theorem}[Unconditional exact support-envelope bracket]
\label{thm:unconditional-support-envelope-bracket}
Let \(C=\RS_\F(D,k)\), let \(D\subseteq\B\subseteq\F\), put
\(q=\abs\F>n\), fix \(\varnothing\ne\Gamma\subseteq\F\), and let
\(B^*=\floor{\eps\abs\Gamma}\).  For \(k+1\le a\le n\), define
\[
 \begin{split}
 L(a)&=\ceil{\binom na\abs\B^{-(a-k-1)}},\\
 P(a)&=\left\lceil\frac{\abs\Gamma}{q}
       \left\lceil\frac{L(a)(q-n)}
       {q-n+k(L(a)-1)}\right\rceil\right\rceil,\\
 U(a)&=\min\left\{\abs\Gamma,\binom na\right\}.
 \end{split}                                       \tag{SB1}
\]
If \(a_-<a_+\) satisfy
\[
                         P(a_-)>B^*,\qquad U(a_+)\le B^*, \tag{SB2}
\]
then the global target threshold exists and
\[
                         a_-<a^*\le a_+.            \tag{SB3}
\]
If \(a_+=a_-+1\), then \(a^*=a_+\) exactly.  In particular,
\[
                  \binom n{k+1}\le B^*\quad\Longrightarrow\quad
                  a^*=k+1.                         \tag{SB4}
\]
These are unconditional finite comparisons; \textup{(SB2)} is the
literal target reserve and contains no asymptotic placeholder.
\end{theorem}

\begin{proof}
The exact prefix list and pole-line construction give
\(B_{C,\Gamma}^{\rm MCA}(a)\ge P(a)\) by
\cref{prop:simple-pole-lower}.  The exact support atlas gives
\(B_{C,\Gamma}^{\rm MCA}(a)\le U(a)\) by
\cref{prop:exact-support-upper}.  Thus \(a_-\) is unsafe and \(a_+\) is
safe; monotonicity proves \textup{(SB3)} and the adjacent conclusion.
For \textup{(SB4)}, the first permitted agreement is already safe.
\end{proof}



\section{Conclusion}\label{sec:conclusion}

The finite content of the paper is the quadratic mean-overlap theorem: a
tangent construction supplies \(r+1\) slopes, and a line--secant overlap
argument proves that no received line supplies more.  We are unaware of
a prior finite MDS theorem of this form, and its small parameter pockets
are not certified by the specialized estimate
\eqref{eq:bchks-specialized}.  The example
\((n,k,r)=(32,15,6)\), propagated to positive-density smooth and circle
prime families, isolates this finite contribution.

At constant relative radius the prior BCHKS CA theorem, combined with
the below-half equality of ACFY/WHIR, already gives the stronger global
conclusion; the exact sparse decomposition here supplies a self-contained
and challenge-restricted route.  Making the literature implication
integral completes the exact staircase through every compact subinterval
below half the minimum distance and widens the target compiler
correspondingly.  The circle statements add no hidden coding assumption:
monomial and stereographic equivalences preserve each bad-slope set, not
only its cardinality.

The target comparison remains genuinely discrete.  If
\(b=\floor{\eps\abs\Gamma}\), the largest safe grid point is
\((b-1)/n\), while the real-radius safe set can be half open with unsafe
supremum \(b/n\).  The certified Proth and Mersenne-31 rows illustrate
this endpoint bookkeeping, but the four Proth thresholds are not claimed
as new consequences.

No theorem here reaches the unrestricted near-capacity smooth/circle
frontier.  The fixed-length prime families have positive density only
with length held fixed, and growing dyadic prime-field families
necessarily have density zero.  Progress beyond half distance still
requires new received-line or list-decoding information; no computational
certificate, profile ledger, or unproved equidistribution statement has
been promoted to such an input in this paper.

\appendix
\section{Verification details}\label{sec:verification-details}

This appendix records the elementary coding-theoretic and integer
calculations used above in a form that can be checked independently of the
geometric proofs.  Some of these facts are standard, but including their
short derivations fixes all normalization conventions and makes the four
large finite rows reproducible from the decimal data printed in the paper.

\subsection{Reed--Solomon and MDS conventions}

Let \(D=\{x_1,\ldots,x_n\}\subseteq\F\) consist of distinct points and
let \(1\le k<n\).  Evaluation on \(D\) is injective on
\(\F[X]_{<k}\): a nonzero polynomial of degree below \(k\) cannot vanish
at the \(n\ge k\) distinct points of \(D\).  It follows that
\(\RS_\F(D,k)\) has dimension \(k\).  If two codewords are distinct,
their difference is the evaluation of a nonzero polynomial of degree at
most \(k-1\), hence has at most \(k-1\) zero coordinates and at least
\(n-k+1\) nonzero coordinates.  Conversely, the evaluation of
\(\prod_{i=1}^{k-1}(X-x_i)\) has exactly \(k-1\) zero coordinates.
Thus the minimum distance is
\[
                         d_{\min}=n-k+1.
\]
This proves directly that the code is MDS.

For \(x\in D\), put
\[
 \lambda_x=\left(\prod_{y\in D\setminus\{x\}}(x-y)\right)^{-1}.
\]
The Lagrange interpolation formula for a polynomial \(f\) of degree at
most \(n-1\) shows that its leading coefficient is
\(\sum_{x\in D}\lambda_x f(x)\).  Applying this to
\(f(X)=X^j\), \(0\le j\le n-2\), gives
\[
                         \sum_{x\in D}\lambda_xx^j=0.
\]
Let \(R=n-k\) and let \(H\) be the \(R\)-by-\(n\) matrix whose column at
\(x\) is
\[
                    h_x=\lambda_x(1,x,\ldots,x^{R-1})^{\mathsf T}.
\]
If \(P\in\F[X]_{<k}\), then the \(j\)-th entry of
\(H(P(x))_{x\in D}^{\mathsf T}\) is
\(\sum_x\lambda_xx^jP(x)=0\), because \(j+\deg P\le n-2\).
Thus \(H\) annihilates the code.  Every \(R\)-by-\(R\) submatrix of
\(H\) is a nonzero diagonal scaling of a Vandermonde matrix, so it is
invertible.  Hence \(H\) has rank \(R\), its kernel has dimension \(k\),
and its kernel is exactly \(\RS_\F(D,k)\).

The same Vandermonde calculation proves that every set of at most \(R\)
columns is linearly independent.  In particular, for
\(E\subseteq D\), \(|E|\le R\), the space
\(V_E=\Span\{h_x:x\in E\}\) has dimension \(|E|\).  A word \(u\) agrees
with a codeword outside \(E\) if and only if its syndrome lies in \(V_E\).
Indeed, writing \(u=c+e\) with \(e\) supported on \(E\) gives one
direction.  Conversely, if
\(Hu^{\mathsf T}=\sum_{x\in E}e_xh_x\), subtracting the word whose
coordinates on \(E\) are \(e_x\) leaves zero syndrome and therefore a
codeword.  This is the precise error-support form of the syndrome
normalization used in \cref{prop:syndrome-line-normal-form}.

For completeness, the fixed-support slope uniqueness used repeatedly in
the paper is also an MDS-free linear observation.  Suppose a support
\(S\) explains \(u_0+\gamma u_1\) and \(u_0+\gamma' u_1\) for distinct
\(\gamma,\gamma'\).  Subtracting the two explaining codewords and dividing
by \(\gamma-\gamma'\) explains \(u_1\) on \(S\); subtracting
\(\gamma\) times this explanation from the first explains \(u_0\) there.
Thus \(S\) is common.  Consequently a noncommon support contributes at
most one finite MCA-bad slope.

\begin{lemma}[Monotonicity and exact endpoint conversion]
\label{lem:appendix-endpoint}
For fixed \(C\) and \(\Gamma\), the integer numerator
\(B_{C,\Gamma}^{\rm MCA}(a)\) is nonincreasing in \(a\).  Let
\(B^*=\floor{\eps|\Gamma|}\), and suppose \(a_0\) is the first safe
agreement, while \(a_0-1\) is unsafe.  Then the largest safe grid radius is
\(1-a_0/n\), whereas the safe real radii form
\[
                 \left[0,\frac{n-a_0+1}{n}\right),
\]
so their supremum is one grid step larger and is not attained.
\end{lemma}

\begin{proof}
Every witness at agreement at least \(a+1\) is also a witness at
agreement at least \(a\), proving monotonicity.  For a real radius
\(\delta\), the closed Hamming ball permits precisely
\(r=\floor{\delta n}\) errors and therefore requires agreement
\(n-r\).  The condition \(n-\floor{\delta n}\ge a_0\) is equivalent to
\(\floor{\delta n}\le n-a_0\), or
\(\delta<(n-a_0+1)/n\).  The largest grid point in that interval is
\((n-a_0)/n=1-a_0/n\).
\end{proof}

This lemma explains the half-open intervals in the main theorem.  In
particular, at target \(2^{-128}\), writing
\(B=\floor{q/2^{128}}\), exact numerator \(B\) at error radius \(B-1\)
and numerator at least \(B+1\) at radius \(B\) give the safe set
\([0,B/n)\), rather than a closed interval.

\subsection{The quadratic boundary}

Put \(R=n-k\) and
\[
                  F_{n,k}(r)=r^2-n(3r-R).
\]
The mean-overlap hypothesis is \(F_{n,k}(r)\ge0\).  Its two roots are
\[
       \frac{3n\pm\sqrt{9n^2-4nR}}2
       =\frac{3n\pm\sqrt{n(5n+4k)}}2.
\]
Since \(F_{n,k}(0)=nR>0\) and
\(F'_{n,k}(r)=2r-3n<0\) for \(0\le r\le n\), the admissible initial
interval ends at
\[
 r_{\rm quad}(n,k)=
 \floor{\frac{3n-\sqrt{n(5n+4k)}}2}.
\]
If \(k/n\to\rho\), division by \(n\) gives
\[
 \frac{r_{\rm quad}(n,k)}n
 \longrightarrow
 \alpha_\rho=\frac{3-\sqrt{5+4\rho}}2.
\]
The derivative observation also shows that checking
\(F_{n,k}(B-1)\ge0>F_{n,k}(B)\) proves the exact identity
\(r_{\rm quad}(n,k)=B-1\); no floating-point approximation is involved.

At the four challenge rates the limiting constants are, for orientation,
\[
\begin{array}{c|c|c}
 \rho&\alpha_\rho&\text{decimal value}\\ \hline
 1/2&(3-\sqrt7)/2&0.1771243444\ldots\\
 1/4&(3-\sqrt6)/2&0.2752551286\ldots\\
 1/8&(6-\sqrt{22})/4&0.3273960600\ldots\\
 1/16&(6-\sqrt{21})/4&0.3543560762\ldots
\end{array}
\]
The decimals are not used in any proof; the exact signs in
\eqref{eq:quadratic-row-certificates} locate the finite boundaries.

\subsection{Proth certificates and budget divisions}

We spell out the certificate format.  Suppose
\(p=u2^s+1\), where \(u\) is odd and \(u<2^s\), and suppose an integer
\(a_0\) satisfies
\[
                         a_0^{(p-1)/2}\equiv-1\pmod p.
\]
If a prime \(\ell\) divides \(p\), then the order of \(a_0\) modulo
\(\ell\) has exact \(2\)-part \(2^s\).  Hence \(2^s\mid\ell-1\), so
\(\ell\ge2^s+1\).  The inequality \(u<2^s\) gives
\(p=u2^s+1<2^{2s}+1\), and therefore \(2^s>\sqrt p-1\).
If \(p\) were composite it would have a prime divisor at most
\(\sqrt p\), contradicting the lower bound \(\ell\ge2^s+1>\sqrt p\).
Thus \(p\) is prime.  This is Proth's theorem in exactly the form needed
here.

All certificate inputs are collected below.  The equality in the last
column is a residue equality modulo the corresponding \(p\).
\[
\begin{array}{c|r|r|r|r|r}
 \rho&s&u&a_0&\operatorname{bits}(p)&
 a_0^{(p-1)/2}\bmod p\\ \hline
 1/2&92&26766274163673319604503&3&167&p-1\\
 1/4&93&41595378994516821279015&13&169&p-1\\
 1/8&95&24737346889219389259839&5&170&p-1\\
 1/16&97&13387194060291799253121&5&171&p-1
\end{array}
\]
Together with the decimal values of \(p\) printed before
\cref{prop:proth-row-check}, this table is a complete Proth certificate:
the reader checks the single exact identity \(p-1=u2^s\), the parity and
size of \(u\), and one modular exponentiation.  Modular exponentiation is
an integer operation and does not rely on a probabilistic primality test.

For a completely mechanical check, first compute
\(t=a_0^u\bmod p\) by binary square-and-multiply, then square \(t\)
exactly \(s-1\) times modulo \(p\).  The certificate accepts precisely
when the result is \(p-1\).  The algorithm uses
\(O(\log u+s)\) modular multiplications.  The remaining row checks use
only division with remainder.  Writing
\[
                         p=B2^{128}+r_p,
\]
the four exact remainders are
\[
\begin{array}{c|r}
 \rho&r_p\\ \hline
 1/2&1381541083842484386787422633985\\
 1/4&2921538492713497448761933168641\\
 1/8&2495687119199326634196634435585\\
 1/16&20440865928680199099134339186689
\end{array}
\]
and each lies strictly between \(0\) and
\(2^{128}=340282366920938463463374607431768211456\).  Therefore the
displayed \(B\) is exactly \(\floor{p/2^{128}}\).

Finally, \(n=2^e\) with \(e=41,42,43,44\), respectively, while
\(s=92,93,95,97\).  Thus \(n\mid2^s\mid p-1\), and the cyclic group
\(\F_p^\times\) contains a unique subgroup of order \(n\).  Each field
has the printed bit length, all smaller than \(256\), and
\(k=2^{40}\) gives the required four rates.  The exact boundary signs are
\[
\begin{array}{c|r|r}
 \rho&F_{n,k}(B-1)&F_{n,k}(B)\\ \hline
 1/2&5154112775168&-663955886271\\
 1/4&7590647904465&-3182321912768\\
 1/8&13908181940112&-6720484728007\\
 1/16&19335616403905&-20973145690236
\end{array}
\]
and hence \(B-1=r_{\rm quad}(n,k)\) in every row.

\begin{proposition}[Independent certificate verification]
\label{prop:appendix-certificate-verification}
The data in \cref{tab:prize-proth-rows} define four prime fields, four
power-of-two multiplicative subgroups of the stated lengths, and slope
budgets \(B=\floor{p2^{-128}}\).  In every row the exact safe set is
\([0,B/n)\), and the endpoint is unsafe.
\end{proposition}

\begin{proof}
The Proth calculation proves primality, the divisibility calculation
produces the subgroup, and the remainder calculation gives the target
budget.  The two signs locate the quadratic boundary at \(B-1\).
Apply \cref{cor:prize-window-compiler} and then
\cref{lem:appendix-endpoint}.  Every step is an equality or inequality
between the printed integers.
\end{proof}

\begin{thebibliography}{GGSW26}

\bibitem[ABF26]{ABF26}
G. Arnon, D. Boneh, and G. Fenzi,
\newblock Open problems in list decoding and correlated agreement,
\newblock IACR Cryptology ePrint Archive, Report 2026/680, 2026,
revised July 6, 2026.

\bibitem[BCHKS25]{BCHKS25}
E. Ben-Sasson, D. Carmon, U. Hab\"ock, S. Kopparty, and S. Saraf,
\newblock On proximity gaps for Reed--Solomon codes,
\newblock IACR Cryptology ePrint Archive, Report 2025/2055, 2025.

\bibitem[BCIKS20]{BCIKS20}
E. Ben-Sasson, D. Carmon, Y. Ishai, S. Kopparty, and S. Saraf,
\newblock Proximity gaps for Reed--Solomon codes,
\newblock \emph{J. ACM} 70 (2023), no.~5, 1--57,
\newblock doi:10.1145/3614423; conference version in
\emph{61st IEEE Symposium on Foundations of Computer Science
(FOCS 2020)} and extended version, IACR ePrint 2020/654.

\bibitem[BCGM25]{BordageChiesaGuanManzur25}
S. Bordage, A. Chiesa, Z. Guan, and I. Manzur,
\newblock All polynomial generators preserve distance with mutual
correlated agreement,
\newblock IACR Cryptology ePrint Archive, Report 2025/2051, 2025,
 revised May 2026; conference version in CCC 2026.

\bibitem[CS25]{CS25}
E. Crites and A. Stewart,
\newblock On Reed--Solomon proximity gaps conjectures,
\newblock IACR Cryptology ePrint Archive, Report 2025/2046, 2025.

\bibitem[Dav00]{Davenport}
H. Davenport,
\newblock \emph{Multiplicative Number Theory}, third ed.,
\newblock Graduate Texts in Mathematics 74, Springer, New York, 2000.

\bibitem[GG25]{GoyalGuruswami25}
R. Goyal and V. Guruswami,
\newblock Optimal proximity gaps for subspace-design codes and (random)
Reed--Solomon codes,
\newblock IACR Cryptology ePrint Archive, Report 2025/2054, 2025,
revised March 2026.

\bibitem[GGSW26]{GoyalGuruswamiSunWootters26}
R. Goyal, V. Guruswami, Y. Sun, and M. Wootters,
\newblock Locality of curve-decoding and improved proximity gaps,
\newblock arXiv:2607.08516, 2026.

\bibitem[HLP24]{HLP24}
U. Hab\"ock, D. Levit, and S. Papini,
\newblock Circle STARKs,
\newblock IACR Cryptology ePrint Archive, Report 2024/278, 2024.

\bibitem[Kam26]{Kambire26}
A. Kambir\'e,
\newblock Proximity gaps conjecture fails near capacity over prime fields,
\newblock arXiv:2604.09724, 2026.

\bibitem[KKH26]{KKH26}
D. Krachun, S. Kazanin, and U. Hab\"ock,
\newblock Failure of proximity gaps close to capacity,
\newblock IACR Cryptology ePrint Archive, Report 2026/782, 2026.

\bibitem[ACFY25]{WHIR}
G. Arnon, A. Chiesa, G. Fenzi, and E. Yogev,
\newblock WHIR: Reed--Solomon proximity testing with super-fast verification,
\newblock in \emph{EUROCRYPT 2025}, 214--243;
IACR ePrint 2024/1586.

\bibitem[Hab25]{Habock25}
U. Hab\"ock,
\newblock A note on mutual correlated agreement for Reed--Solomon codes,
\newblock IACR Cryptology ePrint Archive, Report 2025/2110, 2025.

\bibitem[PP26]{ProximityPrize}
{Ethereum Foundation},
\newblock The Proximity Prize: \$1M Reed--Solomon Challenge,
\newblock preliminary version, 2026,
\newblock \url{https://proximityprize.org/}, accessed July 11, 2026.

\bibitem[YZ26]{YuanZhu26}
C. Yuan and R. Zhu,
\newblock A syndrome-space approach to proximity gaps and correlated
agreement for random linear codes,
\newblock arXiv:2605.07595, 2026.

\end{thebibliography}

\end{document}
